\documentclass[a4paper, 11pt, twoside]{book}

\usepackage{comment}

% ---------------- Encodage ------------- %
\usepackage[utf8]{inputenc}
\usepackage[T1]{fontenc}
\usepackage[french]{babel}
\usepackage[autolanguage]{numprint}
\usepackage{csquotes}
% ---------------- Encodage ------------- %

%------- Bibliographie ----------
\usepackage[backend=biber, maxbibnames=99, style=science]{biblatex}
\addbibresource{refs.bib}
%------- Bibliographie ----------

% -------------- Mise en page ---------- %
\usepackage[top=1.5cm, bottom=1.5cm, left = 1.5cm, right=3cm]{geometry}

\usepackage{fancyhdr}
\pagestyle{fancy}

\usepackage{lastpage}
%\fancyfoot[C]{\thepage \ / \pageref{LastPage}}
\fancyhead[L]{L'essentiel des mathématiques de première année de Licence}
\fancyhead[R]{Année scolaire 2022-2023}
\usepackage{indentfirst}
\usepackage{setspace}
\usepackage[usenames, dvipsnames]{xcolor}
\usepackage{booktabs}
\usepackage{multicol}
\usepackage{caption}
\usepackage{subcaption}
\usepackage{diagbox}
\usepackage{xspace}

\setlength{\tabcolsep}{10pt}
\usepackage{needspace}
\usepackage[Lenny]{fncychap}
\usepackage{lmodern}
% -------------- Mise en page ---------- %

% ------------ Boxes -------------- %
\usepackage{framed}
\usepackage[most]{tcolorbox}
% ------------- Boxes -------------- %

% ---------------- Maths -------------- %
\usepackage{mathtools, amssymb, amsthm}
\usepackage[makeroom]{cancel}
\usepackage{centernot}
\usepackage{systeme}
\usepackage{stmaryrd}
\usepackage{dsfont}
% ---------------- Maths -------------- %

% ------------- Théorèmes ------------ %
\newtcolorbox{theorem}[1][]{enhanced, colback=white, colframe=green!60!black, title=\textbf{Théorème} : {#1}, breakable, attach boxed title to top text left={yshift*=-3mm}, boxed title style={size=small, colback=green!60!black, sharpish corners}, sharpish corners}

\newtcolorbox{lemma}[1][]{enhanced, colback=white, colframe=red!60!black, title=\textbf{Lemme} : {#1}, breakable, attach boxed title to top text left={yshift*=-3mm}, boxed title style={size=small, colback=red!60!black, sharpish corners}, sharpish corners}

\newtcolorbox{corollary}[1][]{enhanced, colback=white, colframe=green!60!black, title=\textbf{Corollaire} : {#1}, breakable, attach boxed title to top text left={yshift*=-3mm}, boxed title style={size=small, colback=green!60!black, sharpish corners}, sharpish corners}

\newtcolorbox{proposition}[1][]{enhanced, colback=white, colframe=red!60!black, title=\textbf{Proposition} : {#1}, breakable, attach boxed title to top text left={yshift*=-3mm}, boxed title style={size=small, colback=red!60!black, sharpish corners}, sharpish corners}

\newtcolorbox{definition}[1][]{enhanced, colback=white, colframe=blue!60!black, title=\textbf{Définition} : {#1}, breakable, attach boxed title to top text left={yshift*=-3mm}, boxed title style={size=small, colback=blue!60!black, sharpish corners}, sharpish corners}

\newtcolorbox{axiom}[1][]{enhanced, colback=white, colframe=blue!60!black, title=\textbf{Axiome} : {#1}, breakable, attach boxed title to top text left={yshift*=-3mm}, boxed title style={size=small, colback=blue!60!black, sharpish corners}, sharpish corners}

\theoremstyle{definition}
\newtheorem{example}{Exemple}[chapter]
\newtheorem{exercise}{Exercice}[chapter]

\theoremstyle{remark}
\newtheorem*{remark}{Remarque}
\newtheorem*{notation}{Notation}
\newtheorem*{nomenclature}{Nomenclature}
% ------------- Théorèmes ------------ %

% -------- Liens -------
\usepackage{hyperref}

%\hypersetup{
%    colorlinks=true,
%    allcolors=black,
%    pdftitle={Essentiel mathématiques L1},
%}
% -------- Liens -------

% ------------- Images --------------- %
\usepackage{graphicx}
% ------------- Images --------------- %

% ------------ TIKZ -------------- %
\usepackage{tikz}
\usepackage{pgfplots}
\pgfplotsset{compat=1.15}
\usepackage{mathrsfs}
\usetikzlibrary{arrows, angles, quotes, calc}

% -------- TIKZ ------------------ %

% ------------------ Physique ---------------- %
\usepackage{physics}
\usepackage{siunitx}
\AtBeginDocument{\RenewCommandCopy\qty\SI}
% ------------------ Physique ---------------- %

% ------------ Ensembles de nombres ---------- %
\newcommand{\N}{\mathbb{N}}
\newcommand{\Z}{\mathbb{Z}}
\newcommand{\D}{\mathbb{D}}
\newcommand{\Q}{\mathbb{Q}}
\newcommand{\R}{\mathbb{R}}
\newcommand{\C}{\mathbb{C}}
\newcommand{\K}{\mathbb{K}}
% ------------ Ensembles de nombres ---------- %

% ------------ Macros ---------------- %
\DeclareMathOperator{\pgcd}{pgcd}
\DeclareMathOperator{\ppcm}{ppcm}
\DeclareMathOperator{\rg}{rg}
\DeclareMathOperator{\com}{com}
\DeclareMathOperator{\End}{End}
\DeclareMathOperator{\Aut}{Aut}
\DeclareMathOperator{\Vect}{Vect}

\newcommand{\deriv}[2]{\frac{\diffd{#1}}{\diffd{#2}}}
\newcommand{\link}[2]{\href{#1}{\textcolor{blue}{\underline{#2}}}}
\newcommand{\floor}[1]{\lfloor #1 \rfloor}
\newcommand{\ceil}[1]{\lceil #1 \rceil}

\newcommand{\appli}[5]{
	\begin{array}{rrcl}
		#1 \colon & #2 & \to     & #3 \\
		          & #4 & \mapsto & #5
	\end{array}
}
% ------------ Macros ---------- %

%-----------Bloc de code------------%
\usepackage{listings}
\usepackage{color}

\lstset{frame=single,
	language=C++,
	aboveskip=3mm,
	belowskip=3mm,
	showstringspaces=false,
	columns=flexible,
	basicstyle={\small\ttfamily},
	numbers=none,
	numberstyle=\tiny\color{gray},
	keywordstyle=\color{blue},
	commentstyle=\color{dkgreen},
	stringstyle=\color{mauve},
	breaklines=true,
	breakatwhitespace=true,
	tabsize=4
}

\newcommand*{\inlineCode}{\fontfamily{pnc}\selectfont\ttfamily}

\usepackage{comment}
%-----------Bloc de code----------------%

\renewcommand{\leq}{\leqslant}
\renewcommand{\geq}{\geqslant}

\includeonly{chapters/algebre/applications_lineaires}

\begin{document}

\newgeometry{margin=2.5cm}
\begin{titlepage}
	\begin{center}
		\vspace*{\fill}
			\vspace*{1cm}
			\hrule
			\vspace{0.5cm}
			\Large \textbf{L'essentiel des mathématiques de première année de Licence}
			\\
			Raphaël Heng
			\\
            Alyce Théobald
			\vspace{0.5cm}
			\hrule
			\vspace{0.5cm}
			\begin{figure}[!h]
				\centering
				\includegraphics[scale=0.6]{img/logo_lyon1.jpg}
			\end{figure}
			\vspace{0.5cm}
			Université Claude Bernard Lyon 1 
			\\
			Licence 1 - Portail Mathématiques-Informatique 
			\\
			Année universitaire 2022-2023 
			\\
			\vspace*{1cm}
		\vspace*{\fill}
	\end{center}
\end{titlepage}
\restoregeometry

\tableofcontents

\part{Introduction}

\def\arraystretch{1.5}
\par Avant toute chose, nous tenons à préciser que cette mise en page est destinée à une impression sous la forme d'un livre. Ainsi pour profiter d'une bonne lecture sur la version numérique, nous vous recommandons de la lire en mode \og double pages \fg avec les plus grandes marges orientées vers le centre. 
\\
\par \noindent Ce document repose principalement sur les enseignements de nos professeurs Guillaume AUBRUN, Kenji IOHARA et Thomas STROBL
mais il peut contenir des formulations d'autres sources telles que \textbf{Bibmath} \cite{bibmath}, \textbf{Wikipédia} \cite{wikipedia}, \textbf{Exo7} \cite{exo7} ou encore le livre destiné aux élèves de CPGE recommandé sur le site de la licence Mathématiques \cite{livre_prepa}.

\par \noindent Il regroupe l'essentiel des compétences mathématiques à maîtriser à la fin de la première année de Licence. Vous y trouverez les définitions et les théorèmes à connaître accompagnés d'exercices à savoir refaire. Nous essaierons de démontrer le plus de théorèmes possibles, cependant les preuves ne sont pas toute à retenir (cela dépend également de votre orientation : mathématiques ou informatique).
Il se peut également que les outils mathématiques ne soient pas présentées scrupuleusement comme aux cours magistraux, nous avons éventuellement paraphrasé certains passages. Par exemple, il se peut que les notations utilisées ne soient pas les mêmes que celles vues en cours, nous avons préféré utiliser des notations plus courantes en France.
\\
\par \noindent Nous pensons qu'il est intéressant de définir certains mots de vocabulaires définis ci-dessous :
\begin{itemize}
    \item \textbf{Assertion} : Une assertion est une affirmation mathématique qui est soit vraie soit fausse. 
    \item \textbf{Axiome} : Un axiome est une assertion que l'on considère vraie sans démonstration.
    \item \textbf{Définition} : Une définition énonce comment un objet mathématique est construit.
    \item \textbf{Théorème} : Un théorème est une assertion d'importance particulière ayant été démontrée.
    \item \textbf{Corollaire} : Un corollaire est un résultat découlant d'un théorème.
    \item \textbf{Lemme} : Un lemme est un résultat intermédiaire sur lequel on s'appuie pour démontrer un théorème.
    \item \textbf{Proposition} : Une proposition est un résultat simple qui n'est pas associé à un théorème.
    \item \textbf{Conjecture} : Une conjecture est une proposition dont on ignore la véracité.
\end{itemize}
\par \noindent Rappelons également les ensembles de nombres étudiés au lycée :
\[ \text{L'ensemble des entiers naturels : } \N = \{0, 1, \ldots \} \]
\[ \text{L'ensemble des entiers relatifs : } \Z = \{ \ldots, -1, 0, 1, \ldots \} \]
\[ \text{L'ensemble des nombres décimaux : } \D = \left\{ \frac{a}{10^n} : a \in \Z, n \in \N \right\} \]
\[ \text{L'ensemble des nombres rationnels : } \Q = \left\{ \frac{a}{b}  : a \in \Z, b \in \Z^* \right\} \]
\[ \text{L'ensemble des nombres réels : } \R = ]-\infty, +\infty[ \]
\par \noindent Nous définirons l'ensemble des nombres réels plus rigoureusement dans le \autoref{chap:nb_reels}.
\\
\par \noindent Pour désigner un ensemble privé de 0, nous pouvons lui ajouter \og * \fg en exposant. 
\\ 
Par exemple $\N^* = \{1, 2, \ldots\}$.
\\
\par \noindent Définissons également certaines notations qui seront utilisées dans ce document :
\begin{itemize}
    \item $\forall$ : \og Pour tout \fg ou \og Quelque soit \fg.
    \item $\exists$ : \og Il existe \fg.
    \item $\exists!$ : \og Il existe un unique \fg.
    \item $\in$ : \og Appartient à \fg.
    \item $\subset$ : \og Inclus dans \fg.
    \item $\subsetneq$ : \og Inclus dans mais pas égal à \fg.
    \item $P \implies Q$ : \og Si $P$ alors $Q$ \fg.
    \item $P \iff Q$ : \og $P$ équivaut à $Q$ \fg. Autrement dit : \og $P$ si et seulement si $Q$ \fg.
    \item $x \coloneqq y$ : \og $x$ est défini par $y$ \fg. Pour les lecteurs informaticiens, \og $\coloneqq$ \fg se comporte comme le \og = \fg en programmation.
    \item $\equiv$ : Selon le contexte, il peut désigner plusieurs choses \cite{symbole_congru_wikipedia}:
    \begin{itemize}
        \item En arithmétique, il désigne une congruence sur des entiers.
        \item En logique, il désigne une équivalence.
        \item Sinon il désigne une \og identité \fg. C'est-à-dire une égalité qui est vraie quelque soient les valeurs des variables employées.
    \end{itemize}
    \item $\square$ : Quand il est utilisé à la fin d'une démonstration, il signifie : \og Ce qu'il fallait démontrer \fg.
    \item $\llbracket a, b \rrbracket$ : Désigne l'intervalle d'entiers entre $a$ et $b$ inclus.
    \item $[a, b]$ : Désigne l'intervalle de réels entre $a$ et $b$ inclus.
\end{itemize}
\par \noindent Pour finir nous jugeons bon de rappeler quelques raisonnements usuels.

\begin{definition}[Raisonnement par récurrence]
    Il existe plusieurs variantes du raisonnement par récurrence, définissons d'abord la récurrence simple. L'objectif est de montrer qu'une propriété $P_n$ est vraie pour tout entier naturel $n$. 
    \begin{enumerate}
        \item \textbf{Initialisation} : On montre que $P_0$ est vraie.
        \item \textbf{Hérédité} : On suppose que pour un $k$ tel que $0 < k < n,\ P_k$ est vraie et on montre que $P_{k+1}$ est vraie.
    \end{enumerate}
\end{definition}

\begin{definition}[Raisonnement par l'absurde]
    Soit $P$ une assertion. Le raisonnement par l'absurde consiste à montrer que la assertion contraire de $P$, que l'on note $\overline{P}$ dans cette définition, est fausse impliquant que $P$ est vraie.
    Pour ce faire, on suppose que $\overline{P}$ est vraie et on commence à raisonner, s'il l'on arrive à une absurdité ou une contradiction, on a montré que $\overline{P}$ est fausse, impliquant que $P$ est vraie.
\end{definition}

\part{Analyse}

\chapter{Nombres réels}\label{chap:nb_reels}
\def\arraystretch{1}

\begin{definition}[Nombre réel \cite{wikipedia_nb_reel}]
    Un nombre réel est un nombre qui peut être représenté par une partie entière et une liste finie ou infinie de décimales.
\end{definition}

\begin{proposition}[Addition et multiplication sur $\R$]
    On peut définir sur $\R$ une addition (notée \og + \fg) et une multiplication (notée \og $\times$ \fg ou \og $\cdot$ \fg) qui prolonge l'addition et la multiplication de $\N$ et vérifie les règles suivantes pour tous $a, b, c \in \R$ :
    \begin{enumerate}
        \item Commutativité : 
        \begin{enumerate}
        	\item $a + b = b + a$.
        	\item $a \cdot b = b \cdot a$.
        \end{enumerate}
        \item Associativité : 
        \begin{enumerate}
        	\item $a + (b + c) = (a + b) + c$.
        	\item $a \cdot (b \cdot c) = (a \cdot b) \cdot c$.
        \end{enumerate}
        \item Distributivité : 
		\begin{enumerate}
			\item $a \cdot (b + c) = a \cdot b + a \cdot b$.
			\item $(a + b) \cdot c = a \cdot c + b \cdot c$.
		\end{enumerate}
        \item \'Element neutre :
        \begin{enumerate}
        	\item $a + 0 = a$.
        	\item $a \cdot 1 = a$.
        \end{enumerate}
        \item \'Elément absorbant : $a \cdot 0 = 0$.
    \end{enumerate}
\end{proposition}

\begin{proposition}[Relation d'ordre sur $\R$]
    On peut définir une relation d'ordre sur $\R$, notée \og $\leq$ \fg, qui prolonge l'ordre de $\N$ et vérifie les règles suivantes pour tous $a, b, c \in \R$ :
    \begin{enumerate}
        \item Réflexivité : $a \leq a$.
        \item Antisymétrie : $a \leq b \text{ et } b \leq a \implies a = b$.
        \item Transitivité : $a \leq b \text{ et } b \leq c \implies a \leq c$.
        \item Ordre total : $a \leq b \text{ ou } b \leq a$.
        \item Compatibilité avec l'addition : $a + c \leq b + c$.
        \item Compatibilité avec la multiplication par un réel positif : $a \leq b \text{ et } c \geq 0 \implies a \cdot c \leq b \cdot c$.
    \end{enumerate}
\end{proposition}

\begin{definition}[Valeur absolue]
    Soit $x \in \R$, on définit la valeur absolue ainsi :
    \begin{align*}
        \abs{x} =
        \begin{cases}
            x &\text{ si } x \geq 0 \\
            -x &\text{ si } x < 0
        \end{cases}
    \end{align*}
\end{definition}

\begin{proposition}
   	Pour tous $a, b \in \R$ :
    \begin{enumerate}
            \item $\abs{a + b} \leq \abs{a} + \abs{b}$
            \item $\abs{a \cdot b} = \abs{a} \cdot \abs{b}$
            \item $\abs{a - b} \geq \abs{a} - \abs{b}$
            \item $\abs{a} = \sqrt{a^2}$
        \end{enumerate}
\end{proposition}

\begin{definition}[Intervalle réel]
    On dit qu'un intervalle réel $I$ est un ensemble de nombres délimité par deux réels $a$ et $b$.
    \begin{enumerate}
    	\item Pour $I = [a, b]$, on dit que $x \in I$ si et seulement si :
    	\[ a \leq x \leq b \]
    	Ici, $I$ est un \textbf{intervalle fermé}.
    	\item Pour $I = \ ]a, b[$, on dit que $x \in I$ si et seulement si :
    	\[ a < x < b \]
    	Ici, $I$ est un \textbf{intervalle ouvert}.
    	\item Pour $I = [a, b[$, on dit que $x \in I$ si et seulement si : 
    	\[ a \leq x < b \]
    	Ici, $I$ est un \textbf{intervalle semi-ouvert}.
    	\item Pour $I = \ ]a, b]$, on dit que $x \in I$ si et seulement si :
    	\[ a < x \leq b \]
    	Ici, $I$ est un \textbf{intervalle semi-ouvert}.
    \end{enumerate}
\end{definition}

\begin{definition}
    Soient $A \subseteq \R$ et $m \in \R$.
    \begin{enumerate}
        \item On dit que $m$ est un \textbf{majorant} de $A$ si et seulement si :
        \[ \forall x \in A,\ x \leq m \]
        \item On dit que $m$ est un \textbf{minorant} de $A$ si et seulement si :
        \[ \forall x \in A,\ x \geq m \]
    \end{enumerate}
    On dit que $A$ est \textbf{majorée} si elle admet un \textbf{majorant}, \textbf{minorée} si elle admet un \textbf{minorant} et \textbf{bornée} si elle est \textbf{majorée} et \textbf{minorée}.
\end{definition}

\begin{theorem}
    Soit $A \subseteq \R$ tel que $A \neq \varnothing$. 
    \begin{enumerate}
    	\item Si $A$ est \textbf{majorée}, elle admet un \textbf{plus petit majorant} appelé la \textbf{borne supérieure} de $A$, notée : $\sup(A)$.
    	\item Si $A$ est \textbf{minorée}, elle admet un \textbf{plus grand minorant} appelé la \textbf{borne inférieure} de $A$, notée : $\inf(A)$.
    \end{enumerate}
\end{theorem}

\begin{proposition}
    Soient $A \subseteq \R$ tel que $A \neq \varnothing$, $M$ un majorant de $A$ et $m$ un minorant de $A$.  
    \begin{enumerate}
        \item $M = \sup(A) \iff \forall \varepsilon > 0 : \ ]M - \varepsilon, M] \cap A \neq \varnothing$.
        \item $m = \inf(A) \iff \forall \varepsilon > 0 : [m, m + \varepsilon[ \ \cap \ A \neq \varnothing $.
    \end{enumerate}
\end{proposition}
\chapter{Fonctions réelles}
\def\arraystretch{1}

\section{Définitions}
\begin{definition}[Fonction]
	Une fonction $f$ est la donnée de :
	\begin{enumerate}
		\item Un ensemble de départ $E$.
		\item Un ensemble d'arrivée $F$.
		\item Une flèche : $ f : E \to F $ à tout élément $x \in E$ associe un élément $f(x) \in F$.
	\end{enumerate}		
	On appelle \textbf{images} les éléments de $F$ et \textbf{antécédents} les éléments de $E$.
\end{definition}

\begin{definition}[Graphe d'une fonction]
	Soient $E$ et $F$ deux ensembles et $f : E \to F$ une fonction.
	\\
	Le graphe de $f$ est défini comme :
	\[ \operatorname{Gr}(f) = \operatorname{Gph}(f) = \{ (x, f(x)),\ x \in E \} \subseteq E \times F \] 
\end{definition}

\begin{definition}
	Soit $f : E \to F$ une fonction. On dit que :
    \begin{enumerate}
        \item $f$ est \textbf{injective} si et seulement si : 
        \[ \forall x_1, x_2 \in E,\ f(x_1) = f(x_2) \implies x_1 = x_2 \]
        \item $f$ est \textbf{surjective} si et seulement si : 
        \[ \forall y \in F,\ \exists x \in E,\ f(x) = y \]
        \item $f$ est \textbf{bijective} si et seulement si :
        \[ \forall y \in E,\ \exists ! x \in E,\ f(x) = y \]
    \end{enumerate}
\end{definition}

\begin{definition}[Bijection réciproque]
	Lorsque $f : E \to F$ est une bijection. On peut définir $f^{-1} : F \to E$ la bijection réciproque de $f$ qui associe à tout élément de $F$ son unique antécédent dans $E$.
\end{definition}

\begin{proposition}
	Soient $x \in E,\ y \in F$.
    \begin{multicols}{2}
        \begin{enumerate}
            \item $f^{-1} (f(x)) = x$
            \item $f(f^{-1}(y)) = y$
        \end{enumerate}
    \end{multicols}
\end{proposition}

\begin{definition}
	Soient $I \subseteq \R$, $f : I \to \R$ et $T \in \R_+$. On dit que :
    \begin{enumerate}
        \item $f$ est \textbf{paire} si et seulement si : 
        \[ \forall x \in I,\ f(x) = f(-x) \]
        \item $f$ est \textbf{impaire} si et seulement si  :
        \[ \forall x \in I,\ -f(x) = f(-x) \]
        \item $f$ est \textbf{$T$-périodique} si et seulement si :
        \[ \forall (x, n) \in \R \times \Z,\ f(x + nT) = f(x) \]
    \end{enumerate}
\end{definition}

\begin{definition}
	Soient $I \subseteq \R,\ f : I \to \R$. On dit que :
	\begin{enumerate}
	    \item $f$ est \textbf{majorée} ou qu'elle admet un \textbf{majorant} si et seulement si : 
	    \[ \exists M \in \R,\ \forall x \in I,\ f(x) \leq M \]
        \item $f$ est \textbf{minorée} ou qu'elle admet un \textbf{minorant} si et seulement si :
        \[ \exists m \in \R,\ \forall x \in I,\ f(x) \geq m \]
        \item $f$ est \textbf{bornée} si et seulement si elle est \textbf{majorée} et \textbf{minorée}.
	\end{enumerate}
\end{definition}

\begin{definition}
	Soient $I \subseteq \R$ et $f : I \to \R$. On dit que :
    \begin{enumerate}
        \item $f$ est \textbf{croissante} si et seulement si :
        \[ \forall x, y \in I,\ x \leq y \implies f(x) \leq f(y) \]
        \item $f$ est \textbf{décroissante} si et seulement si : 
        \[ \forall x, y \in I,\ x \leq y \implies f(x) \geq f(y) \]
        \item $f$ est \textbf{monotone} si et seulement si elle est \textbf{croissante} ou \textbf{décroissante}.
        \item $f$ est \textbf{strictement croissante} si et seulement si : 
        \[ \forall x, y \in I,\ x < y \implies f(x) < f(y) \]
        \item $f$ est \textbf{strictement décroissante} si et seulement si :
        \[ \forall x, y \in I,\ x < y \implies f(x) > f(y) \]
        \item $f$ est \textbf{strictement monotone} si et seulement si elle est \textbf{strictement croissante} ou \textbf{strictement décroissante}.
    \end{enumerate}
\end{definition}

\section{Opérations sur les fonctions}
\begin{definition}[Opérations sur les fonctions]
	Soient $f, g$ deux fonctions et $A \subseteq \R$.
	\begin{center}
		$
		\appli{f+g}{A}{\R}{x}{f(x)+g(x)}
		\qquad
		\appli{fg}{A}{\R}{x}{f(x)g(x)}
		$
	\end{center}
\end{definition}

\begin{definition}[Composition de fonctions]
	Soient $E, F, G, H \subseteq \R,\ F \subseteq G,\ f : E \to F,\ g : G \to H$.
	\begin{center}
		$
		\appli{g \circ f}{E}{H}{x}{g(f(x))}
		$
	\end{center}
\end{definition}

\begin{definition}[Fonction identité]
    Soient $E, F \subseteq \R$.
	\begin{center}
		$
		\appli{id_E}{E}{F}{x}{x}
		$
	\end{center}
\end{definition}

\begin{proposition}
	Soit $f : E \to F$ une bijection.
    \begin{multicols}{2}
        \begin{enumerate}
            \item $f^{-1} \circ f = id_E$.
            \item $f \circ f^{-1} = id_F$.
        \end{enumerate}
    \end{multicols}
\end{proposition}

\begin{definition}[Image directe]
    Soient $f : E \to F,\ A \subseteq E$.
	\[ f(A) = \{ f(x),\ x \in A \},\ f(A) \subseteq F \]
\end{definition}

\begin{definition}[Image réciproque]
	Soient $f : E \to F,\ B \subseteq F$.
	\[ f^{-1}(B) = \{ \forall x \in E,\ f(x) \in B \},\ f^{-1} \subseteq E \]
\end{definition}

\begin{proposition}
    Soient $f : E \to F,\ A_1, A_2 \subseteq E,\ B_1, B_2 \subseteq F$.
    \begin{enumerate}
        \item $f(A_1 \cup A_2) = f(A_1) \cup f(A_2)$.
        \item $f^{-1} (B_1 \cup B_2) = f^{-1} (B_1) \cup f^{-1} (B_2)$.
        \item $f^{-1} (B_1 \cap B_2) = f^{-1} (B_1) \cap f^{-1} (B_2)$.
    \end{enumerate}
\end{proposition} 


\chapter{Fonctions usuelles}
\def\arraystretch{1}

\begin{definition}[Fonction polynomiale]
	Soient $n \in \N$ et $a_i \in \R$.
	\begin{center}
		$
		\appli{f}{\R}{\R}{x}{\sum_{i=0}^n a_i x^i}
		$
	\end{center}
\end{definition}

\begin{definition}[Fonction partie entière]
	\[ \forall x \in \R,\ \exists ! E(x) \in \Z,\ E(x) \leq x < E(x + 1) \]
	\begin{center}
		$
		\appli{E}{\R}{\Z}{x}{E(x)}
		$
	\end{center}
\end{definition}

\begin{figure}[!h]
	\centering
	\begin{tikzpicture}
		\draw[step=1cm, gray, very thin] (-3, -3) grid (3, 3);
		\draw[thick, ->] (-3, 0) -- (3, 0) node[anchor=north west] {$x$};
		\draw[thick, ->] (0, -3) -- (0, 3) node[anchor=south east] {$y$};
		\draw[thick, red] (1, 1) -- (2, 1);
		\draw[thick, red] (2, 2) -- (3, 2);
		\draw[thick, red] (0, 0) -- (1, 0);
		\draw[thick, red] (-1, -1) -- (0, -1);
		\draw[thick, red] (-2, -2) -- (-1, -2);
		\draw[thick, red] (-3, -3) -- (-2, -3);
		\foreach \x in {-3, -2, -1, 0, 1, 2, 3}
			\draw (\x cm, 1pt) -- (\x cm, -1pt) node[anchor=north] {$\x$};
		\foreach \y in {-3, -2, -1, 0, 1, 2, 3}
			\draw (1pt, \y cm) -- (-1pt, \y cm) node[anchor=east] {$\y$};
	\end{tikzpicture}
	\caption{Graphe de la fonction partie entière}
\end{figure}

\begin{definition}[Fonction puissance]
	Soit $a \in \R$.
	\begin{center}
		$
		\appli{f}{\R_+^*}{\R}{x}{x^a}
		$
	\end{center}
\end{definition}

\begin{proposition}
	Soient $a,\ b,\ x \in \R$.
    \begin{enumerate}
            \item $1^a = 1$.
            \item $x^a \cdot x^b = x^{a + b}$.
            \item $(xy)^a = x^a y^a$.
            \item $(x^a)^b = x^{ab}$.
        \end{enumerate}
\end{proposition}

\section{Fonctions trigonométriques}
\begin{definition}[Fonctions trigonométriques]
	Soit $k \in \Z$ :	
	\begin{center}
		$
		\appli{\cos}{\R}{[-1,1]}{x}{\cos(x)}
		\qquad
		\appli{\sin}{\R}{[-1,1]}{x}{\sin(x)}
		\qquad
		\appli{\tan}{\R \backslash \left\{ \frac{k \pi}{2} \right\}}{\R}{x}{\frac{\sin(x)}{\cos(x)}}
		$
	\end{center}	
\end{definition}

\begin{figure}[h!]
	\centering
	%\includegraphics{images/cos.png}
	\begin{subfigure}{0.3\textwidth}
\begin{tikzpicture}
            \begin{axis}[domain=-5:5, axis lines=middle, grid=both, xlabel=$x$, ylabel=$\cos(x)$, width=6cm]
                \addplot[color=green, samples=100] {cos(deg(x))};
            \end{axis}
    \end{tikzpicture}
		
		\caption{Fonction $\cos{x}$}
	\end{subfigure}
	\begin{subfigure}{0.3\textwidth}
\begin{tikzpicture}
            \begin{axis}[domain=-5:5, axis lines=middle, grid=both, xlabel=$x$, ylabel=$\sin(x)$, width=6cm]
                \addplot[color=green, samples=100] {sin(deg(x))};
            \end{axis}
    \end{tikzpicture}
		\caption{Fonction $\sin{x}$}
	\end{subfigure}
	\begin{subfigure}{0.3\textwidth}
\begin{tikzpicture}
            \begin{axis}[domain=-10:10, restrict y to domain=-10:10, axis lines=middle, grid=both, xlabel=$x$, ylabel=$\tan(x)$, width=6cm]
                \addplot[color=green, samples=100] {tan(deg(x))};
            \end{axis}
    \end{tikzpicture}
		\caption{Fonction $\tan{x}$}
	\end{subfigure}
\end{figure}

\begin{figure}[h!]
	\centering
	\begin{tikzpicture}[scale=3,cap=round,>=latex]
		% draw the coordinates
		\draw[->] (-1.5cm,0cm) -- (1.5cm,0cm) node[right,fill=white] {$x$};
		\draw[->] (0cm,-1.5cm) -- (0cm,1.5cm) node[above,fill=white] {$y$};
		
		% draw the unit circle
		\draw[thick] (0cm,0cm) circle(1cm);
		
		\foreach \x in {0,30,...,360} {
			% lines from center to point
			\draw[gray] (0cm,0cm) -- (\x:1cm);
			% dots at each point
			\filldraw[black] (\x:1cm) circle(0.4pt);
			% draw each angle in degrees
			\draw (\x:0.6cm) node[fill=white] {$\x^\circ$};
		}
		
		% draw each angle in radians
		\foreach \x/\xtext in {
			30/\frac{\pi}{6},
			45/\frac{\pi}{4},
			60/\frac{\pi}{3},
			90/\frac{\pi}{2},
			120/\frac{2\pi}{3},
			135/\frac{3\pi}{4},
			150/\frac{5\pi}{6},
			180/\pi,
			210/\frac{7\pi}{6},
			225/\frac{5\pi}{4},
			240/\frac{4\pi}{3},
			270/\frac{3\pi}{2},
			300/\frac{5\pi}{3},
			315/\frac{7\pi}{4},
			330/\frac{11\pi}{6},
			360/2\pi}
		\draw (\x:0.85cm) node[fill=white] {$\xtext$};
		
		\foreach \x/\xtext/\y in {
			% the coordinates for the first quadrant
			30/\frac{\sqrt{3}}{2}/\frac{1}{2},
			45/\frac{\sqrt{2}}{2}/\frac{\sqrt{2}}{2},
			60/\frac{1}{2}/\frac{\sqrt{3}}{2},
			% the coordinates for the second quadrant
			150/-\frac{\sqrt{3}}{2}/\frac{1}{2},
			135/-\frac{\sqrt{2}}{2}/\frac{\sqrt{2}}{2},
			120/-\frac{1}{2}/\frac{\sqrt{3}}{2},
			% the coordinates for the third quadrant
			210/-\frac{\sqrt{3}}{2}/-\frac{1}{2},
			225/-\frac{\sqrt{2}}{2}/-\frac{\sqrt{2}}{2},
			240/-\frac{1}{2}/-\frac{\sqrt{3}}{2},
			% the coordinates for the fourth quadrant
			330/\frac{\sqrt{3}}{2}/-\frac{1}{2},
			315/\frac{\sqrt{2}}{2}/-\frac{\sqrt{2}}{2},
			300/\frac{1}{2}/-\frac{\sqrt{3}}{2}}
		\draw (\x:1.25cm) node[fill=white] {$\left(\xtext,\y\right)$};
		
		% draw the horizontal and vertical coordinates
		% the placement is better this way
		\draw (-1.25cm,0cm) node[above=1pt] {$(-1,0)$}
		(1.25cm,0cm)  node[above=1pt] {$(1,0)$}
		(0cm,-1.25cm) node[fill=white] {$(0,-1)$}
		(0cm,1.25cm)  node[fill=white] {$(0,1)$};
	\end{tikzpicture}
	\caption{Cercle trigonométrique}
\end{figure}

\begin{proposition}
	Soient $x \in \R$ et $k \in \Z$.
	\begin{enumerate}
		\item $\cos(-x) = \cos(x)$.
		\item $\sin(-x) = -\sin(x)$.
		\item $\tan(-x) = -\tan(x)$.
		\item $\cos(x + 2k\pi) = \cos(x)$.
		\item $\sin(x + 2k\pi) = \sin(x)$.
		\item $\tan(x + k\pi) = \tan(x)$.
		\item $\cos$ est bijective sur $[0,\ \pi]$.
		\item $\sin$ est bijective sur $[-\frac{\pi}{2},\ \frac{\pi}{2}]$.
		\item $\tan$ est bijective sur $]-\frac{\pi}{2},\ \frac{\pi}{2}[$.
	\end{enumerate}
\end{proposition}

\begin{definition}
	On définit $\arccos,\ \arcsin$ et $\arctan$ comme étant les bijections réciproques des fonctions $\cos,\ \sin$ et $\tan$.
\end{definition}

\begin{figure}[h!]
	\centering
	\begin{tikzpicture}
		\begin{axis}[domain=-1:1, axis lines=middle, grid=both, xlabel=$x$, ylabel=$y$, samples=100, legend pos=north west]
			\addplot[color=green]{x};
			\addlegendentry{$x$}
			\addplot[color=red]  {cos(deg(x))};
			\addlegendentry{$\cos(x)$}
			\addplot[color=blue]  {rad(acos(x))};
			\addlegendentry{$\arccos(x)$}
		\end{axis}
	\end{tikzpicture}
	\begin{tikzpicture}
		\begin{axis}[domain=-1:1, axis lines=middle, grid=both, xlabel=$x$, ylabel=$y$, samples=100, legend pos=north west]
			\addplot[color=green]{x};
			\addlegendentry{$x$};
			\addplot[color=red]  {sin(deg(x))};
			\addlegendentry{$\sin(x)$};
			\addplot[color=blue]  {deg(asin(x))};
			\addlegendentry{$\arcsin(x)$};
		\end{axis}
	\end{tikzpicture}
	\begin{tikzpicture}
		\begin{axis}[xmin=-5, xmax=5, ymin=-10, ymax=10, axis lines=middle, grid=both, xlabel=$x$, ylabel=$y$, samples=100, legend pos=north west]
			\addplot[color=green]{x};
			\addlegendentry{$x$};
			\addplot[color=red]  {tan(deg(x))};
			\addlegendentry{$\tan(x)$};
			\addplot[color=blue]  {rad(atan(x))};
			\addlegendentry{$\arctan(x)$};
		\end{axis}
	\end{tikzpicture}
\end{figure}

\begin{proposition}
    \begin{enumerate}
    	\item $\forall x \in [-\frac{\pi}{2},\ \frac{\pi}{2}],\ \arcsin(\sin(x)) = x$.
        \item $\forall x \in ]-\frac{\pi}{2},\ \frac{\pi}{2}[,\ \arctan(\tan(x)) = x$.
        \item $\forall x \in [0,\ \pi],\ \arccos(\cos(x)) = x$.
        \item $\forall x \in [-1,\ 1]$ :
        \begin{enumerate}
                \item $\sin(\arcsin(x)) = x$.
                \item $\cos(\arccos(x)) = x$.
            \end{enumerate}
        \item $\forall x \in \R,\ \tan(\arctan(x)) = x$.
    \end{enumerate}
\end{proposition}

\begin{proposition}
	Soient $a,\ b \in \R$.
        \begin{enumerate}
            \item $\sin(a + b) = \sin(a) \cos(b) + \sin(b) \cos(a)$.
            \item $\cos(a + b) = \cos(a) \cos(b) - \sin(a) \sin(b)$.
        \end{enumerate}
\end{proposition}

\begin{proof}
    Nous pouvons procéder avec des produits scalaires, mais nous allons utiliser les nombres complexes ici.
    \\
    D'une part :
    \[e^{i (a + b)} = \cos(a + b) + i\sin(a + b)\]
    D'autre part :
    \begin{align*}
        e^{i (a + b)} &= e^{ia} \cdot e^{ib} \\
        &= [\cos(a) + i \sin(a)] \cdot [\cos(b) + i \sin(b)] \\
        &= \cos(a) \cos(b) + i\sin(b)\cos(a) + i\sin(a)\cos(b) - \sin(a)\sin(b) \\
        &= \cos(a)\cos(b) - \sin(a)\sin(b) + i [\sin(b)\cos(a) + \sin(a) \cos(b)].
    \end{align*}
    Par identification de la partie réelle et de la partie imaginaire :
    \[ \sin(a + b) = \sin(a) \cos(b) + \sin(b) \cos(a) \]
	\[ \cos(a + b) = \cos(a) \cos(b) - \sin(a) \sin(b) \]
\end{proof}

\begin{proposition}
    $\forall x \in \R,\ \cos^2(x) + \sin^2(x) = 1$.
\end{proposition}

\begin{proof}
    C'est une application du théorème de Pythagore sachant que le rayon du cercle trigonométrique est égal à 1.
\end{proof}

\begin{proposition}
	
    \begin{enumerate}
        \item $\lim_{x \to +\infty} \arctan(x) = \frac{\pi}{2}$.
        \item $\lim_{x \to -\infty} \arctan(x) = -\frac{\pi}{2}$.
    \end{enumerate}
\end{proposition}

\section{Exponentielle et logarithme}

\begin{definition}[Fonction exponentielle]
	\begin{center}
		$
		\appli{\exp}{\R}{\R_+^*}{x}{\exp(x) \equiv e^x}
		$
	\end{center}
\end{definition}

\begin{proposition}
	La fonction exponentielle est \textbf{bijective} et \textbf{strictement croissante} et $\exp(0) = 1$.
    \begin{enumerate}
            \item $\lim_{x \to -\infty} e^x = 0$.
            \item $\lim_{x \to +\infty} e^x = +\infty$.
        \end{enumerate}
    \noindent Soient $x,\ y \in \R$.
    \begin{enumerate}
            \item $\exp(x + y) = \exp(x) \cdot \exp(y)$.
            \item $\exp(-x) = \frac{1}{\exp(x)}$.
            \item $\exp(x - y) = \frac{\exp(x)}{\exp(y)}$.
        \end{enumerate}
\end{proposition}

\begin{definition}[Logarithme néperien]
	\begin{center}
		$
		\appli{\ln}{\R_+^*}{\R}{x}{\ln(x)}
		$
	\end{center}
\end{definition}

\begin{figure}[!h]
	\centering
	\begin{subfigure}{0.45\textwidth}
		\begin{tikzpicture}
			\begin{axis}[domain=0:5, axis lines=middle, grid=both, xlabel=$x$, ylabel=$e^x$, width=7cm]
				\addplot[color=red, samples=50]{e^x};
			\end{axis}
		\end{tikzpicture}  
		\caption{Fonction exponentielle}
	\end{subfigure}
	\begin{subfigure}{0.45\textwidth}
		\begin{tikzpicture}
			\begin{axis}[domain=0:10, axis lines=middle, grid=both, xlabel=$x$, ylabel=$\ln(x)$, width=7.5cm]
				\addplot[color=red, samples=100]{ln(x)};
			\end{axis}
		\end{tikzpicture}
		\caption{Fonction logarithme népérien}
	\end{subfigure}
\end{figure}

\begin{proposition}
	Soient $x \in \R_+^*,\ y \in \R$.
    \begin{multicols}{2}
        \begin{enumerate}
            \item $\exp(\ln(x)) = x$.
            \item $\ln(\exp(y)) = y$.
        \end{enumerate}
    \end{multicols}
\end{proposition}

\begin{proposition}
	La fonction logarithme néperien est \textbf{bijective} et \textbf{strictement croissante} et $\ln(1) = 0$.
    \begin{enumerate}
            \item $\lim_{x \to 0} \ln(x) = -\infty$.
            \item $\lim_{x \to +\infty} \ln(x) = +\infty$.
        \end{enumerate}
    \noindent Soient $x,\ y \in \R$.
    \begin{enumerate}
            \item $\ln(xy) = \ln(x) + \ln(y)$.
            \item $\ln(\frac{1}{x}) = -\ln(x)$.
            \item $\ln(\frac{x}{y}) = \ln(x) - \ln(y)$.
            \item $\ln(x^n) = n\ln(x)$.
        \end{enumerate}
\end{proposition}

\section{Fonctions hyperboliques}

\begin{definition}[Fonctions hyperboliques]
	
	\begin{center}
		$
		\appli{\cosh}{\R}{\R}{x}{\frac{e^x + e^{-x}}{2}}
		\qquad 
		\appli{\sinh}{\R}{\R}{x}{\frac{e^x - e^{-x}}{2}}
		\qquad
		\appli{\tanh}{\R}{\R}{x}{\frac{\sinh(x)}{\cosh(x)} = \frac{e^x - e^{-x}}{e^x + e^{-x}}}
		$
	\end{center}
\end{definition}

\begin{figure}[!h]
	\centering
	\begin{subfigure}{0.3\textwidth}
		%\includegraphics[width=\textwidth, height=3.5cm]{images/ch.png}
		
		\begin{tikzpicture}
			\begin{axis}[domain=-5:5, axis lines=middle, grid=both, xlabel=$x$, ylabel=$\cosh(x)$, width=5cm]
				\addplot[color=red, samples=100]  {cosh(x)};
			\end{axis}
		\end{tikzpicture}
		\caption{Fonction cosh}
	\end{subfigure}
	\begin{subfigure}{0.3\textwidth}
		%\includegraphics[width=\textwidth, height=3.5cm]{images/sh.png}
		\begin{tikzpicture}
			\begin{axis}[domain=-5:5, axis lines=middle, grid=both, xlabel=$x$, ylabel=$\sinh(x)$, width=5cm]
				\addplot[color=red, samples=100]  {sinh(x)};
			\end{axis}
		\end{tikzpicture}
		\caption{Fonction sinh}
	\end{subfigure}
	\begin{subfigure}{0.3\textwidth}
		%\includegraphics[width=\textwidth, height=3.5cm]{images/tanh.png}
		\begin{tikzpicture}
			\begin{axis}[domain=-5:5, axis lines=middle, grid=both, xlabel=$x$, ylabel=$\tanh(x)$, width=5cm]
				\addplot[color=red, samples=100]  {tanh(x)};
			\end{axis}
		\end{tikzpicture}
		\caption{Fonction tanh}
	\end{subfigure}
\end{figure}

\begin{proposition}
	Soit $x \in \R$.
	\begin{enumerate}
		\item $\cosh(-x) = \cosh(x)$.
		\item $\sinh(-x) = \sinh(x)$.
		\item $\tanh(-x) = \tanh(x)$.
	\end{enumerate}
\end{proposition}

\begin{proposition}
	Soient $x,\ y \in \R$.
    \begin{enumerate}
        \item $\cosh(x + y) = \cosh(x) \cosh(y) + \sinh(x) \sinh(y)$.
        \item $\sinh(x + y) = \cosh(x) \sinh(y) + \sinh(x) \cosh(y)$.
    \end{enumerate}
\end{proposition}

\begin{proof}
    Calcul direct avec les définitions de $\cosh$ et de $\sinh$.
\end{proof}

\chapter{Suites réelles}
\def\arraystretch{1}

\section{Définitions}

\begin{definition}[Suite réelle]
    On appelle \textbf{suite réelle} une fonction de $\N \longrightarrow \R$. 
    On note $(u_n)_{n \in \N}$ la fonction $x \longmapsto u_n$.
\end{definition}

\begin{definition}
	Soit $(u_n)_{n \in \N}$ une suite.
	\begin{enumerate}
		\item $(u_n)_{n \in \N}$ est dite \textbf{croissante} si et seulement si pour tout $n \in \N$ :
		\[ u_n \leq u_{n+1} \]
		\item $(u_n)_{n \in \N}$ est dite \textbf{décroissante} si et seulement si pour tout $n \in \N$ :
		\[ u_n \geq u_{n+1} \]
		\item $(u_n)_{n \in \N}$ est dite \textbf{monotone} si et seulement si elle est \textbf{croissante} ou \textbf{décroissante}.
		\item $(u_n)_{n \in \N}$ est dite \textbf{strictement croissante} si et seulement si pour tout $n \in \N$ :
		\[ u_n < u_{n+1} \]
		\item $(u_n)_{n \in \N}$ est dite \textbf{strictement décroissante} si et seulement si pour tout $n \in \N$ :
		\[ u_n > u_{n+1} \]
		\item $(u_n)_{n \in \N}$ est dite \textbf{strictement monotone} si et seulement si elle est \textbf{strictement croissante} ou \textbf{strictement décroissante}.
		\item $(u_n)_{n \in \N}$ est dite \textbf{constante} si et seulement si pour tout $n \in \N$ :
		\[ u_n = u_0 \]
		\item $(u_n)_{n \in \N}$ est dite \textbf{stationnaire} si et seulement si elle est \textbf{constante} à partir d'un certain rang. 
		\[ \exists N \in \N,\ \forall n \geq N : u_n = u_N \]
	\end{enumerate}
\end{definition}

\section{Suites usuelles}
\begin{definition}[Suite arithmétique]
    Soit $r \in \R^*$.
    On définit une \textbf{suite arithmétique} $(u_n)_{n \in \N}$ de la manière suivante :
    \begin{align*}
        (u_n)_{n \in \N} =
        \begin{cases}
            u_{n+1} = u_n + r \\ 
            u_n = u_0 + nr
        \end{cases}
    \end{align*}
\end{definition}

\begin{proposition}
    Soient $(u_n)_{n \in \N}$ une suite arithmétique et $r \in \R^*$.
    \begin{align*}
        \sum_{k = 0}^{n} u_k = (n+1)u_0 + r \frac{n(n+1)}{2}
    \end{align*}
\end{proposition}

\begin{definition}[Suite géométrique]
    Soit $q \in \R^*$.
    On définit une \textbf{suite géométrique} $(u_n)_{n \in \N}$ de la manière suivante : 
    \begin{align*}
        (u_n)_{n \in \N} = 
        \begin{cases}
            u_{n+1} = qu_n \\
            u_n = u_0 q^n
        \end{cases}
    \end{align*}
\end{definition}

\begin{proposition}
    Soient $(u_n)_{n \in \N}$ une suite géométrique et $q \in \R^*$.
    \begin{align*}
        \sum_{k = 0}^{n} u_k = 
        \begin{cases}
            u_0 \frac{1 - q^{n+1}}{1 - q} \text{ si } q \neq 1\\
            u_0(n+1) \text{ si } q = 1
        \end{cases}
    \end{align*}
\end{proposition}

\begin{definition}[Suite arithmético-géométrique]
    Soient $r, q \in \R^*$. 
    \\
    On définit une \textbf{suite arithmético-géométrique} $(u_n)_{n \in \N}$ de la manière suivante :
    \begin{align*}
        (u_n)_{n \in \N} = 
        \begin{cases}
            u_{n+1} = qu_n + r \\ 
            u_n = a + (u_0 - a) q^n,\ a = \frac{r}{1 - q}
        \end{cases}
    \end{align*}
\end{definition}

\begin{enumerate}
	\item On commence par résoudre $a = qa + r \iff a = \frac{r}{1 - q}$.
	\item Ensuite, on pose $v_n = u_n - a$ et $v_{n+1} = u_{n+1} - a$ qui est une suite géométrique, ainsi $v_n = v_0 q^n$.
	\item Finalement $u_n = v_n + a \iff u_n = v_0 q^n + a \iff u_n = (u_0 - a) q^n + a$.
\end{enumerate}

\begin{example}
	Soit $(u_n)_{n \in \N}$ définie par :
	\[
	\begin{cases}
		u_0 &= 3 \\
		u_{n+1} &= \frac{1}{4} u_n + 3
	\end{cases}
	\]
	\begin{enumerate}
		\item On cherche tout d'abord à résoudre $a = \frac{1}{4} a + 3$.
		\begin{align*}
			a &= \frac{1}{4} a + 3 \\
			\frac{3}{4} a &= 3 \\
			a &= 4
		\end{align*}
		\item Posons $v_n = u_n - 4$, on a $v_{n+1} = u_{n+1} - 4$.
		\begin{align*}
			v_{n+1} &= u_{n+1} - 4 \\
			v_{n+1} &= \frac{1}{4} u_n + 3 - 4 \\
			v_{n+1} &= \frac{1}{4} u_n - 1 \\
			v_{n+1} &= \frac{1}{4} \left( u_n - 4 \right) \\
			v_{n+1} &= \frac{1}{4} v_n
		\end{align*}
		Ainsi on a montré que $v_n$ est une suite géométrique de raison $\frac{1}{4}$.
		\\
		On a $v_0 = u_0 - 4 = 3 - 4 = -1$, puis :
		\[ v_n = -\frac{1}{4^n} \]
		\item Finalement, $u_n = v_n + 4 = 4 - \frac{1}{4^n}$.
	\end{enumerate}
\end{example}

\section{Convergence d'une suite}
\begin{definition}
    Soient $(u_n)_{n \in \N}$ une suite et $\ell \in \R$. 
    \\
    On dit que $(u_n)_{n \in \N}$ \textbf{tend} vers $\ell$ si et seulement si :
    \[ \forall \varepsilon > 0,\ \exists N \in \N,\ \forall n \geq N : \abs{u_n - \ell} \leq \varepsilon \]
    On note alors :
    \begin{align*}
        u_n \xrightarrow[n \to +\infty]{} \ell \text{ ou } \lim_{n \to +\infty} u_n = \ell
    \end{align*}
\end{definition}

\begin{definition}
    Soit $(u_n)_{n \in \N}$ une suite réelle. On dit que :
    \begin{enumerate}
        \item $(u_n)_{n \in \N}$ \textbf{converge} si et seulement si : 
        \[ \exists \ell \in \R,\ \forall \varepsilon > 0,\ \exists N \in \N,\ \forall n \geq N : \abs{u_n - \ell} \leq \varepsilon \]
        \item $(u_n)_{n \in \N}$ \textbf{diverge} si et seulement si : 
        \[ \forall \ell \in \R,\ \exists \varepsilon > 0,\ \forall N \in \N,\ \exists n \geq N : \abs{u_n - \ell} > \varepsilon \]
        \item $(u_n)_{n \in \N}$ \textbf{ne converge pas vers} $\ell \in \R$ si et seulement si : 
        \[ \exists \varepsilon > 0,\ \forall N \in \N,\ \exists n \geq N : \abs{u_n - \ell} > \varepsilon \]
    \end{enumerate}
\end{definition}

\begin{theorem}
    La limite d'une suite convergente $(u_n)_{n \in \N}$ est \textbf{unique}.
\end{theorem}

\begin{proof}
    Procédons à un raisonnement par l'absurde.
    \\
    On suppose que $\ell_1 \neq \ell_2$. Posons $\varepsilon = \frac{1}{3} \abs{\ell_1 - \ell_2} > 0$.
    \\
    Par définition de la limite :
    \begin{align*}
        &\exists N_1 \in \N,\ \forall n \geq N_1,\ \abs{u_n - \ell_1} \leq \varepsilon 
        &
        &\exists N_2 \in \N,\ \forall n \geq N_2,\ \abs{u_n - \ell_2} \leq \varepsilon
    \end{align*}
    Posons $N = \max(N_1, N_2)$, si $n \geq N$ alors :
    \begin{align*}
        \abs{u_n - \ell_1} \leq \varepsilon \text{ et } \abs{u_n - \ell_2} \leq \varepsilon
    \end{align*}
    \[ \abs{\ell_1 - \ell_2} = \abs{\ell_1 - u_n + u_n - \ell_2} \leq \abs{u_n - \ell_1} + \abs{u_n - \ell_2} \]
    Alors : 
    \begin{align*}
        \abs{u_n - \ell_1} + \abs{u_n - \ell_2} &\leq \varepsilon + \varepsilon = \frac{2}{3} \abs{\ell_1 - \ell_2} \\
        \abs{\ell_1 - \ell_2} &\leq \frac{2}{3} \abs{\ell_1 - \ell_2} \\
        \frac{1}{3} \abs{\ell_1 - \ell_2} \leq 0 \\
        \varepsilon \leq 0
    \end{align*}
    ce qui est absurde. Ainsi on a montré que $\ell_1 = \ell_2$.
\end{proof}

\begin{theorem}
    Toute suite convergente est \textbf{bornée}.
\end{theorem}

\begin{proof}
    Supposons qu'une suite $(u_n)_{n \in \N}$ converge vers $\ell \in \R$.
    Posons $\varepsilon = 1$. 
    \\
    Par définition de la limite :
    \[ \exists N \in \N,\ \forall n \geq N,\ \abs{u_n - \ell} \leq 1 \iff \ell - 1 \leq u_n \leq \ell + 1 \]
    Posons $M = \max(u_0, \ldots, u_{N-1}, \ell + 1)$ et $m = \min(u_0, \ldots, u_{N - 1}, \ell - 1)$.
    \begin{align*}
        \forall n \in \N,\ 
        \begin{cases}
            m \leq u_n \leq M, &\text{ si } n < N \\
            \ell - 1 \leq u_n \leq \ell + 1, &\text{ si } n > N
        \end{cases}
    \end{align*}
    Sachant que $m \leq \ell - 1$ et $M \geq \ell + 1$.
    On a :
    \[ \forall n \in \N,\ m \leq u_n \leq M \]
    ce qui signifie que $(u_n)_{n \in \N}$ est bornée.
\end{proof}

\begin{theorem}
    Soient $\ell_1, \ell_2 \in \R$, $(u_n)_{n \in \N}$ et $(v_n)_{n \in \N}$ deux suites convergentes telles que :
    \[ \lim_{n \to +\infty} u_n = \ell_1 \text{ et } \lim_{n \to +\infty} v_n = \ell_2 \]
    Alors :
    \[ \lim_{n \to +\infty} (u_n + v_n) = \ell_1 + \ell_2 \]
\end{theorem}

\begin{proof}
    Soit $\varepsilon > 0$.
    \\
    \\
    Par définition de la limite :
    \begin{align*}
        &\exists N_1 \in \N,\ \forall n \geq N_1,\ \abs{u_n - \ell_1} \leq \varepsilon &
        &\exists N_2 \in \N,\ \forall n \geq N_2,\ \abs{v_n - \ell_2} \leq \varepsilon
    \end{align*}
    Posons $N = \max(N_1, N_2)$. Si $n \geq N$, alors :
    \begin{align*}
        \abs{u_n - \ell_1} \leq \varepsilon \text{ et } \abs{v_n - \ell_2} \leq \varepsilon
    \end{align*}
    Puis : 
    \begin{align*}
        \abs{u_n + v_n - (\ell_1 + \ell_2)} = \abs{u_n - \ell_1 + v_n - \ell_2} &\leq \abs{u_n - \ell_1} + \abs{v_n - \ell_2} \\
        \abs{u_n - \ell_1} + \abs{v_n - \ell_2} &\leq \varepsilon + \varepsilon = 2 \varepsilon
    \end{align*}
    Posons $\varepsilon' = 2\varepsilon$.
    Ainsi : 
    \[ \abs{u_n - \ell_1} + \abs{v_n - \ell_2} \leq \varepsilon' \]
    C'est-à-dire que $\lim_{n \to +\infty} (u_n + v_n) = \ell_1 + \ell_2$.
\end{proof}

\begin{theorem}
    Soient $\ell_1, \ell_2 \in \R$, $(u_n)_{n \in \N}$ et $(v_n)_{n \in \N}$ deux suites convergentes telles que :
    \[ \lim_{n \to + \infty} u_n = \ell_1 \text{ et } \lim_{n \to +\infty} v_n = \ell_2 \]
    Alors :
    \[ \lim_{n \to +\infty} (u_n \cdot v_n) = \ell_1 \cdot \ell_2 \]
\end{theorem}

\begin{proof}
    Soit $\varepsilon > 0$.
    \\
    Comme $(u_n)_{n \in \N}$ converge, elle est bornée :
    \[ \exists M \in \R, \forall n \in \N,\ \abs{u_n} \leq M \]
    Par définition de la limite :
    \begin{align*}
        &\exists N_1,\ \forall n \geq N_1,\ \abs{u_n - \ell_1} \leq \varepsilon &
        &\exists N_2,\ \forall n \geq N_2,\ \abs{v_n - \ell_2} \leq \varepsilon
    \end{align*}
    Posons $N = \max(N_1, N_2)$. Si $n \geq N$ alors :
    \[ \abs{u_n - \ell_1} \leq \varepsilon \text{ et } \abs{v_n - \ell_2} \leq \varepsilon \]
    Puis : 
    \begin{align*}
        \abs{u_n \cdot v_n - \ell_1 \cdot \ell_2} &= \abs{u_n \cdot v_n - u_n \cdot \ell_2 + u_n \cdot \ell_2 - \ell_1 \cdot \ell_2} \\
        &= \abs{u_n(v_n - \ell_2) + \ell_2 (u_n - \ell_1)} \\ 
        &\leq \abs{u_n} \abs{v_n - \ell_2} + \abs{\ell_2} \abs{u_n - \ell_1} \\ 
        &\leq M \varepsilon + \abs{\ell_2} \varepsilon = (M + \abs{\ell_2}) \varepsilon
    \end{align*}
    Posons $\varepsilon' = (M + \abs{\ell_2}) \varepsilon$. Ainsi :
    \[ \abs{u_n \cdot v_n - \ell_1 \cdot \ell_2} \leq \varepsilon' \]
    C'est-à-dire que $\lim_{n \to +\infty} (u_n \cdot v_n) = \ell_1 \cdot \ell_2$.
\end{proof}

\begin{theorem}
    Soient $(u_n)_{n \in \N}$ et $(v_n)_{n \in \N}$ deux suites convergentes telles que pour tout $n \in \N,\ u_n \leq v_n$. 
    \[ \lim_{n \to +\infty} u_n \leq \lim_{n \to +\infty} v_n \]
\end{theorem}

\begin{proof}
    Soient $\ell_1, \ell_2 \in \R$.
    Posons :
    \begin{align*}
        \ell_1 &= \lim_{n \to +\infty} u_n & \ell_2 &= \lim_{n \to +\infty} v_n
    \end{align*}
    On raisonne par l'absurde en supposant que $\ell_1 > \ell_2$.
    \\
    Posons : 
    \[ \varepsilon = \frac{\ell_1 - \ell_2}{3} > 0 \]
    \begin{align*}
        &\exists N_1,\ \forall n \geq N_1,\ \abs{u_n - \ell_1} \leq \varepsilon & 
        &\exists N_2,\ \forall n \geq N_2,\ \abs{v_n - \ell_2} \leq \varepsilon
    \end{align*}
    Autrement dit :
    \begin{align*}
        &\forall n \geq N_1,\ u_n \geq \ell_1 - \varepsilon &
        &\forall n \geq N_2,\ v_n \leq \ell_2 + \varepsilon
    \end{align*}
    Posons $N = \max(N_1, N_2)$. Si $n \geq N$ alors :
    \begin{align*}
        v_n \leq \ell_2 + \varepsilon < \ell_1 - \varepsilon \leq u_n \implies v_n < u_n
    \end{align*}
    ce qui est absurde. Ainsi $\ell_1 \leq \ell_2$.
\end{proof}

\begin{corollary}
    Soient $(u_n)_{n \in \N}$, $\ell \in \R$ tels que $\lim_{n \to +\infty} = \ell$ et $m, M \in \R$.
    \begin{enumerate}
            \item Si pour tout $n \in \N,\ u_n \leq M$ alors $\ell \leq M$.
            \item Si pour tout $n \in \N,\ u_n \geq m$ alors $\ell \geq m$.
    \end{enumerate}
\end{corollary}

\begin{theorem}[Théorème des gendarmes]
    Soient $\ell \in \R$, $(u_n)_{n \in \N}, (v_n)_{n \in \N}, (w_n)_{n \in \N}$ des suites telles que pour tout $n \in \N : u_n \leq v_n \leq w_n $.
    \\
    Si $\lim_{n \to +\infty} u_n = \ell$ et $\lim_{n \to +\infty} w_n = \ell$ alors $\lim_{n \to + \infty} v_n = \ell$.
\end{theorem}

\begin{proof}
    \begin{align*}
        &\exists N_1 \in \R,\ \forall n \geq N_1,\ \abs{u_n - \ell} \geq \varepsilon &
        &\exists N_2 \in \R,\ \forall n \geq N_2,\ \abs{w_n - \ell} \geq \varepsilon
    \end{align*}
    Posons $N = \max(N_1, N_2)$. Si $n \geq N$ alors :
    \begin{align*}
        \abs{u_n - \ell} \leq \varepsilon \text{ et } \abs{w_n - \ell} \leq \varepsilon
    \end{align*}
    ce qui revient à dire : 
    \begin{align*}
        &\ell - \varepsilon \leq u_n - \ell \leq \ell + \varepsilon & &\ell - \varepsilon \leq w_n - \ell \leq \ell + \varepsilon
    \end{align*}
    Sachant que :
    \[ u_n \leq v_n \leq w_n \]
    on a :
    \begin{align*}
        \ell - \varepsilon \leq u_n - \ell \leq v_n - \ell \leq w_n - \ell \leq \ell + \varepsilon
    \end{align*}
    et donc finalement :
    \[ \ell - \varepsilon \leq v_n - \ell \leq \ell + \varepsilon \iff \abs{v_n - \ell} \leq \varepsilon \]
    C'est-à-dire $\lim_{n \to +\infty} v_n = \ell$.
\end{proof}

\begin{example}
	Cherchons à déterminer $\lim_{x \to +\infty} \frac{\sin(x)}{x}$.
	\\
	On sait tout d'abord que :
	\[ -1 \leq \sin(x) \leq 1 \]
	puis :
	\[ \frac{-1}{x} \leq \frac{\sin(x)}{x} \leq \frac{1}{x} \]
	On a d'une part :
	\[ \lim_{x \to +\infty} \frac{-1}{x} = 0 \]
	d'autre part :
	\[ \lim_{x \to +\infty} \frac{1}{x} = 0 \]
	par le théorème des gendarmes :
	\[ \lim_{x \to + \infty} \frac{\sin(x)}{x} = 0 \]
\end{example}

\begin{theorem}
    \begin{enumerate}
        \item Toute suite croissante majorée converge.
        \item Toute suite décroissante minorée converge.
    \end{enumerate}
\end{theorem}

\begin{theorem}[Théorème des suites adjacentes]
    Soient $(u_n)_{n \in \N}$ et $(v_n)_{n \in \N}$ deux suites telles que :
    \begin{enumerate}
            \item $(u_n)_{n \in \N}$ est croissante.
            \item $(v_n)_{n \in \N}$ est décroissante.
            \item $\lim_{n \to +\infty} (v_n - u_n) = 0$
        \end{enumerate}
    \noindent Alors $(u_n)_{n \in \N}$ et $(v_n)_{n \in \N}$ convergent vers la même limite.
\end{theorem}

\begin{proof}
    Posons $w_n = v_n - u_n$.
    \\ 
    On sait que $(v_n)_{n \in \N}$ est décroissante et que $(u_n)_{n \in \N}$ est croissante. 
    \\ 
    Ainsi $v_{n+1} - v_n \leq 0$ et $u_{n+1} - u_n \geq 0$.
    Etudions la variation de $(w_n)_{n \in \N}$.
    \begin{align*}
        w_{n+1} - w_n &= v_{n+1} - v_n - (u_{n+1} - u_n) < 0
    \end{align*}
    Ainsi $(w_n)_{n \in \N}$ est décroissante et sa limite est $0$.
    On a alors  :
    \[ w_n \geq 0 \iff  v_n - u_n \geq 0 \iff v_n \geq u_n \]
    D'après les monotonies de $(v_n)_{n \in \N}$ et $(u_n)_{n \in \N}$, on a l'encadrement suivant : 
    \[ u_0 \leq u_n \leq v_n \leq v_0 \]
    $(u_n)_{n \in \N}$ est majorée par $v_0$ et est croissante, donc elle converge vers une limite $\ell_1$.
    \\ 
    $(v_n)_{n \in \N}$ est minorée par $u_0$ et est croissante, donc elle converge vers une limite $\ell_2$.
    \\
    D'une part :
    \[ \lim_{n \to +\infty} w_n = 0 \]
    D'autre part :
    \[ \lim_{n \to +\infty} w_n = \lim_{n \to +\infty} (v_n - u_n) = \ell_2 - \ell_1 \]
    Donc :
    \[ \ell_2 - \ell_1 = 0 \iff \ell_2 = \ell_1 \]
\end{proof}

\section{Suites extraites}
\begin{definition}[Extraction]
    Une extraction est une fonction $\varphi : \N \to \N$ qui est strictement croissante.
\end{definition}

\begin{definition}[Suite extraite]
    Une suite extraite ou une sous-suite d'une suite $(u_n)_{n \in \N}$ est une suite de la forme $(u_{\varphi(n)})_{n \in \N}$ où $\varphi$ est une extraction.
\end{definition}

\begin{proposition}
    Soit $(u_n)_{n \in \N}$ une suite et $(u_{\varphi(n)})_{n \in \N}$ une de ses sous-suites.
    \begin{itemize}
        \item Si $(u_n)_{n \in \N}$ est croissante, alors $(u_{\varphi(n)})_{n \in \N}$ aussi.
        \item Si $(u_n)_{n \in \N}$ est décroissante, alors $(u_{\varphi(n)})_{n \in \N}$ aussi.
        \item Si $(u_n)_{n \in \N}$ est majorée, alors $(u_{\varphi(n)})_{n \in \N}$ aussi.
        \item Si $(u_n)_{n \in \N}$ est minorée, alors $(u_{\varphi(n)})_{n \in \N}$ aussi.
        \item Si $(u_n)_{n \in \N}$ est converge, alors $(u_{\varphi(n)})_{n \in \N}$ aussi.
    \end{itemize}
\end{proposition}

\begin{proposition}
    Soit $(u_n)_{n \in \N}$ une suite, alors :
    \[ (u_n)_{n \in \N} \text{ converge} \iff (u_{2n})_{n \in \N} \text{ et } (u_{2n + 1})_{n \in \N} \text{ convergent vers la même limite} \]
\end{proposition}

\begin{theorem}[Théorème de Ramsey]
    Toute suite admet une sous-suite monotone.
\end{theorem}

\begin{proof}
    Soit $(u_n)_{n \in \N}$ une suite. Soit $E = \{ n \in \N,\ \forall m \geq n,\ u_m \leq u_n \}$.
    \\ 
    \textbf{Cas 1 :} $E$ est fini, donc majoré par un entier $N,\ \forall n \geq N : n \notin E$ donc $\exists m > n,\ u_m > u_n$. On définit alors par récurrence une extraction $\varphi : \N \to \N$ en posant $\varphi(0) = N + 1$, puis, étant donnés $\varphi(0) < \varphi(1) < \cdots < \varphi(K)$, on choisit $\varphi(K + 1)$ tel que $u_{\varphi(K+1)} > u_{\varphi(K)}$ et la suite extraite $(u_{\varphi(n)})_{n \in \N}$ est croissante.
    \\ 
    \textbf{Cas 2 :} $E$ est infini. On pose $E = \{ \varphi(n) : n \in \N \}$ avec $\varphi : \N \to \N$.
    \[ \forall k \in \N,\ \varphi(k) \in E, \text{ comme } \varphi(K + 1) > \varphi(K),\ u_{\varphi(K+1)} \leq u_{\varphi(K)} \]
    et la sous-suite $(u_{\varphi(n)})_{n \in \N}$ est décroissante.
\end{proof}

\begin{theorem}[Théorème de Bolzano-Weierstrass]
    Toute suite bornée admet une sous-suite convergente.
\end{theorem}

\begin{proof}
    Soit $(u_n)_{n \in \N}$ une suite bornée. D'après le théorème de Ramsey, il existe une sous-suite monotone $(u_{\varphi(n)})_{n \in \N}$. Comme $(u_{\varphi(n)})_{n \in \N}$ est monotone et bornée, alors elle converge.
\end{proof}

\section{Limites infinies}

\begin{definition}[Limites infinies]
    Soit $(u_n)_{n \in \N}$ une suite.
    \begin{enumerate}
        \item $\lim_{n \to +\infty} u_n = +\infty \iff \forall A \in \R,\ \exists N \in \N,\ \forall n \geq N : u_n \geq A$.
        \item $\lim_{n \to -\infty} u_n = -\infty \iff \forall A \in \R,\ \exists N \in \N,\ \forall n \geq N : u_n \leq A$.
    \end{enumerate}
\end{definition}

\begin{theorem}
    Soit $(u_n)_{n \in \N}$ une suite.
    \begin{enumerate}
        \item Si elle est \textbf{croissante} alors :
        \begin{itemize}
                \item ou bien elle converge.
                \item ou bien elle tend vers $+\infty$.
            \end{itemize}
        \item Si elle est \textbf{décroissante} alors :
        \begin{itemize}
                \item ou bien elle diverge.
                \item ou bien elle tend vers $-\infty$.
            \end{itemize}
    \end{enumerate}
\end{theorem}

\begin{proof}
    Démontrons les propriétés si $(u_n)_{n \in \N}$ est croissante. \\ 
    On distingue deux cas :
    \begin{enumerate}
        \item Si $(u_n)$ est majorée, elle converge, d'après le théorème de convergence monotone car elle est croissante et majorée.
        \item Si $(u_n)$ n'est pas majorée, montrons qu'elle tend vers $+\infty$.
        Soit $A$ un réel. Comme $(u_n)_{n \in \N}$ n'est pas majorée :
        \[ \exists N \in \N,\ u_N \geq A \]
        \[ \forall n \geq N,\ u_n \geq u_N \geq A \]
    \end{enumerate}
    On utilise un raisonnement analogue si $(u_n)_{n \in \N}$ est décroissante.
\end{proof}

\begin{theorem}[Limites par comparaison]
   Soient $(u_n)_{n \in \N}$ et $(v_n)_{n \in \N}$ deux suites telles que $u_n \leq v_n$.
   \begin{enumerate}
   	\item $\lim_{n \to +\infty} u_n = +\infty \implies \lim_{n \to +\infty} v_n = +\infty$.
   	\item $\lim_{n \to +\infty} v_n = -\infty \implies \lim_{n \to +\infty} u_n = -\infty$.
   \end{enumerate}
\end{theorem}

\begin{definition}[Suite de Cauchy]
	Une suite $(u_n)_{n \in \N}$ est de Cauchy si et seulement si :
	\[ \forall \varepsilon > 0,\ \exists N \in \N,\ \forall n_1, n_2 \in \N : \abs{u_{n_1} - u_{n_2}} \leq \varepsilon \]
\end{definition}

\def\arraystretch{1.5}

\begin{table}[!h]
    \centering
    \begin{tabular}{cc}
         \toprule
         Hypothèses & Conclusion \\ 
         \midrule
         \og $+\infty +\infty$ \fg & $+\infty$ \\ 
         \og $-\infty -\infty$ \fg & $-\infty$ \\
         \og $+\infty + \ell$ \fg & $+\infty$ \\
         \og $-\infty + \ell$ \fg & $-\infty$ \\
         \og $-\infty \cdot \ell > 0$ \fg & $-\infty$ \\ 
         \og $-\infty \cdot \ell < 0$ \fg & $+\infty$ \\ 
         \og $+\infty \cdot \ell > 0$ \fg & $+\infty$ \\ 
         \og $+\infty \cdot \ell < 0$ \fg & $-\infty$ \\ 
         \og $+\infty \cdot +\infty$ \fg & $+\infty$ \\ 
         \og $-\infty \cdot -\infty$ \fg & $+\infty$ \\ 
         \og $-\infty \cdot +\infty$ \fg & $-\infty$ \\  
         \og $\frac{0}{\pm \infty}$ \fg & $0$ \\
         \og $\frac{+\infty}{0}$ \fg & $\begin{cases}
         	+ \infty \text{ si } 0^+ \\
         	- \infty \text{ sinon}
         \end{cases}$ \\
         \og $\frac{-\infty}{0}$ \fg & $\begin{cases}
         	- \infty \text{ si } 0^+ \\ 
         	+ \infty \text{ sinon}
         \end{cases}$ \\
         \og $\infty - \infty$ \fg & FI \\
         \og $0 \cdot \infty$ \fg  & FI \\
         \og $\frac{0}{0}$ \fg & FI \\ 
         \og $\frac{\infty}{\infty}$ \fg & FI \\
         \bottomrule
    \end{tabular}
    \caption{Limites infinies ($\ell \in \R$) et formes indéterminées}
    \label{tab:limites_infinies_et_fi}
\end{table}

\chapter{Continuité et limites de fonctions}
\def\arraystretch{1}

\begin{definition}[Limite d'une fonction]
    \begin{enumerate}
        \item En un point $a \in \R$, $\ell \in \R$ :
        \begin{enumerate}
            \item $ \lim_{x \to a} f(x) = \ell \iff 
	\forall \varepsilon > 0,\ \exists \delta > 0,\ \abs{x - a} \leq \delta \implies \abs{f(x) - \ell} \leq \varepsilon $.
            \item $ \lim_{x \to a} f(x) = +\infty \iff 
	\forall A \in \R,\ \exists \delta > 0,\ \abs{x - a} \leq \delta \implies f(x) \geq A $.
            \item $ \lim_{x \to a} f(x) = -\infty \iff 
	\forall A \in \R,\ \exists \delta > 0,\ \abs{x - a} \leq \delta \implies f(x) \leq A $.
        \end{enumerate}
        \item En l'infini, $\ell \in \R$ : 
        \begin{enumerate}
            \item $ \lim_{x \to +\infty} f(x) = \ell \iff 
	\forall \varepsilon > 0,\ \exists A \in \R,\ x \geq A \implies \abs{f(x) - \ell} < \varepsilon $.
            \item $ \lim_{x \to -\infty} f(x) = \ell \iff 
	\forall \varepsilon > 0,\ \exists A \in \R,\ x \leq A \implies \abs{f(x) - \ell} \leq \varepsilon $.
            \item $\lim_{x \to +\infty} f(x) = +\infty \iff 
	\forall A \in \R,\ \exists B \in \R,\ x \geq B \implies f(x) \geq A $.
            \item $ \lim_{x \to +\infty} f(x) = -\infty \iff 
	\forall A \in \R,\ \exists B \in \R,\ x \geq B \implies f(x) \leq A $.
            \item $ \lim_{x \to -\infty} f(x) = +\infty \iff 
	\forall A \in \R,\ \exists B \in \R,\ x \leq B \implies f(x) \geq A $.
            \item $\lim_{x \to -\infty} f(x) = -\infty \iff 
	\forall A \in \R,\ \exists B \in \R,\ x \leq B \implies f(x) \leq A$.
        \end{enumerate}
    \end{enumerate}
\end{definition}

\begin{theorem}
	Soit $I \subseteq \R,\ f : I \longrightarrow \R,\ a \in I \cup \{-\infty, +\infty\},\ \ell \in \R \cup \{-\infty, +\infty\}$
	\[ \lim_{x \to a} f(x) = \ell \iff \forall (u_n)_{n \in \N} : \lim_{n \to +\infty} u_n = a \implies \lim_{n \to +\infty} f(u_n) = \ell \]
\end{theorem}

\begin{definition}[Limite à gauche et à droite]
    Soient $I \subseteq \R,\ a \in I,\ f : I \longrightarrow \R $.
    \begin{enumerate}
        \item $ \lim_{\substack{x \to a \\ x < a}} f(x) = \lim_{x \to a^-} f(x) = \ell \iff \forall \varepsilon > 0,\ \exists \delta > 0,\ a - \delta < x < a \implies \abs{f(x) - \ell} \leq \varepsilon $.
        \item $\lim_{\substack{x \to a \\ x > a}} f(x) = \lim_{x \to a^+} f(x) = \ell \iff \forall \varepsilon > 0,\ \exists \delta > 0,\ a < x < a + \delta \implies \abs{f(x) - \ell} \leq \varepsilon$.
    \end{enumerate}
\end{definition}

\begin{definition}[Continuité]
	Soient $I \subseteq \R,\ a \in I,\ f : I \longrightarrow \R$.
        \\ 
        On dit que $f$ est \textbf{continue} si et seulement si :
        \[\lim_{x \to a} f(x) = f(a).\] 
	On dit que $f$ est continue sur $I$ si elle est continue en tout point de $I$.\\
	On peut également définir la continuité à gauche et à droite.
\end{definition}

\begin{remark}
	Les théorèmes d'opérations avec les limites, de comparaison et des gendarmes sont analogues à ceux vus  dans le chapitre sur les suites réelles.
\end{remark}

\begin{theorem}[Composition de limites]
	Soient $I, J \subseteq \R$ et $f : I \longrightarrow J,\ g : J \longrightarrow \R,\ a \in I$ tels que :
	\begin{enumerate}
		\item $\lim_{x \to a} f(x) = y \in I$.
		\item $\lim_{y \to z} g(z) = \ell$ existe.
	\end{enumerate}
	\[ \lim_{x \to a} g(f(x)) = \ell \]
\end{theorem}



\begin{theorem}[Théorème des valeurs intermédiaires]
	$\forall a, b \in \R,\ a < b,\ f : [a, b] \longrightarrow \R$ une fonction continue.
	\[ \forall y \in [f(a), f(b)],\ \exists c \in [a, b] : f(c) = y \]
\end{theorem}

\begin{proof}
	On utilise la borne supérieure.
	\\
	Soit $E = \{ x \in I \mid f(x) \leq y \}$. $a \in E$ donc $E \neq \varnothing$. On sait que $E \subseteq I$ donc $E$ est majoré.
	\\
	Posons $c = \sup(E)$.
	\\
	Puisque $c = \sup(E)$, il existe une suite $(c_n)_{n \in \N}$ d'éléments de $E$ telle que $\lim_{n \to +\infty} c_n = c$. 
	\\
	Comme $f$ est continue, on a :
	\[ \lim_{n \to +\infty} f(c_n) = f(c) \] 
	Puisque $c_n \in E,\ f(c_n) \leq y$. En passant à la limite, on a :
	\[ f(c) \leq y \]
	Montrons maintenant que $f(c) \geq y$.
	\begin{itemize}
		\item Si $c = b$, on a bien :
		\[ f(c) = f(b) = \geq y \]
		\item Si $c < b$, pour $n$ assez grand :
		\[ c < c + \frac{1}{n} \leq b \]
		Sachant que $c = \sup(E)$, $c + \frac{1}{n} \notin E$, on a donc :
		\[ f \left( c + \frac{1}{n} \right) > y \]
		On a $\lim_{n \to +\infty} c + \frac{1}{n} = c$ et $f$ étant continue :
		\[ \lim_{n \to +\infty} f \left( c + \frac{1}{n} \right) = f(c) \]
		Sachant que $f \left( c + \frac{1}{n} \right) > y$, en passant à la limite :
		\[ f(c) \geq y \]
	\end{itemize}
\end{proof}

\begin{theorem}
	$\forall a, b \in \R,\ a < b,\ f :\ ]a, b[ \ \longrightarrow \R$. 
	Si $f$ est croissante.
	\begin{enumerate}
		\item $f$ admet une limite en $b$, qui est finie si et seulement si $f$ est \textbf{majorée}.
		\item $f$ admet une limite en $a$, qui est finie si et seulement si $f$ est \textbf{minorée}.
	\end{enumerate}
	Si $f$ est décroissante.
	\begin{enumerate}
		\item $f$ admet une limite en $b$, qui est finie si et seulement si $f$ est \textbf{minorée}.
		\item $f$ admet une limite en $a$, qui est finie si et seulement si $f$ est \textbf{majorée}.
	\end{enumerate}
	$\forall x_0 \in \ ]a, b[$, $f$ a une limite à gauche et à droite en $x_0$ et :
	\[ \lim_{x \to x_0^-} f(x) \leq f(x_0) \leq \lim_{x \to x_0^+} f(x) \]
\end{theorem}

\begin{theorem}
	$\forall a, b \in \R,\ a < b,\ f : [a, b] \longrightarrow \R$ une fonction continue.
	\begin{enumerate}
	    \item $ f \text{ strictement croissante} \implies f : [a, b] \longrightarrow [f(a), f(b)] \text{ est une bijection} $.
            \item $ f \text{ strictement décroissante} \implies f : [a, b] \longrightarrow [f(b), f(a)] \text{ est une bijection} $.
	\end{enumerate}
\end{theorem}

\begin{theorem}
	Soit $I \subseteq \R$ et $f : I \longrightarrow \R$ une injection continue.
	\\
	Alors $f$ est strictement monotone, donc bijective. Si on pose $J = f(I)$, $f^{-1} : J \to I$ est continue.
\end{theorem}

\begin{definition}[Segment]
	Un segment est un intervalle fermé borné.
\end{definition}

\begin{theorem}
	Soient $a, b \in \R$ tels que $a < b$ et $f : [a, b] \longrightarrow \R$ une fonction continue. Alors $f$ est bornée sur $[a, b]$ et elle atteint ses bornes.
	\[ \exists m, M \in \R,\ \forall x \in [a, b] : m \leq f(x) \leq M \text{ et } \exists x_0, x_1 \in [a, b] : f(x_0) = m \text{ et } f(x_1) = M \]
\end{theorem}

\begin{definition}[Prolongement par continuité]
	Soient $I \subseteq \R$, $x_0 \in I,\ f : I\backslash\{x_0\} \longrightarrow \R$.
	On suppose que pour $\ell \in \R$, $\lim_{x \to x_0} f(x) = \ell$ existe. Alors la fonction :
	\begin{center}
		$
		\appli{\overset{\sim}{f}}{I}{\R}{x}{
		\begin{cases}
			f(x) \text{ si } x \neq x_0 \\
			\ell \text{ sinon}
		\end{cases}		
		}
		$
	\end{center}
\end{definition}
\chapter{Dérivabilité et accroissements finis}
\def\arraystretch{1}

\section{Dérivabilité, théorèmes de Rolle et des accroissements finis}
\begin{definition}
	Soient $I \subseteq \R,\ f : I \to \R$, $f$ est dérivable en $a \in I$ si et seulement s'il existe $\ell \in \R$ tel que :
	\begin{align*}
		  \lim_{x \to a} \frac{f(x) - f(a)}{x - a} &= \ell
	\end{align*}
    Une autre manière de définir la dérivabilité :
    \begin{align*}
        \lim_{h \to 0} \frac{f(a + h) - f(a)}{h} &= \ell
    \end{align*}
	On note $f'(a) = \ell$ la dérivée de $f$ en $a$.
	Ainsi une fonction dérivable est une fonction dérivable en tout point de $I$.
	On peut également vérifier la limite à gauche et à droite de $a$.
\end{definition}

\begin{proposition}
	Soient $I \subseteq \R,\ f : I \to \R$ et $a \in I$, si $f$ est \textbf{dérivable} en $a$ alors elle est \textbf{continue} en $a$.
\end{proposition}

\begin{proof}
	Quand $f$ est dérivable en $a$, le développement limité suivant existe :
	\[ f(x) \underset{x \to a}{=} f(a) + f'(a)(x - a) + o(x - a) \]
	En passant à la limite :
	\begin{align*}
		\lim_{x \to a} f(x) &= \lim_{x \to a} \left( f(a) + f'(a)(x - a) + o(x - a) \right) \\
		&= \lim_{x \to a} f(a) + \underbrace{\lim_{x \to a} \left( f'(a)(x - a) + o(x - a) \right)}_{= 0} \\
		\lim_{x \to a} f(x) &= f(a)
	\end{align*}
	ce qui veut dire que $f$ est continue en $a$.
\end{proof}

\begin{theorem}
	Soient $I \subseteq \R$, $a \in I$, $\lambda \in \R$ et $f, g$ deux fonctions dérivables en $a$.
	\begin{enumerate}
		\item $f + \lambda g$ est dérivable en $a$ et $(f + \lambda g)'(a) = f'(a) + \lambda g'(a)$.
		\item $f g$ est dérivable en $a$ et $(fg)'(a) = f'(a) g(a) + f(a) g'(a)$.
	\end{enumerate}
\end{theorem}

\begin{proof}
	\leavevmode
	\begin{enumerate}
		\item On calcule $\lim_{x \to a} \frac{f(x) + \lambda g(x) - (f(a) + \lambda g(a))}{x - a}$.
		\begin{align*}
			\lim_{x \to a} \frac{f(x) + \lambda g(x) - (f(a) + \lambda g(a))}{x - a} &= \lim_{x \to a} \frac{f(x) + \lambda g(x) - f(a) - \lambda g(a)}{x - a} \\
			&= \lim_{x \to a} \frac{f(x) - f(a)}{x - a} + \lambda \lim_{x \to a} \frac{g(x) - g(a)}{x - a} \\
			&= f'(a) + \lambda g'(a)
		\end{align*}
		\item On calcule $\lim_{x \to a} \frac{(fg)(x) - (fg)(a)}{x - a}$. On utilise le fait qu'une fonction dérivable en $a$ est continue en $a$ et la définition de la continuité en $a$.
		\begin{align*}
			\lim_{x \to a} \frac{(fg)(x) - (fg)(a)}{x - a} &= \lim_{x \to a} \frac{f(x) g(x) - f(a) g(a)}{x - a} \\
			&= \lim_{x \to a} \left[ \frac{f(x) g(x) - f(a) g(a)}{x - a} + g(x) \frac{f(x) - f(a)}{x - a} \right] \\
			&= \lim_{x \to a} \frac{f(x)g(x) - f(a)g(a)}{x - a} + \lim_{x \to a} g(x) \frac{f(x) - f(a)}{x - a} \\
			&= f(a) \lim_{x \to a} \frac{g(x) - g(a)}{x - a} + g(a) \lim_{x \to a} \frac{f(x) - f(a)}{x - a} \\
			&= f(a) g'(a) + f'(a) g(a)
		\end{align*}
	\end{enumerate}
\end{proof}

\begin{theorem}[\cite{derivation_wikiversite}]
	Soient $I, J \subseteq \R$, $f : I \to \R$ et $g : J \to \R$ telles que $f(I) \subseteq J$ et $a$ un point de $I$. \\
	Si $f$ est dérivable au point $a$ et $g$ est dérivable au point $f(a)$ alors la composée $g \circ f$ est dérivable au point $a$ et :
	\[ (g \circ f)'(a) = g'(f(a)) f'(a) \]
\end{theorem}

\begin{proof}
	\cite{derivation_wikiversite}. Notons $b = f(a)$. Puisque $g$ est dérivable en $b$, il existe une fonction $u : J \to \R$ telle que :
	\[ u(b) = \lim_{y \to b} u(y) = g'(b) \]
	et $\forall y \in J :$ 
	\[ g(y) - g(b) = u(y) (y - b) \]
	En particulier, $f$ est continue au point $a$ car elle y est dérivable: 
	\[ \lim_{x \to a} u(f(x)) = g'(b) \]
	et $\forall x \in I :$
	\[ g(f(x)) - g(f(a)) = u(f(x)) (f(x) - f(a)) \]
	Le taux de variation au point $a$ de la fonction $g \circ f$ s'exprime alors sous la forme :
	\[ \frac{g(f(x)) - g(f(a))}{x - a} = u(f(x)) \cdot \frac{f(x) - f(a)}{x - a} \]
	et quand $x$ tend vers $a$, cette expression tend vers $g'(b) \cdot f'(a) = g'(f(a)) f'(a)$.
\end{proof}

\begin{definition}[Maximum, minimum]
    Soient $I \subseteq \R,\ f : I \to \R$ et $a \in I$.
    \begin{enumerate}
        \item On dit que $a$ est un \textbf{maximum} si et seulement si : $\forall x \in I,\ f(a) \geq f(x)$.
        \item On dit que $a$ est un \textbf{minimum} si et seulement si : $\forall x \in I,\ f(a) \leq f(x)$.
    \end{enumerate}
    On appelle extremum un point qui est soit un maximum soit un minimum.
    \begin{enumerate}
        \item On dit que $a$ est un \textbf{maximum local} si et seulement si : $\exists \varepsilon > 0,\ a \text{ est un maximum de } f_{|]a - \varepsilon, a + \varepsilon[}$.
        \item On dit que $a$ est un \textbf{minimum local} si et seulement si : $\exists \varepsilon > 0,\ a \text{ est un minimum de } f_{|]a - \varepsilon, a + \varepsilon[}$.
    \end{enumerate}
\end{definition}

\begin{theorem}
	Soient $I \subseteq \R,\ f : I \to \R$ une fonction dérivable et $a \in \overset{\circ}{I}$, $\overset{\circ}{I}$ désigne l'intérieur de $I$. 
	\[ a \text{ est un extremum local} \implies f'(a) = 0 \]
\end{theorem}

\begin{proof}
	Soient $a$ un extremum local et $f_{|]a - \varepsilon, a + \varepsilon[}$.
	\\
	\begin{enumerate}
		\item Quand $h > 0$ : 
		\[ \frac{f(a + h) - f(a)}{h} \leq 0,\ \lim_{h \to 0} f'(a) \leq 0 \]
		\item Quand $h < 0$ :
		\[ \frac{f(a + h) - f(a)}{h} \geq 0,\ \lim_{h \to 0} f'(a) \geq 0 \]
	\end{enumerate}
	alors $f'(a) = 0$.
\end{proof}

\begin{theorem}[Théorème de Rolle]
    Soient $a, b \in \R$ tels que $a < b$ et $f : [a, b] \to \R$ une fonction telle que :
    \begin{enumerate}
            \item $f$ est continue sur $[a, b]$.
            \item $f$ est dérivable sur $]a, b[$.
            \item $f(a) = f(b)$.
        \end{enumerate}
    \par \noindent Il existe un $c \in \ ]a, b[$ tel que : \[ f'(c) = 0 \]
\end{theorem}

\begin{proof}\cite{exo7_analyse1}
	\leavevmode
	\begin{enumerate}
		\item Si $f$ est constante, alors n'importe quel $c \in \ ]a, b[$ convient.
		\item Sinon il existe $x_0 \in [a, b]$ tel que $f(x_0) \neq f(a)$. 
		\\
		Supposons par exemple $f(x_0) > f(a)$. Alors $f$ est continue sur $[a, b]$, donc elle admet un maximum en un point $c \in [a, b]$. Mais $f(c) \geq f(x_0) > f(a)$ donc $c \neq a$. De même comme $f(a) = f(b)$ alors $c \neq b$. Ainsi $c \in \ ]a, b[$. En $c$, $f$ est dérivable et admet un maximum local donc $f'(c) = 0$.
	\end{enumerate}
\end{proof}

\begin{theorem}[Théorème des accroissements finis]
    Soient $a, b \in \R$ tels que $a < b$ et $f : [a, b] \to \R$, telle que : 
    \begin{enumerate}
            \item $f$ est continue sur $[a, b]$.
            \item $f$ est dérivable sur $]a, b[$.
        \end{enumerate}
    \par \noindent Il existe un $c \in ]a, b[$ tel que :
	\[ f'(c) = \frac{f(b) - f(a)}{b - a} \]
\end{theorem}

\begin{proof}\cite{exo7_analyse1}
	Posons $\ell = \frac{f(b) - f(a)}{b - a}$ et $g(x) = f(x) - \ell (x - a)$.
	 \\
	 Alors $g(a) = f(a),\ g(b) = f(b) - \frac{f(b) - f(a)}{b - a} \cdot (b - a) = f(a)$.
	 \\
	  Par le théorème de Rolle, il existe un $c \in \ ]a, b[$ tel que $g'(c) = 0$. Or $g'(x) = f'(x) - \ell$. Ce qui donne $f'(c) = \frac{f(b) - f(a)}{b - a}$.
\end{proof}

\begin{corollary}[Inégalité des accroissements finis]
	Soient $a,b \in \R$ tels que $a < b$, $f : \ ]a, b[ \to \R$ une fonction dérivable sur $]a, b[$ et $M$ une constante telle que pour tout $x \in \ ]a, b[,\ \abs{f'(x)} \leq M$.
	\[ \abs{\frac{f(b) - f(a)}{b - a}} \leq M \]
\end{corollary}

\begin{proof}
	D'après le théorème des accroissements finis, il existe un $c \in \ ]a, b[$ tel que $\frac{f(b) - f(a)}{b - a}$. Or pour tout $x \in \ ]a, b[$, $\abs{f'(x)} \leq M$ donc $\abs{f'(c)} \leq M$ et donc :
	\[ \abs{\frac{f(b) - f(a)}{b - a}} \leq M \]
\end{proof}

\begin{proposition}[\cite{exo7_analyse1}]
    Soient $f : [a, b] \to \R$ une fonction continue sur $[a, b]$ et dérivable sur $]a, b[$.
    \begin{enumerate}
            \item $\forall x \in \ ]a, b[,\ f'(x) \geq 0 \iff f \text{ croissante}$.
            \item $\forall x \in \ ]a, b[,\ f'(x) \leq 0 \iff f \text{ décroissante}$.
            \item $\forall x \in \ ]a, b[,\ f'(x) = 0 \iff f \text{ constante}$.
            \item $\forall x \in \ ]a, b[,\ f'(x) > 0 \implies f \text{ strictement croissante}$.
            \item $\forall x \in \ ]a, b[,\ f'(x) < 0 \implies f \text{ strictement décroissante}$ .
        \end{enumerate}
\end{proposition}

\begin{proof}
	Montrons le 1. 
	\\
	\boxed{\implies} : Supposons que pour $x \in \ ]a, b[,\ f'(x) \geq 0$. \\
	Soient $x, y \in \ ]a, b[$ tels que $x \leq y$. D'après le théorème des accroissements finis, il existe un $c \in \ ]x, y[$ tel que $\frac{f(x) - f(y)}{x - y} = f'(c)$.
	\[ f(x) - f(y) = f'(c)(x - y) \]
	Or $f'(x) \geq 0$ pour $x \in \ ]a, b[$ donc $f'(c) \geq 0$ et $x \leq y$ donc $x - y \leq 0$ et $f(x) - f(y) \leq 0$ donc  $f(x) \leq f(y)$.
	\\
	\boxed{\impliedby} : Supposons que pour $x, y \in \ ]a, b[$ tels que $x \leq y$ et $f(x) \leq f(y)$. \\
	On a donc :
	\begin{align*}
		f(y) - f(x) &\geq 0 \\
		\frac{f(y) - f(x)}{y - x} &\geq 0 
	\end{align*}
	On sait que :
	\[ \lim_{y \to x} \frac{f(y) - f(x)}{y - x} = f'(x) \]
	donc :
	\[ f'(x) \geq 0 \]
\end{proof}

\begin{definition}
	Soient $I \subseteq \R,\ f : I \to \R$ et $n \in \N$.
	\begin{enumerate}
		\item $f \in \mathcal{D}^n(I, \R)$ signifie que $f$ est $n$ fois dérivable sur $I$.
		\item $f \in \mathcal{C}^n(I, \R)$ signifie que $f$ est $n$ fois dérivable et que sa dérivée $n$-ième est continue.
		\item $f \in \mathcal{C}^{\infty}(I, \R) = \mathcal{D}^{\infty}(I, \R)$ signifie que $f \in \mathcal{C}^n(I, \R),\ \forall n \in \N$. On dit que les fonctions $\mathcal{C}^{\infty}$ sont des \textbf{fonctions lisses}.
	\end{enumerate}
\end{definition}

\begin{proposition}
	$\forall f, g \in \mathcal{C}^n(I, \R) \implies f + g, f \cdot g, f \circ g \in \mathcal{C}^n(I, \R)$
\end{proposition}

\section{Convexité}
\begin{definition}
    Soient $I \subseteq \R,\ f : I \to \R$.
    \begin{enumerate}
        \item On dit que $f$ est \textbf{convexe} si et seulement si : 
        \[ \forall x, y \in I,\ \lambda \in [0, 1] : f(\lambda x + (1 - \lambda)y) \leq \lambda f(x) + (1 - \lambda) f(y) \]
        \item On dit que $f$ est \textbf{concave} si et seulement si : 
        \[ \forall x, y \in I,\ \lambda \in [0, 1] : f(\lambda x + (1 - \lambda)y) \geq \lambda f(x) + (1 - \lambda) f(y) \]
    \end{enumerate}
\end{definition}

\par Géométriquement, $f$ est convexe signifie que son graphe passe sous les cordes de $f$ et que les tangentes passent sous le graphe. $f$ est concave signifie que son graphe au-dessus des cordes de $f$ et que les tangentes passent par-dessus le graphe.

\begin{theorem}
	Soient $I \subseteq \R,\ f \in \mathcal{D}^2(I, \R)$.
        \begin{enumerate}
                \item $f \text{ convexe} \iff f'' \geq 0$.
                \item $f \text{ concave} \iff f'' \leq 0$.
            \end{enumerate}
\end{theorem}

\begin{theorem}[Inégalité de Jensen]
	Soient $I \subseteq \R$, $f : I \to \R$ une fonction convexe, $x_i \in I$ et $\lambda_i \in [0, 1]$ tels que $\sum_{i=1}^n \lambda_i = 1$.
	\[ f \left( \sum_{i=1}^{n} \lambda_i x_i \right) \leq \sum_{i=1}^{n} \lambda_i f(x_i) \]
	Si $f$ est concave, 
	\[ f \left( \sum_{i=1}^{n} \lambda_i x_i \right) \geq \sum_{i=1}^{n} \lambda_i f(x_i) \]
\end{theorem}

\begin{proof}\cite{inegalite_jensen_bibmath}
	On procède par récurrence pour montrer $P(n) : f \left( \sum_{i=1}^{n} \lambda_i x_i \right) \leq \sum_{i=1}^{n} \lambda_i f(x_i)$. Le principe de la preuve est similaire pour l'autre inégalité.
	\begin{enumerate}
		\item \textbf{Initialisation :} Pour $n = 1$ et $n = 2$ il s'agit de la définition d'une fonction convexe.
		\item \textbf{Hérédité :} Supposons que $P(n)$ vraie pour un $n > 2$. Soient $x_1, \ldots, x_{n+1} \in I$ et $\lambda_1, \ldots, \lambda_{n+1} \in [0,1]$ tels que $\lambda_1 + \cdots + \lambda_n + \lambda_{n+1} = 1$. On veut estimer :
		\[ f(\lambda_1 x_1 + \cdots + \lambda_n x_n + \lambda_{n+1} x_{n+1} ) \]
		On pose :
		\[ 
		\begin{cases}
			\lambda_n' = \lambda_n + \lambda_{n+1} \\
			\lambda_n' x_n' = \lambda_n x_n + \lambda_{n+1} x_{n+1}
		\end{cases}
		\]
		Alors $\lambda_n' \in [0, 1]$. En effet :
		\[ \lambda_n \geq 0,\ \lambda_{n+1} \geq 0 \implies \lambda_n' \geq 0 \]
		et 
		\[ \lambda_n' = 1 - (\lambda_1+ \cdots + \lambda_{n-1}) \leq 1 \]
		On a aussi $x_n' \in I$. En effet, si $x_n \leq x_{n+1}$, alors :
		\[ x_n' = \frac{\lambda_n}{\lambda_n'} x_n + \frac{\lambda_{n+1}}{\lambda_n'} x_{n+1} \leq \frac{\lambda_n}{\lambda_n + \lambda_{n+1}} x_{n+1} + \frac{\lambda_{n+1}}{\lambda_n + \lambda_{n+1}} x_{n+1} \leq x_{n+1} \]
		De même, 
		\[ x_n' \geq x_n \]
		On a 
		\begin{align*}
			f(\lambda_1 x_1 + \cdots + \lambda_n x_n + \lambda_{n+1} x_{n+1}) &= f(\lambda_1 x_1 + \cdots + \lambda_{n-1} x_{n-1} + \lambda_n' x_n') \\
			&\leq \lambda_1 f(x_1) + \cdots  + \lambda_{n-1} f(x_{n-1}) + \lambda_n' f(x_n') 
		\end{align*}
		Puisque que $f$ est convexe,
		\begin{align*}
			f(x_n') &= f \left(\frac{\lambda_n}{\lambda_n'} x_n + \frac{\lambda_{n+1}}{\lambda_n'} x_{n+1}\right) \\
			&\leq \frac{\lambda_n}{\lambda_n'} f(x_n) + \frac{\lambda_{n+1}}{\lambda_n'} f(x_{n+1})
		\end{align*}
		On conclut que :
		\[ f(\lambda_1 x_1 + \cdots + \lambda_n x_n + \lambda_{n+1} x_{n+1}) \leq \lambda_1f(x_1) + \cdots  +\lambda_{n+1} f(x_{n+1}) \]
		Donc $P(n+1)$ est vraie.
	\end{enumerate}
\end{proof}

\begin{proposition}
	Soient $I \subseteq \R,\ f : I \to \R$ et $a \in I$, la tangente de $f$ en $a$ est :
	\[ \mathcal{T}_a(x) = f(a) + f'(a)(x - a) \]
\end{proposition}

\begin{proposition}
	Soient $a, b \in \R$ tels que $a < b,\ f : [a, b] \to \R$.
	La corde $c$ reliant les points $(a, f(a))$ et $(b, f(b))$ est définie par l'équation suivante :
	\[ c = \frac{f(b) - f(a)}{b - a} (x - a) + f(a) \]
\end{proposition}

\begin{proposition}
	Soient $I \subseteq \R,\ f \in \mathcal{D}^2(I, \R),\ a \in I$.
	\begin{align*}
		\begin{cases}
			f'(a) = 0 \\
			f''(a) < 0
		\end{cases}
		\implies a \text{ est un maximum local, si } f''(a) > 0,\ a \text{  est un minimum local}
	\end{align*}
\end{proposition}

\begin{definition}[Suite récurrente]
	Soient $I \subseteq \R,\ f : I \to \R$ et $u_0 \in I$.
	Si $\forall n \in \N,\ u_n \in I$ alors on peut définir :
	\[ u_{n + 1} = f(u_n) \]
\end{definition}

\begin{lemma}
	Soit $I \subseteq \R$. S'il existe un $\ell \in I$ tel que $\lim_{n \to +\infty} u_n = \ell$ et si $f$ est continue en $\ell$ alors 
	\[ f(\ell) = \ell \]
\end{lemma}

\begin{definition}
	Soient $I \subseteq \R,\ f : I \to \R$. On dit que $f$ est stable sur $I$ si et seulement si : 
	\[ f(I) \subseteq I \]
\end{definition}

\begin{proposition}
	Soit $I \subseteq \R$.
	\begin{enumerate}
		\item Si $f$ est \textbf{croissante} sur $I$ alors 
		$
		(u_n)_{n \in \N} = 
		\begin{cases}
			u_0 \in I \\
			u_{n + 1} = f(u_n)
		\end{cases}
		$
		est monotone.
		\[ u_1 \geq u_0 \iff (u_n) \text{ croissante} \]
		\[ u_1 \leq u_0 \iff (u_n) \text{ décroissante} \]
		\item Si $f$ est \textbf{décroissante} sur $I$, alors les suites extraites $(v_n)_{n \in \N} = u_{2n}$ et $(w_n)_{n \in \N} = u_{2n + 1}$ sont monotones, l'une est \textbf{croissante}, l'autre est \textbf{décroissante}.
	\end{enumerate}
\end{proposition}

\begin{definition}[Point fixe]
	Soient $f : \mathcal{D}_f \to \R$ et $x \in \mathcal{D}_f,\ f(x) = x$. On dit que $x$ est un \textbf{point fixe} de $f$.
\end{definition}

\begin{definition}[Coefficient de convergence]
    Soient $(u_n)_{n \in \N}$ tel que $\lim_{n \to +\infty} u_n = \ell \in \R$, $\varepsilon_n = \abs{u_n - \ell}$ et supposons que $\lim_{n \to +\infty} \frac{\varepsilon_{n + 1}}{\varepsilon_n} = K \in \R_+$. 
    On appelle $K$ coefficient de la convergence.
    \begin{itemize}
        \item Si $K = 1$ la convergence est \textbf{lente}.
        \item Si $K = 0$ la convergence est \textbf{rapide}.
        \item Si $0 < K < 1$ la convergence est \textbf{géométrique}.
    \end{itemize}
\end{definition}

\begin{definition}[Fonction contractante]
	Soient $I \subseteq \R,\ f : I \to \R$.
    On dit que $f$ est \textbf{contractante} si et seulement si :
    \[ \exists k \in \ ]0, 1[,\ \forall x, y \in I : \abs{f(x) - f(y)} \leq k \abs{x - y} \]
\end{definition}

\begin{theorem}[Théorème du point fixe]
	Soient $I$ un intervalle fermé et $f : I \to I$ une fonction contractante et continue et $(u_n)_{n \in \N}$ sa suite récurrente associée.
	\begin{enumerate}
        \item Il existe un unique point fixe $\ell \in I$.
        \item $\lim_{n \to +\infty} u_n = \ell$.
        \item La convergence de $(u_n)_{n \in \N}$ est géométrique.
    \end{enumerate}
\end{theorem}

\begin{proof}
	Soient $a, b$ des réels tels que $a < b$.
	\begin{enumerate}
		\item \textbf{Existence} : Il existe un point fixe d'après le théorème des valeurs intermédiaires pour $g(x) = f(x) - x$.
		\begin{enumerate}
			\item $g(a) = f(a) - a \geq 0$.
			\item $g(b) = f(b) - b \leq 0$.
		\end{enumerate}
		\[ \exists c \in [a, b] : g(c) = f(c) - c = 0 \iff f(c) = c \]
		\item \textbf{Unicité} : Soient $\ell_1, \ell_2$ des points fixes.
		Supposons que $\ell_1 \neq \ell_2$ .
		\begin{align*}
			\abs{\ell_1 - \ell_2} &= \abs{f(\ell_1) - f(\ell_2)} \\
								  &\leq k \abs{\ell_1 - \ell_2} \\
								  &< \abs{\ell_1 - \ell_2}
		\end{align*}
		ce qui est contradictoire.
	\end{enumerate}
\end{proof}

\begin{theorem}[Suites récurrentes linéaires d'ordre 2]
	Soit $u_{n+2} = a u_{n+1} + b u_n$ pour $a, b \in \R$.
	On pose l'équation suivante pour $r \in \R$ :
	\[ r^2 - ar - b = 0 \qquad \Delta = a^2 + 4b \]
	Pour $\lambda, \mu, \alpha \in \R$ :
	\begin{itemize}
        \item $\Delta > 0 \implies u_n = \lambda r_1^n + \mu r_2^n$ avec $r_1, r_2$ tels que :
        \begin{align*}
            r_1 &= \frac{a - \sqrt{\Delta}}{2} & r_2 &= \frac{a + \sqrt{\Delta}}{2}
        \end{align*}
        \item $\Delta = 0 \implies u_n = (\lambda + \mu n) r_0^n$ avec $r_0$ tel que :
        \begin{align*}
            r_0 = \frac{a}{2}
        \end{align*}
        \item $\Delta < 0 \implies u_n = \lambda r^n \cos(n \alpha) + \mu r^n \sin(n \alpha)$ avec les racines de la forme $re^{i\alpha}$ et $re^{-i\alpha}$.
    \end{itemize}    
    On trouve $\lambda, \mu$ grâce aux conditions sur les deux premiers termes de la suite.
\end{theorem}

\chapter{Intégration}
\def\arraystretch{1}

\begin{definition}[Fonction en escalier]
	Soient $a, b \in \R$ tels que $a < b$, $n \in \N$ et $f : [a, b] \longrightarrow \R$.
	\\
	$f$ est une fonction en escalier s'il existe une division de $[a, b]$ :
	\[ a = x_0 < \cdots < x_n = b \]
	telle que pour $0 \leq i \leq n,\ x_i \in \R,\  f_{| ]x, x_{i+1}[}$ est constante. On appelle ça une \textbf{subdivision}.
	Autrement dit c'est une fonction constante par morceaux.
\end{definition}

\begin{definition}[Intégrale d'une fonction en escalier]
	Soient $f : [a, b] \longrightarrow \R$ une fonction en escalier et $a_i$ les morceaux de la subdivision.
	\\
	Alors on définit :
	\[ \int_{a}^{b} f(x) \ \diffd x = \sum_{i = 1}^{n} f(a_i) (a_{i+1} - a_i) \]
\end{definition}

\begin{definition}
	Soit $f : [a, b] \longrightarrow \R$. 
	\begin{enumerate}
		\item $f$ est \textbf{Riemann-intégrable} si :
		$\forall \varepsilon > 0,\ \exists e, E$ des fonctions en escalier telles que :
		\begin{enumerate}
			\item $e \leq f \leq E$.
			\item $\int_{a}^{b} (E(x) - e(x)) \ \diffd x < \varepsilon$.
		\end{enumerate}
		\item Soit $f$ \textbf{Riemann-intégrable} sur $[a, b]$.
		\[ \int_{a}^{b} f(x) \ \diffd x = \underset{e \leq f}{\sup} \int_{a}^{b} e(x) \ \diffd x = \underset{E \geq f}{\inf} \int_{a}^{b} E(x) \ \diffd x \]
	\end{enumerate}
\end{definition}

\begin{lemma}
	Soient $f, g : [a, b] \longrightarrow \R$ deux fonctions Riemann-intégrables telles que $f(x) \leq g(x)$, alors :
	\begin{enumerate}
		\item $\int_{a}^{b} f(x) \ \diffd x \leq \int_{a}^{b} g(x) \ \diffd x$.
		\item $\abs{\int_{a}^{b} f(x) \ \diffd x} \leq \int_{a}^{b} \abs{f(x)} \ \diffd x$.
	\end{enumerate}
\end{lemma}

\begin{theorem}[Théorème fondamental de l'analyse]
	Soient $f \in \mathcal{C}^0([a, b])$ et $x \in [a, b]$.
	\begin{align*}
		F(x) = \int_a^x f(t) \ \diffd t
	\end{align*}
	alors $F \in \mathcal{C}^1(]a, b[)$ et pour tout $x \in \ ]a, b[,\ F'(x) = f(x)$.    
\end{theorem}

\begin{corollary}
    Si $F \in \mathcal{C}^1(]a, b[)$ tel que pour tout $x \in \ ]a, b[,\ F'(x) = f(x)$ alors
	\begin{align*}
		\int_a^b f(x) \ \diffd x = [F(x)]_a^b = F(b) - F(a) 
	\end{align*}
\end{corollary}

\begin{proposition}
	Soient $\lambda, a, b \in \R$ et $f, g \in \mathcal{C}^0([a, b])$.
	\begin{align*}
		\int_a^b f(x) + \lambda g(x) \ \diffd x = \int_a^b f(x) \ \diffd x + \lambda \int_a^b g(x) \ \diffd x 
	\end{align*}
\end{proposition}

\begin{proof}
    \begin{align*}
        \int_{a}^{b} f(x) + \lambda g(x) \ \diffd x &= [F(x) + \lambda G(x)]_a^b \\
                                                    &= (F(b) + \lambda G(b)) - (F(a) + \lambda G(a)) \\
                                                    &= F(b) - F(a) + \lambda (G(b) - G(a)) \\
                                                    &= \int_{a}^{b} f(x) \ \diffd x + \lambda \int_{a}^{b} g(x) \ \diffd x
    \end{align*}
\end{proof}

\begin{proposition}
	Soient $a, b \in \R$ et $c \in \ ]a, b[$.
	\begin{align*}
		\int_a^b f(x) \ \diffd x = \int_a^c f(x) \ \diffd x + \int_c^b f(x) \ \diffd x
	\end{align*}
\end{proposition}

\begin{theorem}[Théorème de la moyenne]
    Soit $f \in \mathcal{C}^0([a, b])$. $\exists c \in \ ]a, b[$  tel que :
    \[ \frac{\int_{a}^{b} f(x) \ \diffd x}{b - a} = f(c) \]
\end{theorem}

\begin{theorem}[Intégration par parties]
	Soient $u, v \in \mathcal{C}^1([a, b])$ alors
	\begin{align*}
		\int_a^b u'(x) v(x) \ \diffd x = [u(x)v(x)]_a^b - \int_a^b u(x) v'(x) \ \diffd x
	\end{align*}
  \end{theorem}

\begin{proof}
    \begin{align*}
        (uv)'(x) &= u'(x) v(x) + u(x)v'(x) \\
        \iff u'(x)v(x) &= (uv)'(x) - u(x)v'(x)
    \end{align*}
    Par croissance et linéarité de l'intégrale :
    \begin{align*}
        \int_a^b u'(x)v(x) \ \diffd x &= \int_a^b (uv)'(x) \ \diffd x - \int_a^b u(x)v'(x) \ \diffd x \\
        \int_a^b u'(x)v(x) \ \diffd x &=  [u(x)v(x)]_a^b - \int_a^b u(x)v'(x) \ \diffd x 
    \end{align*}
\end{proof}

\begin{example}
	Soit $I = \int_{0}^{1} x e^x \ \diffd x$.
	\\
	Procédons par intégration par parties, posons :
	\[
	\begin{cases}
		u'(x) = e^x,\ v(x) = x \\
		u(x) = e^x,\ v'(x) = 1
	\end{cases}
	\]
	On a ensuite :
	\begin{align*}
		I &= [xe^x]_0^1 - \int_{0}^{1} e^x \ \diffd x \\
		  &= e - [e^x]_0^1 \\
		  &= e - (e - 1) \\
		  &= 1
	\end{align*}
\end{example}

\begin{theorem}[Intégration par changement de variable]
	Soient $I \subseteq \R,\ f \in \mathcal{C}^0(I)$ et $\varphi \in \mathcal{C}^1([a, b])$ tel que $\varphi ([a, b]) \subseteq I$ alors 
	\begin{align*}
		\int_a^b f(\varphi(x)) \cdot \varphi'(x) \ \diffd x = \int_{\varphi(a)}^{\varphi(b)} f(x) \ \diffd x
	\end{align*}
\end{theorem}

\begin{proof}
    \begin{align*}
        (f \circ \varphi)'(x) &= f' \circ \varphi(x) \cdot \varphi'(x) \\
        \int_a^b (f \circ \varphi)'(x) \ \diffd x &= \int_a^b f' \circ \varphi(x)  \cdot \varphi'(x) \ \diffd x \\
         &= [F(\varphi(x))]_a^b \\
         &= F(\varphi(b)) - F(\varphi(a)) \\
         &= [F(x)]_{\varphi(a)}^{\varphi(b)} \\
         &= \int_{\varphi(a)}^{\varphi(b)} f(x) \ \diffd x 
    \end{align*}
\end{proof}

\begin{example}
	Soit $I = \int_{1}^{e} \frac{(\ln(x))^n}{x} \ \diffd x$ pour $n \in \N$.
	\\
	Procédons par changement de variable, posons :
	\[
	\begin{cases}
		u = \ln(x) \\
		\diffd u = \frac{1}{x} \ \diffd x
	\end{cases}
	\]
	\begin{align*}
		I &= \int_{\ln(1)}^{\ln(e)} u^n \ \diffd u \\
		  &= \int_{0}^{1} u^n \ \diffd u \\
		  &= \left[ \frac{u^{n+1}}{n+1} \right]_0^1 \\
		  &= \frac{1}{n + 1}
	\end{align*}
\end{example}

Lorsque nous sommes confrontés à une intégrale de fonctions trigonométriques, on peut se ramener à une intégrale de fraction rationnelle en posant le changement de variable suivant :
\[ u = \tan(\frac{x}{2}) \]
\begin{multicols}{3}
    \begin{enumerate}
        \item $\cos(x) = \frac{1 - u^2}{1 + u^2}$
        \item $\sin(x) = \frac{2u}{1 + u^2}$
        \item $\diffd x = \frac{2}{1 + u^2} \diffd u$
    \end{enumerate}
\end{multicols}
\begin{proof}
	\leavevmode
    \begin{enumerate}
        \item \[ \cos(x) = \cos^2\left( \frac{x}{2} \right) - \sin^2\left( \frac{x}{2} \right) \]
            \begin{align*}
                \cos^2\left( \frac{x}{2} \right) - \sin^2\left( \frac{x}{2} \right) &= \cos^2\left( \frac{x}{2} \right) - \sin^2\left( \frac{x}{2} \right) \cdot \frac{1 + \tan^2 \left( \frac{x}{2} \right)}{1 + \tan^2 \left( \frac{x}{2} \right)}
            \end{align*}
            Posons $A(x) = \left( \cos^2\left( \frac{x}{2} \right) - \sin^2\left( \frac{x}{2} \right) \right) \left( 1 + \tan^2\left( \frac{x}{2} \right) \right)$.
            \begin{align*}
                A(x) &= \cos^2\left( \frac{x}{2} \right) - \sin^2\left( \frac{x}{2} \right) + \left[ \cos^2\left( \frac{x}{2} \right) - \sin^2\left( \frac{x}{2} \right) \right] \frac{\sin^2 \left( \frac{x}{2} \right)}{\cos^2 \left( \frac{x}{2} \right)} \\
                &= \cos^2 \left( \frac{x}{2} \right) - \frac{\sin^4 \left( \frac{x}{2} \right)}{\cos^2 \left( \frac{x}{2} \right)} \\
                &= \frac{\cos^4 \left( \frac{x}{2} \right) - \sin^4 \left( \frac{x}{2} \right)}{\cos^2 \left( \frac{x}{2} \right)}\\
                &= \frac{\left[ \cos^2\left( \frac{x}{2} \right) - \sin^2 \left( \frac{x}{2} \right) \right] \left[ \cos^2\left( \frac{x}{2} \right) + \sin^2 \left( \frac{x}{2} \right) \right]}{\cos^2 \left( \frac{x}{2} \right)}
            \end{align*}
             On sait que $\forall x \in \R, \cos^2(x) + \sin^2(x) = 1$.
             Ainsi : 
             \begin{align*}
                 A(x) = 1 - \frac{\sin^2\left( \frac{x}{2} \right)}{\cos^2\left( \frac{x}{2} \right)} = 1 - \tan^2 \left( \frac{x}{2} \right)
             \end{align*}
             Ainsi en posant $u = \tan(\frac{x}{2})$, on retrouve bien :
             \begin{align*}
                 \cos^2\left( \frac{x}{2} \right) - \sin^2\left( \frac{x}{2} \right) = \cos(x) = \frac{1 - u^2}{1 + u^2}
             \end{align*}
     \item \[ \sin(x) =  2 \sin(\frac{x}{2}) \cos(\frac{x}{2}) \]
        \begin{align*}
            2\sin(\frac{x}{2})\cos(\frac{x}{2}) &= 2\sin(\frac{x}{2})\cos(\frac{x}{2}) \cdot \frac{1 + \tan^2 \left( \frac{x}{2} \right)}{1 + \tan^2 \left( \frac{x}{2} \right)} \\
            &= \frac{2\sin(\frac{x}{2})\cos(\frac{x}{2}) \cdot \left( 1 + \tan^2 \left( \frac{x}{2} \right) \right)}{1 + \tan^2 \left( \frac{x}{2} \right)}
        \end{align*}
        Posons $A(x) = 2\sin(\frac{x}{2})\cos(\frac{x}{2}) \cdot \left( 1 + \tan^2 \left( \frac{x}{2} \right) \right)$.
        \begin{align*}
            A(x) &= 2\sin(\frac{x}{2})\cos(\frac{x}{2}) + 2\sin(\frac{x}{2})\cos(\frac{x}{2}) \tan^2 \left( \frac{x}{2} \right) \\
            &= 2\sin(\frac{x}{2})\cos(\frac{x}{2}) + 2\sin(\frac{x}{2})\cos(\frac{x}{2}) \frac{\sin^2 \left( \frac{x}{2} \right)}{\cos^2 \left( \frac{x}{2} \right)} \\
            &= 2\sin(\frac{x}{2})\cos(\frac{x}{2}) + \frac{2 \sin^3 \left( \frac{x}{2} \right)}{\cos(\frac{x}{2})} \\ 
            &= 2 \left( \sin(\frac{x}{2}) \left[ \cos(\frac{x}{2}) + \frac{\sin^2 \left( \frac{x}{2} \right)}{\cos(\frac{x}{2})} \right] \right) \\
            &= 2 \left( \sin(\frac{x}{2}) \left[ \frac{\cos^2 \left(\frac{x}{2}\right) + \sin^2 \left( \frac{x}{2} \right)}{\cos(\frac{x}{2})} \right] \right)
        \end{align*}
        On sait que $\forall x \in \R, \cos^2(x) + \sin^2(x) = 1$. Ainsi :
        \begin{align*}
            A(x) &= 2 \frac{\sin(\frac{x}{2})}{\cos(\frac{x}{2})} \\
            &= 2 \tan(\frac{x}{2})
        \end{align*}
        On a donc :
        \begin{align*}
            2\sin(\frac{x}{2})\cos(\frac{x}{2}) = \sin(x) = \frac{2 \tan(\frac{x}{2})}{1 + \tan^2 \left( \frac{x}{2} \right)}
        \end{align*}
        En posant $u = \tan(\frac{x}{2})$ on retrouve :
        \[ \sin(x) = \frac{2u}{1 + u^2} \]
        \item 
        \begin{align*}
            u = \tan(\frac{x}{2}) &\iff \arctan(u) = \frac{x}{2} \\
            &\iff x = 2 \arctan(u) \\
            &\iff \diffd x = \frac{2}{1 + u^2} \diffd u
        \end{align*}
    \end{enumerate}
\end{proof}




\chapter{Equations différentielles linéaires}
\def\arraystretch{1}

\noindent Pour résoudre une équation différentielle, nous allons suivre ces 3 étapes :
\begin{enumerate}
	\item Déterminer les solutions de l'équation homogène associée.
	\item Déterminer une solution particulière de l'équation différentielle.
	\item Combiner les solutions précédentes pour obtenir la solution générale.
\end{enumerate}

\section{\'Equations différentielles d'ordre 1}

\begin{definition}[\'Equation différentielle d'ordre 1]
	Soient $y \in \mathcal{C}^1(\R)$ et $a, b \in \mathcal{C}^0(\R)$. \\
	Une équation différentielle d'ordre 1 est une équation de la forme
	\[ y' + a(x)y = b(x) \]
\end{definition}

\begin{theorem}
	Soient $y \in \mathcal{C}^1(\R)$ et $a \in \mathcal{C}^0(\R)$, les solutions de :
	\[ y' + a(x)y = 0 \]
	sont de la forme $C \in \R,\ A'(x) = a(x)$:
	\[ y = C \cdot \exp(-A(x)) \]
\end{theorem}

\begin{proof}
    Il faut montrer l'inclusion dans les deux sens. \\
    \boxed{\supseteq} :
	Définissons
	\[ y = C e^{-A(x)} \]
	on aurait donc
	\[ y' = -C a(x)e^{-A(x)} \]
	puis
	\[ y' + a(x)y(x) = -C a(x) e^{-A(x)} + a(x)Ce^{-A(x)} = 0 \]
    \boxed{\subseteq} : \cite{bibmath_resolution_eq_diff} :
    Supposons $y$ solution de $y' + a(x) y = 0$.
    Alors il existerait un $C \in \R$ tel que $y(x) = Ce^{-A(x)}$.
    Posons $f(x) = y(x)e^{A(x)}$.
    \begin{align*}
        f'(x) &= y'(x)e^{A(x)} + y(x)e^{A(x)}A'(x) \\
        &= -a(x)y(x)e^{A(x)} + a(x)y(x)e^{A(x)} \\
        &= 0
    \end{align*}
    Cela implique que $f(x) = C,\ C \in \R$. Donc :
    \begin{align*}
        C = y(x)e^{A(x)} \iff y(x) = Ce^{-A(x)}
    \end{align*}
\end{proof}

\begin{example}
    $(E_1) : y' - 2xy = 0$.
    En appliquant le théorème précédent on obtient les solutions :
    \begin{align*}
        y_0 = Ce^{x^2}
    \end{align*}
\end{example}

Pour trouver une solution particulière d'une équation différentielle d'ordre 1. Nous utilisons $y_{h}(x)$, sauf qu'ici, $C$ n'est plus une constante mais une fonction. Cette méthode est appelée \textbf{variation de la constante}.
    \begin{align*}
        \begin{cases}
            y_{p} = C(x)e^{-A(x)} \\
            y_{p}' = C'(x)e^{-A(x)} - C(x)a(x)e^{-A(x)}
        \end{cases}
    \end{align*}
    On obtient alors en remplaçant dans l'équation générale :
    \begin{align*}
        &y_{p}' + a(x) y_{p} = b(x) \\
        \iff &C'(x)e^{-A(x)} - C(x)a(x) e^{-A(x)} + a(x)C(x)e^{-A(x)} = b(x) \\
        \iff &C'(x) e^{-A(x)} = b(x) \\
        \iff &C(x) = \int b(x) e^{A(x)}
    \end{align*}

\begin{example}
    $(E_1) : y' - 2xy = \exp(x^2 - x)$.
    On utilise donc la solution homogène trouvée précédemment avec la variation de la constante.
    \begin{align*}
        \begin{cases}
            y_p = C(x) e^{x^2} \\ 
            y_p' = C'(x) e^{x^2} + C(x)2x e^{x^2}
        \end{cases}
    \end{align*}
    On remplace dans l'équation $(E_1)$.
    \begin{align*}
        y_p' - 2x y_p &= \exp(x^2 - x) \\
        C'(x) e^{x^2} + C(x)2x e^{x^2} - 2x C(x)e^{x^2} &= \exp(x^2 - x) \\
        C'(x)e^{x^2} &= \exp(x^2 - x) \\ 
        C'(x) &= \frac{\exp(x^2 - x)}{e^{x^2}} \\
        C'(x) &= \exp(x^2 - x - x^2) \\
        C'(x) &= \exp(-x) \\
        C(x) &= -\exp(-x)
    \end{align*}
    Ainsi une solution particulière de $(E_1)$ :
    \[ y_p = -\exp(-x) \exp(x^2) = -\exp(x^2 - x) \]
\end{example}

\section{\'Equations différentielles d'ordre 2}

\begin{definition}[\'Equations différentielle d'ordre 2]
	Soient $y \in \mathcal{C}^2(\R)$, $b \in \mathcal{C}^0(\R)$ et $p, q \in \R$. \\
	Une équation différentielle d'ordre 2 est une équation de la forme
	\[ y'' + py' + q = b(x) \]
\end{definition}

\begin{theorem}
  Soient $y \in \mathcal{C}^2(\R),\ p, q \in \R$ et $(E)$ l'équation suivante :
  \[ (E) : y'' + py' + qy = 0, \]
  On s'intéresse d'abord à cette équation associée :
  \[ \ \lambda^{2} + p\lambda + q = 0 \]
  Pour $\Delta = p^2 - 4q$. 
  \begin{enumerate}
	\item Cas 1 :
		  $\Delta > 0$
		  \[ y = C_{1} e^{\lambda_{1} x} + C_{2} e^{\lambda_{2} x},\ \lambda_{i} C_{i} \in \R \]
        \[ \lambda_{1/2} = \frac{-p \pm \sqrt{\Delta}}{2} \]
	\item Cas 2:
		  $\Delta = 0$
		  \[ y = (C_{1} + C_{2}x) e^{\lambda x},\ C_{i}, \lambda \in \R \]
        \[ \lambda = \frac{-p}{2} \]
	\item Cas 3 :
		  $\Delta < 0$
		  $,\ \lambda_{1} = a + ib,\ a, b \in \R,\ \lambda_{2} = a - ib = \overline{\lambda_{1}}$
		  \[ y = e^{ax}(C_1 \cos(\abs{b}x) + C_2 \sin(\abs{b}x)), C_i \in \R \]
        \[ \lambda_{1/2} = \frac{-p \pm i \sqrt{-\Delta}}{2} \]
  \end{enumerate}
\end{theorem}

\par La variation de la constante est difficilement applicable sur les équations différentielles de second ordre, nous devons trouver d'autres méthodes.
Toutes les équations de second ordre qu'on étudiera dans ce chapitre auront pour second membre une composé de fonctions polynomiales, exponentielles et trigonométriques. Ainsi, nous pouvons utiliser ces propriétés.
\\
\par \noindent Voici une méthode pour trouver une solution particulière d'une équation de type :
    \[ y'' + py' + qy = b(x) \]
    pour $p, q \in \R,\ b \in \mathcal{C}^0(\R),\ y \in \mathcal{C}^2(\R)$. \\
    Soient $\alpha, \beta \in \R$, $m \in \N$, $P, P_1, P_2, Q_1, Q_2 \in \R[X]$. 
    \begin{itemize}
        \item \cite{exo7_analyse1} Si $b(x) = e^{\alpha x}(P_1(x) \cos(\beta x) + P_2(x) \sin(\beta x) )$
        \[y_{p} = x^{m} e^{\alpha x}(Q_1(x) \cos(\beta x) + Q_2(x) \sin(\beta x)) ,\ \deg(Q_1), \deg(Q_2) \leq \max\left\{ \deg(P_1), \deg(P_2) \right\} \]
        \[ m = \begin{cases}
            0 \text{ si } \alpha + i\beta \text{ n'est pas racine de l'équation caractéristique} \\
            1 \text{ sinon}
        \end{cases} \]
        \item Si $b(x) = P(x) e^{\alpha x}$ 
        \[ y_p = x^m Q(x)e^{\alpha x},\ \deg(Q) \leq \deg(P) \]
        avec $m$ l'ordre de multiplicité (voir \autoref{def:ordre_mult}) de la racine $\alpha$ par rapport à l'équation caractéristique associée.
\end{itemize}

\begin{remark}
    Il existe des propriétés analogues pour les équations différentielles du premier ordre, parfois cela est plus rapide qu'avec la variation de la constante.
\end{remark}

\begin{example}
    $(E) : y'' - 2y' + 3y = 9x^2 e^{2x} + 4e^x$.
    Tout d'abord, résolvons l'équation homogène associée :
    \[ (E_0) : y'' - 2y' + 3y = 0 \]
    \[ \lambda^2 - 2\lambda + 3 = 0 \]
    \[ \Delta = (-2)^2 - 4 \cdot 3 = -8 < 0 \]
    \begin{align*}
        \lambda_1 &= \frac{2 - i\sqrt{8}}{2} & \lambda_2 &= \frac{2 + i\sqrt{8}}{2} \\
        &= 1 - i\sqrt{2} & &= 1 + i\sqrt{2}
    \end{align*}
    Ainsi les solutions de $(E_0)$ sont :
    \[ y_0 = e^x \left(C_1 \cos(\sqrt{2}x) + C_2 i \sin(\sqrt{2}x) \right) \]
    Trouvons une solution particulière de l'équation :
    \[ (E_1) : y'' - 2y' + 3y = 9x^2 e^{2x} \]
    On remarque que le second membre (le \og $b(x)$ \fg) est de la forme $P(x) e^{\alpha x}$ avec $\deg(P) = 2$ et $\alpha \notin \{1 \pm i\sqrt{2}\}$, ainsi la solution particulière est de la forme $y_1 = (ax^2 + bx + c)e^{2x},\ a, b, c \in \R$.
    On a donc :
    \begin{align*}
        \begin{cases}
            y_1 = (ax^2 + bx + c)e^{2x} \\
            y_1' = (2ax + b)e^{2x} + 2(ax^2 + bx + c)e^{2x} = (2ax + b)e^{2x} + 2y_1 \\
            y_1'' = 2a e^{2x} + 2(2ax + b)e^{2x} + 2y_1'
        \end{cases}
    \end{align*}
    En remplaçant dans $(E_1)$ :
    \begin{align*}
        2ae^{2x} + 2(2ax + b)e^{2x} + 2y_1' - 2[(2ax + b)e^{2x} + 2y_1] + 3y_1 &= 9x^2 e^{2x} \\
        2ae^{2x} + (4ax + 2b)e^{2x} + 2y_1' - (4ax + 2b)e^{2x} - 4y_1 + 3y_1 &= 9x^2 e^{2x} \\
        2ae^{2x}+ 2y_1' - y_1 &= 9x^2e^{2x} \\
        2ae^{2x} + 2[(2ax + b)e^{2x} + 2y_1] - y_1 &= 9x^2e^{2x} \\
        2ae^{2x} + (4ax + 2b)e^{2x} + 4 y_1 - y_1 &= 9x^2 e^{2x} \\
        2ae^{2x} + (4ax + 2b)e^{2x} + 3y_1 &= 9x^2e^{2x} \\
        2ae^{2x} + (4ax + 2b)e^{2x} + 3[(ax^2 + bx + c)e^{2x}] &= 9x^2 e^{2x} \\ 
        2ae^{2x} + (4ax + 2b)e^{2x} + (3ax^2 + 3bx + 3c)e^{2x} &= 9x^2 e^{2x} \\
        (2a + 4ax + 2b + 3ax^2 + 3bx + 3c)e^{2x} &= 9x^2 e^{2x} \\
        [3ax^2 + (4a + 3b)x + (2a + 2b + 3c)]e^{2x} &= 9x^2 e^{2x}
    \end{align*}
    On procède par identification :
    \begin{align*}
        \systeme{
            3a = 9,
            4a + 3b = 0,
            2a + 2b + 3c = 0
        }
        \iff 
        \systeme{
            a = 3,
            4a + 3b = 0,
            2a + 2b + 3c = 0
        }
        \iff 
        \systeme{
            a = 3,
            3b = -12,
            2a + 2b + 3c = 0
        }
        \\ 
        \iff 
        \systeme{
            a = 3,
            b = -4,
            2a + 2b + 3c = 0
        }
        \iff 
        \systeme{
            a = 3,
            b = -4,
            3c = 2
        }
        \iff 
        \systeme{
            a = 3,
            b = -4,
            c = \frac{2}{3}
        }
    \end{align*}
    \[ y_1 = \left(3x^2 - 4x + \frac{2}{3} \right)e^{2x} \]
    Maintenant trouvons une solution particulière de l'équation :
    \[ (E_2) : y'' - 2y' + 3y = 4e^x \]
    Ici on a encore une forme $P(x)e^{\alpha x}$ avec $P(x) = 1, \deg(P) = 0$ et $\alpha \notin \{ 1 \pm i \sqrt{2} \}$ ainsi $y_2 = k e^{x}, k \in \R$
    \begin{align*}
        \begin{cases}
            y_2 = k e^{x} \\
            y_2' = k e^{x} \\
            y_2'' = k e^{x}
        \end{cases}
    \end{align*}
    \begin{align*}
        k e^{x} -2k e^{x} + 3k e^{x} &= 4e^x \\ 
        2k e^x &= 4e^x \\
        k e^x = 2e^x
    \end{align*}
    \[ y_2 = 2e^x \]
    Solution générale :
    \begin{align*}
        y &= y_0 + y_1 + y_2 \\
        y &= e^x \left(C_1 \cos(\sqrt{2}x) + C_2 i \sin(\sqrt{2}x) \right) + \left(3x^2 - 4x + \frac{2}{3} \right)e^{2x} + 2e^x \\
        y &= \left(C_1 \cos(\sqrt{2}x) + C_2 i \sin(\sqrt{2}x) + 2 \right)e^x + \left(3x^2 - 4x + \frac{2}{3} \right)e^{2x}
    \end{align*} 
\end{example}
\chapter{Développements limités et formules de Taylor}

Dans cette partie, en l'absence de précisions supplémentaires, $I$ désigne un intervalle de $\R$.

\begin{notation}
    Dans cette section, on définit la notation suivante : 
    \[ \overline{I} \coloneqq I \cup \{\pm \infty\} \]
\end{notation}

\par Il arrive parfois des situations où l'on se retrouve avec des formes indéterminées de type \og $\frac{0}{0}$ \fg ou \og $\frac{\infty}{\infty}$ \fg lorsque nous essayons de calculer les limites, le théorème suivant permet de lever l'indétermination assez simplement.
Plus tard dans ce chapitre, nous pourrons également utiliser les développements limités pour lever les indéterminations.
\section{Règle de l'Hôpital}
\begin{theorem}[Règle de l'Hôpital]\cite{regle_hopital_bibmath}
    Soient $a, b, \ell \in \overline{\R}$ tels que $a < b$ et $f, g : ]a, b[ \to \R$ deux fonctions dérivables telles que $g'$ ne s'annule pas. 
    \begin{enumerate}
        \item Si $\lim_{x \to a} f(x) = \lim_{x \to a} g(x) = 0$ et si $\lim_{x \to a} \frac{f'(x)}{g'(x)} = \ell$, alors $\lim_{x \to a} \frac{f(x)}{g(x)} = \ell$.
        \item Si $\lim_{x \to a} g(x) = \pm \infty$ et si $\lim_{x \to a} \frac{f'(x)}{g'(x)} = \ell$, alors $\lim_{x \to a} \frac{f(x)}{g(x)} = \ell$.
    \end{enumerate}
\end{theorem}


\begin{example}
    $\lim_{x \to 0} \frac{\sin(2x)}{x}$, on a ici une forme indéterminée \og $\frac{0}{0}$ \fg.
    On remarque que $\sin(2x)$ et $x$ sont dérivables en 0 et $x' = 1 \neq 0$, on utilise donc la règle de l'Hôpital pour lever l'indétermination.
    \begin{align*}
        \lim_{x \to 0} \frac{\sin(2x)}{x} &= \lim_{x \to 0} \frac{\sin'(2x)}{x'} \\
        &= \lim_{x \to 0} \frac{2 \cos(x)}{1} \\
        &= 2 \lim_{x \to 0} \cos(x) \\
        &= 2
    \end{align*}
\end{example}

\begin{definition} 
    Soit $a \in \R$.
    On dit que $x$ est au voisinage d'un point $a$ si et seulement si :
    \[ \exists \varepsilon > 0, \text{ tel que } x \in ]a - \varepsilon, a + \varepsilon[ \]
\end{definition}

\section{Relations de négligeabilité, domination, d'équivalence}
\begin{definition}
    Soient $I \subset \R$, $f, g : I \to \R$, $a \in \overline{I}$ et $\varepsilon$ telle que $\lim_{x \to a} \varepsilon(x) = 0$.
    \begin{enumerate}
        \item On dit que $f$ est \emph{dominée} par $g$ au voisinage de $a$ s'il existe un $B \in \R_+$ tel que $\abs{f(x)} \leq B \abs{g(x)}$ au voisinage de $a$. 
        On écrit alors $f \underset{a}{=} O(g)$ ou $f = O_a(g)$.
        \item On dit que $f$ est \emph{négligeable} devant $g$ au voisinage de $a$ si $f(x) = g(x) \cdot \varepsilon(x)$.
        On écrit alors $f \underset{a}{=} o(g)$ ou $f = o_a(g)$.
        \item On dit que $f$ est \emph{équivalente} à $g$ au voisinage de $a$ si $f(x) = g(x) (1 + \varepsilon(x))$.
        On écrit alors $f \underset{a}{\sim} g$.
    \end{enumerate}
\end{definition}

\begin{proposition}\leavevmode
    \begin{enumerate}
        \item Si $\lim_{x \to a} \frac{f(x)}{g(x)} = 0$ alors $f$ est négligeable devant $g$.
        \item Si $\lim_{x \to a} \frac{f(x)}{g(x)} = 1$ alors $f$ est équivalente à $g$.
        \item Si $\lim_{x \to a} \frac{f(x)}{g(x)}$ est bornée, alors $f$ est dominée par $g$.
    \end{enumerate}   
\end{proposition}

\begin{proposition}
\begin{multicols}{2}
    \begin{enumerate}
        \item $o(1) + o(1) = o(1)$
        \item $\forall \lambda \in \R,\ \lambda \cdot o(1) = o(1)$
        \item $\forall n \in \N^*,\ (o(1))^n = o(1)$
        \item $\forall \alpha > 0,\ (o(1))^{\alpha} = o(1)$
        \item $\forall \alpha \in \R,\ (1 + o(1))^{\alpha} = 1 + o(1)$
        \item $O(1) + O(1) = O(1)$
        \item $\forall \lambda \in \R,\ \lambda \cdot O(1) = O(1)$
        \item $\forall n \in \N^*,\ (O(1))^n = O(1)$
        \item $o(1) \cdot O(1) = o(1)$
    \end{enumerate}
\end{multicols}
\end{proposition}

\begin{proposition}
\begin{multicols}{2}
    \begin{enumerate}
        \item $\ln(x) \underset{+\infty}{\sim} o(x)$
        \item $\forall \alpha, \beta \in \R_+^*,\ (\ln(x))^{\beta} = o(x^{\alpha})$
        \item $\forall \alpha, \beta \in \R_+^*,\ x^{\beta} = o(e^{\alpha x})$
    \end{enumerate}
\end{multicols}
\end{proposition}

\begin{proposition}
    Soit $f$ une fonction polynomiale. 
    \begin{enumerate}
        \item Un équivalent de $f$ en l'infini est un son monôme de plus haut degré.
        \item Un équivalent de $f$ en 0 est son monôme de plus bas degré.
    \end{enumerate}
\end{proposition}

\section{Développements limités}

\begin{definition}[Polynôme de Taylor]
    Soit $f \in \mathcal{D}^n(I, \R)$ alors son polynôme de Taylor en $x_0$ est :
    \[ T_{n, x_0}^f(x) = \sum_{k = 0}^{n} \frac{f^{(k)}(x_0)}{k!} (x - x_0)^k \]
\end{definition}

\begin{theorem}[Formule de Taylor-Young]
    Soient $f \in \mathcal{C}^{n}(I, \R),\ x_0 \in I$.
    \begin{align*}
        f(x) = \sum_{k = 0}^{n} \frac{f^{(k)}(x_0)}{k!} (x - x_0)^k + o((x - x_0)^n) 
    \end{align*}
\end{theorem}

\begin{theorem}[Formue de Taylor-Lagrange]
    Soient $f \in \mathcal{C}^{n+1}(I, \R),\ x_0 \in I$. 
    \begin{align*}
        \exists c \in 
        \begin{cases}
            ]x_0, x[ \text{ si } x > x_0 \\
            ]x, x_0[ \text{ si } x < x_0
        \end{cases},\
        f(x) = \sum_{k = 0}^{n} \frac{f^{(k)}(x_0)}{k!} (x - x_0)^k + 
        \frac{f^{(n+1)}(c)}{(n+1)!} (x - x_0)^{n+1}
    \end{align*}
\end{theorem}

\begin{theorem}[Formule de Taylor avec reste intégral]
    Soient $f \in \mathcal{C}^{n+1}(I, \R),\ x_0 \in I$.
    \begin{align*}
        f(x) = \sum_{k = 0}^{n} \frac{f^{(k)}(x_0)}{k!} (x - x_0)^k + \int_{x_0}^{x} \frac{f^{(n+1)}(t)}{n!} (x - t)^n \diffd t
    \end{align*}
\end{theorem}

\begin{corollary}[Inégalité de Taylor-Lagrange]
    Si $\forall x \in I,\ \abs{f^{(n+1)}(x)} \leq M$ alors
    \[ \abs{\int_{x_0}^{x} \frac{f^{(n+1)}(t)}{n!} (x - t)^n \ \diffd t} \leq M \frac{\abs{x - x_0}^{n+1}}{(n+1)!} \]
\end{corollary}

\begin{definition}[Développement limité]
    Un polynôme $P_n(x)$ de degré $n$ satisfaisant :
    \[ f(x) = P_n(x) + o((x - x_0)^n) \]
    est un développement limité d'ordre $n$ de la fonction $f$.
\end{definition}

\begin{remark}
    Il est courant d'abréger développement limité par DL s'il n'y a pas d'ambiguïté.
\end{remark}

\begin{proposition}
    Si une fonction admet un développement limité, alors il est unique.
\end{proposition}

\subsection{Opérations sur les développements limités}
\begin{proposition}
    Soient $(c_0, \ldots, c_n) \in \R^{n+1}, (d_0, \ldots, d_n) \in \R^{n+1}$ et $f, g$ deux fonctions admettant des développements limités en 0 telles que :
    \[ f(x) = c_0 + c_1x + \cdots + c_n x^n + o(x^n) \]
    \[ g(x) = d_0 + d_1x + \cdots + d_n x^n + o(x^n) \]
    \begin{enumerate}
        \item Addition : $f + g$ admet un développement limité en 0 à l'ordre $n$.
        \[ f(x) + g(x) = (c_0 + d_0) + (c_1 + d_1)x + \cdots + (c_n + d_n) x^n + o(x^n) \]
        \item Multiplication : $f \cdot g$ admet un développement limité en 0 à l'ordre $n$.
        \[ (c_0 + c_1 x + \cdots + c_n x^n) \cdot (d_0 + d_1 x + \cdots + d_n x^n) \]
        où l'on conserve les monômes de degré inférieur ou égal à $n$.
        \item Composition : Si $g(0) = 0$ alors la fonction $f \circ g$ admet un développement limité en 0 à l'ordre $n$.\\
        Posons $C(x) \coloneqq c_0 + c_1x + \cdots + c_n x^n$ et $D(x) \coloneqq d_0 + d_1x + \cdots + d_n x^n$.\\
        Sa partie polynomiale est le polynôme tronqué à l'ordre $n$ (on conserve les monômes de degré inférieur ou égal à $n$) de la composition $C(D(x))$. 
    \end{enumerate}
\end{proposition}


\part{Algèbre}
\chapter{Calcul Algèbrique}
\def\arraystretch{1}

\par \noindent Dans cette partie, $\K$ désigne soit $\N$, soit $\Z$, soit $\Q$, soit $\R$.

\begin{proposition}
    Soient $a, b, c \in \K$ :
    \begin{enumerate}
    	\item Associativité :
    	\begin{enumerate}
    		\item $a + (b + c) = (a + b) + c$.
    		\item $a \cdot (b \cdot c) = (a \cdot b) \cdot c$.
    	\end{enumerate}
    	\item Commutativité : 
    	\begin{enumerate}
    		\item $a + b = b + a$.
    		\item $a \cdot b = b \cdot a$.
    	\end{enumerate}
    	\item \'Elément neutre :
    	\begin{enumerate}
    		\item $a + 0 = a$.
    		\item $a \cdot 1 = a$.
    	\end{enumerate}
    	\item \'Elément absorbant : $a \cdot 0 = 0$.
    	\item \'Elément symétrique si $\K \neq \N$ : $\exists a' \in \K : a + a' = 0$
    	\item \'Elément inverse si $\K \in \{\Q, \R\}$ : $\exists a' \in \K : a \cdot a' = 1$.
    	\item Distributivité :
    	\begin{enumerate}
    		\item $a \cdot (b + c) = a \cdot b + a \cdot c$.
    		\item $(a + b) \cdot c = a \cdot c + b \cdot c$.
    	\end{enumerate}
    \end{enumerate}
\end{proposition}

\begin{proposition}[Opérations sur les fractions]
    Soient $a, c \in \Z,\ b, d \in \Z^*$.
    \begin{enumerate}
    	\item $\frac{a}{b} + \frac{c}{d} = \frac{ad + bc}{bd}$.
    	\item $\frac{a}{b} \cdot \frac{c}{d} = \frac{a \cdot c}{b \cdot d}$.
    \end{enumerate}
\end{proposition}

\begin{proof}
	Soient $a, a', c, c' \in \N$ et $b, b', d, d' \in \N^*$ tels que :
	\begin{align*}
		\frac{a}{b} &= \frac{a'}{b'} & \frac{c}{d} &= \frac{c'}{d'}
	\end{align*}
	ce qui revient à dire :
	\begin{align*}
		ab' &= a'b & cd' &= c'd
	\end{align*}
	\begin{enumerate}
		\item Montrons que :
		\[  \]
		ce qui revient à montrer que :
		\begin{align*}
			&\frac{ad + bc}{bd} = \frac{a'd' + b'c'}{b'd'} \\
			\iff &(ad + bc)(b'd') = (a'd' + b'c')(bd)
			\\
			\iff &(ad + bc)(b'd') - (a'd' + b'c')(bd) = 0
		\end{align*}
		\begin{align*}
			(ad + bc)(b'd') - (a'd' + b'c')(bd) &= adb'd' + bcb'd' - a'd'bd - b'c'bd \\
			&= (ab' - a'b)(dd') + (cd' - c'd)(bb')
		\end{align*}
		Sachant que $ab' = a'b$ et $cd' = c'd$, on a $ab' - a'b = 0$ et $cd' - c'd = 0$.
		\\
		Ainsi :
		\[ (ab' - a'b)(dd') + (cd' - c'd)(bb') = 0 \]
		\item Montrons que :
		\begin{align*}
			&\frac{ac}{bd} = \frac{a'c'}{b'd'} \\
			\iff &(ac)(b'd') = (a'c')(bd) \\
			\iff &(ac)(b'd') - (a'c')(bd) = 0
		\end{align*}
		\begin{align*}
			(ac)(b'd') - (a'c')(bd) &= acb'd' - a'c'bd + a'bcd' - a'bcd' \\
			&= (ab' - a'b)(cd') + (cd' - c'd)(a'b) 
		\end{align*}
		Sachant que $ab' = a'b$ et $cd' = c'd$, on a $ab' - a'b = 0$ et $cd' - c'd = 0$.
		\\
		Ainsi :
		\[ (ab' - a'b)(cd') + (cd' - c'd)(a'b) = 0 \]
	\end{enumerate}
\end{proof}

\begin{definition}[Somme]
    Soient $a_k \in \R$, $m, n \in \N$ tels que $m \leq n$.
    \[ \sum_{k = m}^n a_k = a_m + a_{m + 1} + \cdots + a_n \]
\end{definition}

\begin{proposition}[Linéarité de la somme]
    $\forall m, n \in \N,\ m \leq n,\ a_k, b_k, \lambda \in \R$ :
    \[ \sum_{k = m}^{n} (a_k + \lambda b_k) = \sum_{k = m}^n a_k + \lambda \sum_{k = m}^n b_k \]
\end{proposition}

\begin{proof}
	Nous pouvons le vérifier en développant les sommes.
\end{proof}

\begin{proposition}[Somme téléscopique]
    $\forall m, n \in \N,\ m \leq n,\ a_k \in \R$ :
    \[ \sum_{k = m}^n (a_k - a_{k - 1}) = a_n - a_{m - 1} \]
\end{proposition}

\begin{proof}
	Nous pouvons le vérifier en développant la somme.
\end{proof}

\begin{proposition}
	$\forall n \in \N$ :
	\begin{align*}
		\sum_{k = 0}^{n} k &= \frac{n(n+1)}{2} & 
		\sum_{k = 0}^{n} k^2 &= \frac{n(n+1)(2n+1)}{6} &
		\sum_{k = 0}^{n} k^3 &= \frac{n^2(n+1)^2}{4}
	\end{align*}
\end{proposition}

\begin{proof}
	On procède par récurrence simple.
\end{proof}

\begin{definition}
	Pour tout $n \in \N$ :
	\[ n! = n \cdot (n - 1) \cdot \ldots \cdot 2 \cdot 1 \]
\end{definition}

\begin{proposition}
	$\forall n, k \in \N,\ n \geq k$ :
	\[ \binom{n}{p} = \frac{n!}{p!(n - p)!} \]
\end{proposition}

\begin{proposition}
	Pour tous $n, k \in \N$ tels que $n \geq k$.
	\[
	\binom{n}{k - 1} + \binom{n}{k} = \binom{n+1}{k}
	\]
\end{proposition}

\begin{proof}
	D'une part :
	\[ \binom{n}{k - 1} = \frac{n!}{(k-1)!(n - (k - 1))!} \]
	D'autre part :
	\[ \binom{n}{k} = \frac{n!}{k!(n - k)!} \]
	On a ensuite : 
	\begin{align*}
		\binom{n}{k - 1} + \binom{n}{k} &= 
		\frac{n!}{(k-1)!(n - (k - 1))!} + \frac{n!}{k!(n - k)!} \\
		&= \frac{n!}{(k-1)!(n - k + 1))!} + \frac{n!}{k!(n - k)!} \\
		&= \frac{n!k}{k!(n-k+1)!} + \frac{n!(n-k+1)}{k!(n-k+1)!} \\
		&= \frac{n!k + n!(n - k + 1)}{k!(n - k + 1)!} \\
		&= \frac{n!(k + n - k + 1)}{k!(n - k + 1)!} \\
		&= \frac{n!(n+1)}{(n - k + 1)!} \\
		&= \frac{(n+1)!}{k!(n + 1 - k)!} \\
		&= \binom{n+1}{k}
	\end{align*}
\end{proof}
	
\begin{proposition}[Binôme de Newton]
	$\forall a, b \in \K,\ n \in \N$ :
	\begin{align*}
		(a + b)^n = \sum_{k = 0}^{n} \binom{n}{k} a^kb^{n-k}
	\end{align*}
\end{proposition}

\begin{proof}
	Procédons par récurrence sur $n \in \N$ pour montrer $P(n) : \forall a, b \in \K,\ n \in \N : (a + b)^n = \sum_{k = 0}^{n} \binom{n}{k} a^kb^{n-k}$.
	\begin{enumerate}
		\item \textbf{Initialisation} : Pour $n = 0$.
		\\
		D'une part :
		\[ (a + b)^0 = 1 \]
		D'autre part :
		\[ \sum_{k = 0}^{0} \binom{0}{k} a^k b^{-k} = 1 \]
		$P(0)$ est vraie.
		\item \textbf{Hérédité} : Supposons que $P(n)$ est vraie pour un $n \geq 1$ fixé.
		\\
		\begin{align*}
			(a + b)^n &= \sum_{k=0}^{n} \binom{n}{k} a^k b^{n-k} \\
			(a + b)^{n+1} &= (a+b) \sum_{k=0}^{n} \binom{n}{k} a^k b^{n-k} \\
			&= \sum_{k = 0}^{n} \binom{n}{k} a^{k+1} b^{n-k} + \sum_{k=0}^{n} \binom{n}{k} a^k b^{n - k + 1}
		\end{align*}
		Procédons à un changement de variable $j = k = 1$.
		\begin{align*}
			(a + b)^{n+1} &= \sum_{j = 1}^{n + 1} \binom{n}{j - 1} a^j b^{n - j + 1} + \sum_{k = 0}^{n} \binom{n}{k} a^k b^{n-k+1} \\
			&= \sum_{k = 1}^{n + 1} \binom{n}{k - 1} a^k b^{n - k + 1} + \sum_{k = 0}^{n} \binom{n}{k} a^k b^{n - k + 1} \\
			&= \binom{n}{n} a^{n+1} + \binom{n}{0} b^{n+1} + \sum_{k = 1}^{n} \binom{n}{k - 1} a^k b^{n - k + 1} + \sum_{k = 1}^{n} \binom{n}{k} a^k b^{n - k + 1} \\
			&= a^{n+1} + b^{n+1} + \sum_{k = 1}^{n} \left[ \binom{n}{k - 1} + \binom{n}{k} \right] a^k b^{n - k + 1} \\
			&= a^{n+1} + b^{n+1} + \sum_{k=1}^{n} \binom{n+1}{k} a^k b^{n - k + 1} \\
			&= \sum_{k=0}^{n+1} \binom{n + 1}{k} a^k b^{n + 1 - k}
		\end{align*}
		Ce qui montre que $P(n+1)$ est vraie. 
		\\
		Par principe de récurrence, la propriété $P(n)$ est vraie pour tout $n \in \N$.
	\end{enumerate}
\end{proof}

\begin{definition}[Produit]
    $\forall m, n \in \N,\ m \leq n,\ a_k \in \R$ :
    \[ \prod_{k = m}^{n} = a_m \cdot a_{m+1} \cdot \ldots \cdot a_n \]
\end{definition}
\chapter{Ensembles}
\def\arraystretch{1}

Nous allons tout d'abord donner une définition intuitive d'un ensemble : Un ensemble $E$ est une collection d'objets appelés éléments. Si $E$ contient un élément $x$, on dit que $x$ appartient à $E$, noté $x \in E$.

\begin{definition}[Ensemble vide]
  L'ensemble vide noté $\varnothing$ est l'ensemble ne contenant aucun élément.
\end{definition}

\begin{definition}[Inclusion]
  Soient $E$ et $F$ deux ensembles.
  $$F \subseteq E \iff \forall x \in F,\ x \in E$$
  On dit que $F$ est inclu dans $E$.
\end{definition}

\begin{definition}[\'Egalité d'ensembles]
  Soient $E$ et $F$ deux ensembles.
  $$E = F \iff E \subseteq F \text{ et } F \subseteq E$$
\end{definition}

\begin{definition}[Singleton]
  Un singleton est un ensemble ne contenant qu'un seul élément.
\end{definition}

\begin{definition}[Réunion d'ensembles]
  Soient $E$ et $F$ deux ensembles.
  \[ E \cup F = \{ x : x \in E \text{ ou } x \in F \} \]
  On lit \og $E$ union $F$ \fg.
\end{definition}

\begin{definition}[Intersection d'ensembles]
  Soient $E$ et $F$ deux ensembles.
  \[ E \cap F = \{ x : x \in E \text{ et } x \in F \} \]
  On lit \og $E$ inter $F$ \fg.
\end{definition}

\begin{definition}[Complémentaire d'un ensemble]
  Soient $E$ et $F$ deux ensembles. 
  \\
  Le complémentaire de $F$ dans $E$ noté $E \backslash F$ ou $\complement_E F$ est défini par :
  \[ E \backslash F = \{ x : x \in E,\ x \notin F \} \]
  Si on parle du complémentaire sans préciser d'ensemble, on peut se permettre de noter $E^C$ ou $\complement E$.
\end{definition}

\begin{proposition}
	Soient $A, B, C$ et $E$ des ensembles.
	\begin{enumerate}
		\item La réunion et l'intersection sont commutatives et assocatives.
		\item \'Elément neutre :
		\begin{enumerate}
			\item $A \cup \varnothing = A$.
			\item $A \cap A = A$.
		\end{enumerate}
		\item $A \subseteq E \iff A \cap E = E \cap A = A$.
		\item Distributivité :
		\begin{enumerate}
			\item $A \cup (B \cap C) = (A \cup B) \cap (A \cup C)$.
			\item $A \cap (B \cup C) = (A \cap B) \cup (A \cap C)$.
		\end{enumerate}
	\end{enumerate}
\end{proposition}

\begin{proposition}[Lois de Morgan]
  Soient $A$ et $B$ deux ensembles.
  \begin{align*}
    (A \cup B)^{C} &= A^{C} \cap B^{C} & (A \cap B)^{C} &= A^{C} \cup B^{C}
  \end{align*}
\end{proposition}

\begin{proof}
	Soient $A$ et $B$ deux ensembles et $x$ un élément quelconque. 
	\begin{enumerate}
		\item \boxed{\subseteq} : Par définition du complémentaire :
		\[ x \in (A \cup B)^C \iff x \notin (A \cup B) \]
		$x \notin A$ car $A \subseteq (A \cup B)$ et $x \notin B$ car $B \subseteq (A \cup B)$, ainsi $ \in A^C$ et $x \in B^C$. 
		\\
		Par définition de l'intersection :
		\[ x \in (A^C \cup B^C) \]
		et donc :
		\[ (A \cup B)^C \subseteq (A^C \cap B^C) \]
		\boxed{\supseteq} : Par définition de l'intersection :
		\begin{align*}
			x \in (A^C \cap B^C) &\iff x \in A^C \text{ et } x \in B^C \\
			&\iff x \notin A \text{ et } x \notin B \\
			&\iff x \notin (A \cup B)^C
		\end{align*}
		d'où :
		\[ (A^C \cap B^C) \subseteq (A \cup B)^C \]
		\item \boxed{\subseteq} : Par définition du complémentaire :
		\begin{align*}
			x \in (A \cap B)^C &\iff x \notin (A \cap B) \\
			&\iff x \notin A \text{ et } x \notin B \\
			&\iff x \in A^C \text{ et } x \in B^C \\
			&\iff x \in (A^C \cap B^C)
		\end{align*}
		Sachant que :
		\[ (A^C \cap B^C) \subseteq (A^C \cup B^C) \]
		On a :
		\[ x \in (A^C \cap B^C) \implies x \in (A^C \cup B^C) \]
		d'où :
		\[ (A \cap B)^C \subseteq (A^C \cup B^C) \]
		\boxed{\supseteq} : Par définition de la réunion :
		\begin{align*}
			x \in (A^C \cup B^C) &\iff x \in A^C \text{ ou } x \in B^C \\
			&\iff x \notin A \text{ ou } x \notin B \\
			&\iff x \notin (A \cap B) \\
			&\iff x \in (A \cap B)^C
		\end{align*}
		Ainsi : 
		\[ (A^C \cap B^C) \subseteq (A \cup B)^C \]
	\end{enumerate}
\end{proof}

\begin{definition}[Produit cartésien]
    Soient $E$ et $F$ des ensembles. On définit le produit cartésien :
    \[ E \times F = \{ (x, y),\ x \in E,\ y \in F \} \]
    Par convention : $\underbrace{E \times \cdots \times E}_{n \text{ fois}} = E^n$.
\end{definition}

\chapter{Logique et raisonnements}
\def\arraystretch{1.5}
\section{Logique}
\noindent Notations logiques :
\begin{enumerate}
	\item $\neg$ : \og non \fg.
	\item $\land$ : \og et \fg.
	\item $\lor$ : \og ou \fg.
	\item $\veebar$ : \og ou exclusif \fg.
\end{enumerate}

\begin{definition}[Assertion]
  Une assertion est une affirmation mathématique soit vraie soit fausse.
\end{definition}

\begin{definition}[Prédicat]
  Un prédicat est un énoncé mathématique dont la véracité dépend d'une ou plusieurs variables.
\end{definition}

Soient $P$ et $Q$ deux prédicats.
\begin{table}[!h]
	\centering
	\begin{tabular}{ccccccc}
		\toprule
		$P$ & $Q$ & $P \land Q$ & $P \lor Q$ & $\neg P$ & $P \implies Q$ & $P \veebar Q$ \\
		\midrule
		$V$ & $V$ & $V$ & $V$ & $F$ & $V$ & $F$ \\
		$V$ & $F$ & $F$ & $V$ & $F$ & $F$ & $V$ \\
		$F$ & $V$ & $F$ & $V$ & $V$ & $V$ & $V$ \\
		$F$ & $F$ & $F$ & $F$ & $V$ & $V$ & $F$ \\
		\bottomrule
	\end{tabular}
	\caption{Table de vérité}
\end{table}

\begin{proposition}
  Soient $P$ et $Q$ deux prédicats.  
	\begin{enumerate} 
		\item $(P \implies Q) \land (Q \implies P) \implies (P \iff Q)$.
		\item $P \implies Q \iff \neg P \lor Q$.
		\item $\neg (P \lor Q) \iff \neg P \land \neg Q$.
		\item $\neg (P \land Q) \iff \neg P \lor \neg Q$.
		\item $P \implies Q \iff \neg Q \implies \neg P$.
	\end{enumerate}
\end{proposition}

\begin{remark}
    Les 2. et 3. sont les lois de Morgan, la 4. est la contraposée.
\end{remark}

\par \noindent Voici quelques négations usuelles :
\begin{itemize}
    \item Le contraire de \og $\forall x \in E, P(x)$ \fg est \og $\exists x \in E, \neg P(x)$ \fg.
    \item Le contraire de \og $x < y$ \fg est \og $x \geq y$ \fg.
\end{itemize}

\section{Raisonnements}

\begin{definition}[Raisonnement par récurrence simple]
    Il existe plusieurs variantes du raisonnement par récurrence, définissons d'abord la récurrence simple. L'objectif est de montrer qu'une propriété $P_n$ est vraie pour tout entier naturel $n$. 
    \begin{enumerate}
        \item \textbf{Initialisation} : On montre que $P_0$ est vraie.
        \item \textbf{Hérédité} : On suppose que pour un $k$ tel que $0 < k < n,\ P_k$ est vraie et on montre que $P_{k+1}$ est vraie.
    \end{enumerate}
\end{definition}

\begin{definition}[Raisonnement par analyse-synthèse \cite{analyse_synthese_bibmath}]
	Raisonnement utilisé pour démontrer l'\textbf{existence} et l'\textbf{unicité} d'un objet.
	\begin{enumerate}
		\item \textbf{Analyse} : On suppose que l'objet existe et on cherche les conditions nécessaires que doit vérifier l'objet. Cette partie démontre l'\textbf{unicité}.
		\item \textbf{Synthèse} : On considère l'objet identifié dans la partie analyse et on vérifie qu'il a les propriétés souhaitées. Cette partie démontre l'\textbf{existence}.
	\end{enumerate}
\end{definition}

Voici quelques types de raisonnements utilisés pour démontrer une proposition du type \og Si $P$ vraie alors $Q$ vraie \fg.

\begin{definition}[Raisonnement direct]
	Il suffit de raisonner à partir de l'hypothèse \og $P$ vraie \fg pour montrer que \og $Q$ vraie \fg . 
\end{definition}

\begin{definition}[Raisonnement par contraposition]
	Parfois il est plus simple de montrer que $\neg Q \implies \neg P$.
	Donc on suppose que $\neg Q$ est vraie et on démontrer que $\neg P$ est vraie.
\end{definition}

\begin{definition}[Raisonnement par l'absurde]
	Pour montrer que $P \implies Q$, il est possible de montrer que $P \implies \neg Q$ est absurde.
	\\
	On suppose que $P$ vraie mais que $Q$ fausse, on raisonne jusqu'à aboutir à une contradiction.
\end{definition}
\chapter{Nombres complexes}
\def\arraystretch{1}
On définit l'ensemble des nombres complexes, noté $\C$, comme une extension de l'ensemble des nombres réels.
Cette extension introduit un nouvel élément, noté $i$, appelé \textbf{nombre imaginaire} et défini comme $i^2=-1$.

\section{Vision algébrique des nombres complexes}
\begin{definition}[Forme algébrique d'un nombre complexe]
    Soient $a, b \in \R$ et $z \in \C$, on appelle \textbf{forme algébrique} de $z$ l'expression $z = a + ib$.
    \\
    $a$ est appelé \og \textbf{partie réelle} \fg, notée $\Re(z)$ et $b$ est appelé \og \textbf{partie imaginaire} \fg, notée $\Im(z)$.
\end{definition}

\begin{proposition}
	Soient $z_1, z_2, z_3 \in \C$.
	\begin{enumerate}
		\item Associativité : 
		\begin{enumerate}
			\item $(z_1 + z_2) + z_3 = z_1 + (z_2 + z_3)$.
			\item $(z_1 \cdot z_2) \cdot z_3 = z_1 \cdot (z_1 \cdot z_3)$.
		\end{enumerate}
		\item Commutativité :
		\begin{enumerate}
			\item $z_1 + z_2 = z_2 + z_1$.
			\item $z_1 \cdot z_2 = z_2 \cdot z_1$.
		\end{enumerate}
		\item \'Elément neutre :
		\begin{enumerate}
			\item $0 + i0 = 0$.
			\item $z_1 + 0 = z_1$.
			\item $z_1 \cdot 1 = z_1$
		\end{enumerate}
		\item \'Elément absorbant : $z_1 \cdot 0 = 0$.
		\item \'Elément symétrique : $\exists z' \in \C : z_1 + z' = 0$.
		\item \'Elément inverse : $\exists z' \in \C :  z_1 \cdot z' = 1$.
		\item Distributivité :
		\begin{enumerate}
			\item $(z_1 + z_2) \cdot z_3 = z_1 \cdot z_3 + z_2 \cdot z_3$.
			\item $z_1 \cdot (z_2 + z_3) = z_1 \cdot z_2 + z_1 \cdot z_3$.
		\end{enumerate}
	\end{enumerate}
\end{proposition}

\begin{definition}[Module d'un nombre complexe]
    Soit $z = a + ib$ avec $a,\ b \in \R$. 
    \\
    On définit $\abs{z}$ tel que :
    \[ \abs{z} = \sqrt{a^2 + b^2} \]
    qu'on appelle \textbf{module} de $z$.
\end{definition}

\begin{definition}[Conjugué d'un nombre complexe]
    Soit $z = a + ib$ avec $a,\ b \in \R$. On appelle conjugué de $z$ qu'on note $\overline{z}$ tel que :
    \[ \overline{z} = a - ib \]
\end{definition}

\begin{proposition}
    Soient $z_1, z_2 \in \C$.
    \begin{enumerate}
            \item $\abs{z_1 + z_2} \leq \abs{z_1} + \abs{z_2}$.
            \item $\abs{z_1 - z_2} \geq \abs{z_1} - \abs{z_2}$.
            \item $\abs{z_1 \cdot z_2} = \abs{z_1} \cdot \abs{z_2}$.
            \item Si $z_1 \neq 0$ : $\abs{\frac{1}{z_1}} = \frac{1}{\abs{z_1}}$.
            \item Si $z_2 \neq 0$ : $\abs{\frac{z_1}{z_2}} = \frac{\abs{z_1}}{\abs{z_2}}$.
            \item $\abs{z}^2 = z \cdot \overline{z}$.
            \item $\abs{z} \geq 0$.
            \item $\abs{z} = 0 \iff z = 0$.
            \item $\abs{z} = \abs{\overline{z}} = \abs{-z} = \abs{-\overline{z}}$.
        \end{enumerate}
\end{proposition}

\section{Vision géométrique des nombres complexes}
Il est possible de représenter les nombres complexes sur un plan complexe avec l'axe des ordonnées représentant la partie imaginaire et l'axe des abscisses la partie réelle.

\begin{definition}
	Soit $z \in \C$, on appelle l'\textbf{argument} de $z$, noté $\arg(z)$, l'angle entre l'axe de la partie réelle et la droite issue de l'origine passant par $z$. 
\end{definition}

\begin{figure}[!h]
	\centering
	\begin{tikzpicture}
		\coordinate (O) at (0, 0);
		\coordinate (Z) at (4, 3);
		\coordinate (X) at (5, 0);
		\coordinate (Y) at (0, 5);
		\draw[step=1cm, gray, very thin] (-2, -2) grid (6,6);
		\draw[thick, ->] (O) -- (X) node[anchor=south east] {$\Re(z)$};
		\draw[thick, ->] (O) -- (Y) node[anchor=north west] {$\Im(z)$};
		\draw[->] (O) -- (Z) node[circle, fill, inner sep=0pt, scale=0.5, label=above:$z$]{$z$};
		\draw (1.9, 1.6) node[anchor=south]{$\abs{z}$};
		\path[gray]
		pic["$\arg(z)$" shift={(30pt, 5pt)}, draw, -, angle radius=1cm] {angle = X--O--Z};
	\end{tikzpicture}
	\caption{Vision géométrique des nombres complexes}
\end{figure}

\begin{proposition}
	Soient $z_1, z_2 \in \C$ et $n \in \N$.
	\begin{enumerate}
		\item $\arg(z_1 \cdot z_2) = \arg(z_1) + \arg(z_2)$.
		\item $\arg(z_1^n) = n \arg(z_1)$.
		\item Si $z_1 \neq 0$ : $\arg \left(\frac{1}{z_1} \right) = - \arg(z_1)$.
		\item Si $z_2 \neq 0$ : $\arg \left( \frac{z_1}{z_2} \right) = \arg(z_1) - \arg(z_2)$.
	\end{enumerate}
\end{proposition}

\begin{definition}
	Soient $z \in \C,\ r = \abs{z}$ et $\theta = \arg(z)$, il est possible d'exprimer $z$ dans sa forme trigonométrique :
	\[ z = r \left( \cos(\theta) + i \sin(\theta) \right) \]
\end{definition}

\begin{proposition}
	Soient $z_1 = r_1 \left( \cos(\theta_1) + i \sin(\theta_1) \right)$ et $z_2 = r_2 \left( \cos(\theta_2) + i \sin(\theta_2) \right)$ deux nombres complexes tels que :
	\begin{enumerate}
			\item $r_1 = \abs{z_1}$.
			\item $r_2 = \abs{z_2}$.
			\item $\theta_1 = \arg(z_1)$.
			\item $\theta_2 = \arg(z_2)$.
		\end{enumerate}
	\[ z_1 z_2 = r_1 r_2 \left( \cos(\theta_1+\theta_2) + i \sin(\theta_1+ \theta_2) \right) \]
\end{proposition}

\begin{proof}
	\begin{align*}
		z_1 z_2 &= r_1(\cos(\theta_1) + i \sin(\theta_1)) r_2(\cos(\theta_2) + i \sin(\theta_2)) \\
		&= (r_1 \cos(\theta_1) + i r_1 \sin(\theta_1))(r_2 \cos(\theta_2) + i r_2 \sin(\theta_2)) \\
		&= r_1 r_2 \cos(\theta_1) \cos(\theta_2) + i r_1 r_2 \cos(\theta_1) \sin(\theta_2) + i r_1 r_2 \sin(\theta_1) \cos(\theta_2) - r_1 \sin(\theta_1) \sin(\theta_2) \\
		&= r_1 r_2 (\cos(\theta_1) \cos(\theta_2) - \sin(\theta_1) \sin(\theta_2)) + i (r_1 r_2(\cos(\theta_1) \sin(\theta_2) + \sin(\theta_1)\cos(\theta_2))) \\
		&= r_1 r_2 (\cos(\theta_1 + \theta_2) + i \sin(\theta_1 + \theta_2))
	\end{align*}
\end{proof}

\begin{definition}
	Soient $z \in \C,\ r = \abs{z}$ et $\theta = \arg(z)$ tels que :
	\[ z = r (\cos(\theta) + i \sin(\theta)) \]
	On peut écrire $z$ sous une forme utilisant l'exponentielle :
	\[ z = re^{i \theta} \]
\end{definition}

\begin{proposition}[Formule de Moivre]
	Soient $\theta \in \R$ et $n \in \N$.
	\[ (\cos(\theta) + i \sin(\theta))^n = \cos(n \theta) + i \sin(n \theta) \]
\end{proposition}

\begin{proposition}[Identité d'Euler]
	\[ e^{i\pi} + 1 = 0 \]
\end{proposition}

\begin{proposition}[Formules d'Euler]
	Soit $\theta \in \R$.
	\begin{enumerate}
			\item $\cos(\theta) = \frac{e^{i\theta} + e^{-i\theta}}{2}$.
			\item $\sin(\theta) = \frac{e^{i\theta} - e^{-i\theta}}{2i}$.
		\end{enumerate}
\end{proposition}

\begin{definition}
	Soient $z \in \C$ et $n \in \N$. On appelle \textbf{racine $n$-ième} de $z$ tout $\omega \in \C$ tel que :
	\[ \omega^n = z \] 
\end{definition}

\begin{proposition}
	Soient $z \in \C^*$, $\theta = \arg(z)$ et $\rho = \abs{z}$ tels que :
	\[ z = \rho e^{i\theta} \]
	$z$ admet $n$ racines $n$-ièmes, pour $0 \leq k \leq n - 1$ :
	\[ \omega_k = \rho^{\frac{1}{n}} e^{i \left( \frac{\theta}{n} + \frac{2 k \pi}{n} \right)} \]
\end{proposition}

\begin{proof}
	Soit $z = \rho e^{i \theta}$. \\
	Les racines $n$-ièmes de $z$ sont les nombres $\omega = re^{i \theta'}$ pour $r = \abs{\omega}$ et $\theta' = \arg(\omega)$ tels que $\omega^n = z$.
	\begin{align*}
		(r e^{i \theta'})^n &= \rho e^{i \theta} \\
		r^n (e^{i \theta'})^n &= \rho e^{i \theta} \\
		r^n e^{i n \theta'} &= \rho e^{i \theta}
	\end{align*} 
	Par identification nous avons pour $0 \leq k \leq n - 1$:
	\begin{align*}
		\begin{cases}
			r^n = \rho \\
			n \theta' = \theta + 2 k \pi
		\end{cases}
		\iff 
		\begin{cases}
			r = \rho^{\frac{1}{n}} \\
			\theta' = \frac{\theta}{n} + \frac{2 k \pi}{n}
		\end{cases}
	\end{align*}
	On a finalement : 
	\[ \omega_k = \rho^{\frac{1}{n}} e^{i \left( \frac{\theta}{n} + \frac{2 k \pi}{n} \right) } \]
\end{proof}

\section{Géométrie des nombres complexes}
\begin{proposition}
	Soit :
	\begin{center}
		$
		\appli{f}{\C}{\C}{z}{f(z)}
		$
	\end{center}
	\begin{enumerate}
		\item Soit $a \in \C,\ f(z) = z + a$ : translation d'affixe $a$.
		\item Soit $a \in \R^*,\ f(z) = az$ : homothétie de rapport $a$.
		\item Soient $a \in \C$ et $\theta \in \R,\ f(z) = (z-a)e^{i\theta} + a$ : rotation d'angle $\theta$ et de centre $a$.
		\item Soit $\theta \in \R,\ f(z) = \overline{z} e^{2i\theta}$ : réflexion par rapport à la droite formant un angle $\theta$ avec l'axe des réels.
	\end{enumerate}
\end{proposition}

\begin{proposition}
	
	\begin{enumerate}
		\item L'axe des réels : $\overline{z} = z$.
		\item Un axe formant une angle $\theta$ avec l'axe des réels : $\overline{e^{-i \theta}z} = e^{-i \theta} z$.
		\item L'asymptote verticale de partie réelle $a$ : $z + \overline{z} = 2a$.
	\end{enumerate}
\end{proposition}
\chapter{Arithmétique}
\def\arraystretch{1}

\begin{definition}
    Soient $a \in \Z$ et $b \in \Z^*$.
    \[ a \text{ est un multiple de } b \iff b \text{ est un diviseur de } a \iff b \mid a \iff \exists q \in \Z, a = bq \]
\end{definition}

\section{Divisibilité}

\begin{theorem}[Division euclidienne]
    Soient $a \in \Z$ et $b \in \Z^*$.
    \[ \exists ! (q, r) \in \Z^2,\ 0 \leq r < \abs{b} : a = bq + r  \]
    \begin{enumerate}
    	\item $a$ est appelé le \textbf{dividende}
        \item $b$ est appelé le \textbf{diviseur}
        \item $q$ est appelé le \textbf{quotient}
        \item $r$ est appelé le \textbf{reste}
    \end{enumerate}
\end{theorem}

\begin{proof}\cite{livre_prepa}
	\begin{enumerate}
		\item \textbf{Existence :} Supposons $a \in \N$ et considérons $M = \{n \in \N : nb \leq a\}$ l'ensemble des multiples de $b$ inférieurs à $a$. $M$ est une partie de $\N$. Nous avons deux propriétés : 
		\begin{enumerate}
			\item $M$ est non vide car 0 est un multiple de $b$ inférieur à $a$.
			\item $M$ est majoré par $a$ d'après sa définition.
		\end{enumerate}
		Ainsi $M$ admet un plus grand élément que l'on note $q$, vérifiant :
		\begin{enumerate}
			\item $qb \leq a$ car $q \in M$ 
			\item $(q + 1)b > a$ car $q + 1 > q$ sachant que $q$ est le plus grand élément de $M$, $q + 1 \notin M$.
		\end{enumerate}
		Posons : $r = a - bq$. Sachant que $a \geq bq,\ r \geq 0$. On a $r < b$ car $b = (q + 1)b - qb > a - bq = r$. Supposons que $a \in \Z$. 
		\begin{enumerate}
			\item Si $a$ est positif, on se ramène au cas précédent.
			\item Dans le cas où $a < 0,\ -a \geq 0$, ainsi il existe $(q', r') \in \Z^2$ tel que :
		\[ -a = bq' + r',\ 0 \leq r' < \abs{b} \]
		\[a = b(-q') - r' \]
			\begin{enumerate}
				\item Si $r' = 0$, on pose $q = -q'$ et $r = 0$ et on obtient le couple recherché.
				\item Si $r' \neq 0$, $r' \in \llbracket 1, b-1 \rrbracket$ et $a = b(-q' -1) + (b - r')$, on pose $q = -q' - 1$ et $r = b - r'$ et on obtient le couple recherché. 
			\end{enumerate}
		\end{enumerate}
		\item \textbf{Unicité :} Soient $q,\ q',\ r,\ r' \in \Z$. \\
	On a d'une part : $a = bq + r$ et d'autre part : $a = bq' + r'$.
	On sait que $0 \leq r < b$ et $0 \leq r' < b$ donc :
	\[ b \abs{q' - q} = \abs{r' - r} < b \]
	ce qui n'est possible que si $\abs{q' - q} = 0$ ce qui impliquerait $q = q'$. Ceci entraîne donc $r = r'$.
	\end{enumerate}
\end{proof}

\begin{example}
	Donnons la division euclidienne de $5$ par $2$.
	\[ 5 = 2 \cdot 2 + 1 \]
\end{example}

\section{PGCD et PPCM}

\begin{definition}
	Soient $a, b \in \Z^*$. 
	\begin{enumerate}
		\item L'ensemble des diviseurs de $\N^*$ commun à $a$ et $b$ admet un plus grand élément. C'est le \textbf{plus grand commun diviseur} des entiers $a$ et $b$. On le note $\pgcd(a, b)$.
		\item L'ensemble des diviseurs de $\N^*$ commun à $a$ et $b$ admet un plus petit élément. C'est le \textbf{plus petit commun multiple} des entiers $a$ et $b$. On le note $\ppcm(a, b)$.
	\end{enumerate}
\end{definition}

\begin{theorem}
	Soient $a, b \in \Z^*$ et $d \in \Z$.
	\begin{enumerate}
		\item $a \mid d \text{ et } b \mid d \implies \ppcm(a, b) \mid d$
		\item $d \mid a \text{ et } d \mid b \implies d \mid \pgcd(a, b)$
	\end{enumerate}
\end{theorem}

\begin{proof}
	\leavevmode
	\begin{enumerate}
		\item Posons $\ell = \ppcm(a,b)$. 
		\[ \exists ! (q,r) \in \Z^2,\ d = q \ell + R,\ 0 \leq r < \ell \]
		\[ r = d - q \ell,\ d \text{ et } \ell \text{ sont multiples de } a \text{ et } r \text{ est aussi un multiple de } a \text{ et } b \]
		Par la minimalité de $\ell,\ r = 0 \implies m = q \ell$.
		\item Posons $m = \pgcd(a, b)$. Montrons que :
		\[ \pgcd(m, d) = m \]
		Soit $\ell = \ppcm(m,d)$, $\ell \geq m$, $a$ et $b$ sont multiples de $m$ et $d$. D'après 1. :
		\[ \ell \mid a \text{ et } \ell \mid b,\ \ell \leq m \]
		Sachant que $\ell \geq m$ et $\ell \leq m$, $\ell = m$.
	\end{enumerate}
\end{proof}

\begin{definition}
	Soient $a, b \in \Z^*$. 
	\\	
	On dit que $a$ et $b$ sont \textbf{premiers entre eux} si et seulement si $\pgcd(a, b) = 1$.
\end{definition}

\section{Algorithme d'Euclide}
\begin{proposition}[Algorithme d'Euclide]
	Soient $a, b \in \Z^*$ tels que $\abs{a} > \abs{b}$.
	\[ \exists ! (q, r) \in \Z^2,\ 0 \leq r < \abs{b} : a = bq + r,\  \]
	\[ \pgcd(a, b) = \pgcd(b, a) = \pgcd(b, a - qb) = \pgcd(b, r) \]
	\begin{enumerate}
		\item Si $r = 0$ alors $a = qb$ et donc $\pgcd(a, b) = b$.
		\item Si $r \neq 0$ alors :
		\[ \exists ! (q_1, r_1) \in \Z^2,\ 0 \leq r_1 < r : b = q_1 r + r_1 \]
		Ensuite : 
		\begin{enumerate}
			\item Si $r_1 = 0$ alors $b = q_1 r$ et donc $\pgcd(a, b) = r$.
			\item Si $r_1 \neq 0$ alors :
			\[ \exists ! (q_2, r_2) \in \Z^2,\ 0 \leq r_2 < r_1 : r = q_2 r_1 + r_2 \]
		\end{enumerate}
	\end{enumerate}
	On procède de cette manière jusqu'à obtenir un reste nul, le $\pgcd$ de $a$ et $b$ est le dernier reste non nul.
\end{proposition}

\begin{example}
	Déterminons $\pgcd(216, 126)$.
	\begin{align*}
		216 &= 126 \cdot 1 + 90 \\
		126 &= 90 \cdot 1 + 36 \\
		90  &= 36 \cdot 2 + 18 \\
		36  &= 18 \cdot 2 + 0
	\end{align*}
	\[ \pgcd(216, 126) = 18 \]
\end{example}

\begin{theorem}[Identité de Bézout]
	Soient $a, b \in \Z$.
	\[ \exists (u, v) \in \Z^2,\ au + bv = \pgcd(a, b) \]
\end{theorem}

\begin{example}
	On cherche $u, v \in \Z$ tels que :
	\[ 216u + 126v = 18 \]
	On utilise l'algorithme d'Euclide :
	\begin{align*}
		216 &= 126 \cdot 1 + 90 \\
		126 &= 90 \cdot 1 + 36 \\
		90  &= 36 \cdot 2 + 18 \\
		36  &= 18 \cdot 2 + 0
	\end{align*}
	\begin{align*}
		90 - 36 \cdot 2 &= 18 \\
		90 - (126 - 90) \cdot 2 &= 18 \\
		90 \cdot 3 - 126 \cdot 2 &= 18 \\
		(216 - 126) \cdot 3 - 126 \cdot 2 &= 18 \\
		216 \cdot 3 + 126 \cdot (-5) &= 18
	\end{align*}
\end{example}

\begin{corollary}
	Soient $a, b \in \Z^*$ et $d \in \Z$.
	\[ \exists (u, v) \in \Z^2,\ au + bv = d \iff \pgcd(a,b) \mid d \]
\end{corollary}

\begin{lemma}[Lemme de Gauss]
	Soient $a, b \in \Z^*$ tels que $\pgcd(a, b) = 1$.
	\[ \forall c \in \Z,\ a \mid bc \implies a \mid c \]
\end{lemma}

\begin{proof}
	\begin{align*}
		\pgcd(a, b) = 1 &\implies \exists(u, v) \in \Z^2,\ au + bv = 1 \\
		&\implies a(cu) + b(cv) = c \\ 
		&\implies \pgcd(a, bc) \mid c
	\end{align*}
\end{proof}

\section{Nombres premiers}
\begin{definition}[Nombre premier]
	Soit $p \in \N^*$. 
	\\
	On dit que $p$ est \textbf{premier} si et seulement s'il admet \textbf{exactement} deux diviseurs : $1$ et $p$.
\end{definition}

\begin{theorem}[Théorème d'Euclide]
	Il existe une infinité de nombres premiers.
\end{theorem}

\begin{proof}
	Supposons qu'il existe $k$ nombres premiers $p_1, \ldots, p_k$.
	\[ N = p_1 \cdots p_k + 1 \implies p_i \nmid N \]
\end{proof}

\begin{lemma}
	Soit $n \in \N$ tel que $n \geq 2$.
	\\
	Si $p$ est le plus petit diviseur de $n$ tel que $p > 2$ alors $p$ est premier.
\end{lemma}

\begin{theorem}[Décomposition en facteurs premiers]
	Soit $n \in \N^*$ tel que $n \geq 2$. 
	\\
	Il existe une unique écriture de $n$ sous la forme de :
	\[ p_1^{\alpha_1} \cdots p_k^{\alpha_k} \]
	\begin{enumerate}
		\item Pour $1 \leq i \leq k$, les $p_i$ sont premiers.
		\item $\alpha_i \in \N^*$.
		\item $p_1 < p_2 < \cdots < p_k$.
	\end{enumerate}
\end{theorem}

\begin{example}
	Décomposons $40$ en facteurs premiers.
	\begin{align*}
		40 = 20 \cdot 2 \\
		20 = 10 \cdot 2 \\
		10 = 5 \cdot 2
	\end{align*}
	Ainsi $40 = 2^3 \cdot 5$.
\end{example}

\begin{proposition}
	Soient $a, b \in \Z$ et $i, k \in \N$. Pour déterminer $\pgcd(a, b)$ on peut utiliser leurs décompositions en facteurs premiers.
	\[ a = p_1^{\alpha_1} \cdots p_k^{\alpha_k} \]
	\[ b = n_1^{\beta_1} \cdots n_i^{\beta_i} \]
	$\pgcd(a, b)$ correspond au produit des facteurs premiers communs.
\end{proposition}

\section{Congruences}
\begin{definition}
	Soient $a, b \in \Z$ et $n \in \N$ tel que $n \geq 2$.
	\\
	On dit que $a$ et $b$ sont \textbf{congrus modulo} $n$ s'il existe un $k$ tel que :
	\[ a - b = k n \]
	On écrit généralement $a \equiv b \pmod n$ ou $a \equiv b \ [n]$.
\end{definition}

\begin{proposition}
	Soient $a, b, c, d \in \Z$ et $k, n \in \N$ tel que $n \geq 2$.
	\begin{enumerate}
		\item $a \equiv a \ [n]$ 
		\item $a \equiv b \ [n] \iff b \equiv a \ [n]$ 
		\item $a \equiv b \ [n] \text{ et } b \equiv c  \ [n] \implies a \equiv c \ [n]$
		\item $a \equiv b \ [n] \text{ et } c \equiv d     \ [n] \implies a + c \equiv b + d \ [n] \text{ et } ac \equiv bd \ [n]$
		\item $a \equiv b \ [n] \implies a^k \equiv b^k \ [n]$
	\end{enumerate}
\end{proposition}

\begin{proof}
	Immédiate en utilisant la définition de la congruence.
\end{proof}

\begin{theorem}
	Soient $m_1, m_2 \in \N$ tel que :
	\begin{enumerate}
		\item $m_1 > 1$
		\item $\pgcd(m_1, m_2) = 1$
	\end{enumerate}
	Soient $a_1, a_2 \in \Z$ tel que pour $x \in \Z$ :
	\begin{enumerate}
		\item $x \equiv a_1 \ [m_1]$
		\item $x \equiv a_2 \ [m_2]$
	\end{enumerate}
	Notons $\mathcal{S}$ l'ensemble des solutions :
	$\exists k_0,\ k \in \Z$ tel que :
	\begin{enumerate}
		\item $\mathcal{S} = k_0 + k m_1 m_2$ 
		\item $0 \leq k_0 < m_1 m_2$
	\end{enumerate}
\end{theorem}


\chapter{Polynômes}

\section{Définitions}
\begin{definition}[Polynôme]
    Un \emph{polynôme} est un élément de l'ensemble 
    \begin{align*}
        \K[X] = \{ a_0 1 + a_1 X + a_2 X^2 + \cdots a_n X^n \mid a_i \in \K,\ n \in \N \}
    \end{align*}
    Soit $P \in \K[X]$, si $a_n \neq 0$, $n$ est le degré du polynôme, on le note 
    $ \deg(P) = n. $
\end{definition}

\begin{proposition}
	$\forall \lambda \in \K^*,\ (P, Q) \in (\K[X])^2$. 
	\begin{multicols}{2}
        \begin{enumerate}
    		\item $\operatorname{deg}(\lambda) = 0$.
    		\item $\operatorname{deg}(P \cdot Q) = \operatorname{deg}(P) + \operatorname{deg}(Q)$.
    		\item $\operatorname{deg}(P + Q) = \mathrm{max}\left(\operatorname{deg}(P),\ \operatorname{deg}(Q)\right)$.
            \item $\operatorname{deg}(0) = -\infty$.
    	\end{enumerate}
    \end{multicols}
\end{proposition}

\section{Arithmétique des polynômes}

\begin{theorem}[Division euclidienne de polynômes]
	Soient $P_1, P_2$ deux polynômes non nuls.
	\begin{align*}
		\exists ! (Q, R) \in (\K[X])^2 \text{ tel que } P_1 = P_2 Q + R 
	\end{align*}
	Vérifiant :
	$\operatorname{deg}(R) = -\infty$ ou $0 \leq \operatorname{deg}(R) < \operatorname{deg}(Q)$.
\end{theorem}

\begin{definition}[Polynôme irréductible]
	Un polynôme $P \in \K[X]$ non constant est dit irréductible, s'il n'existe pas $(P_1, P_2) \in (\K[X])^2$ tel que $P = P_1 P_2$ et $\operatorname{deg}(P_1) < \operatorname{deg}(P_2)$.
\end{definition}

\begin{definition}\label{def:ordre_mult}
	Soient $P$ un polynôme non constant, $\alpha \in \K$ et $m \in \N^*$ $\alpha$ est une racine d'ordre de multiplicité $m$ de $P$ si et seulement si
	\begin{align*}
		(X - \alpha)^m \mid P \text{ et } (X - \alpha)^{m+1} \nmid P.
	\end{align*}
\end{definition}

\begin{theorem}
	Soient $P$ un polynôme non constant, $\alpha \in \K$ et $m \in \N^*$
	\begin{align*}
		(X - \alpha)^m \mid P \iff P(\alpha) = P'(\alpha) = \cdots = P^{m-1}(\alpha) = 0.
	\end{align*}
\end{theorem}

\begin{theorem}[Théorème fondamental de l'algèbre]
	Soit $P(X)$ un polynôme à coefficients complexes de degré $n$. $P(X)$ admet $n$ racines complexes. 
\end{theorem}

\section{Fractions rationnelles}
\begin{definition}[Fraction rationnelle]
	$F(X)$ est appelée fraction rationnelle s'il existe $P(X), Q(X) \in \K[X]$ tels que :
	\begin{align*}
		F(X) = \frac{P(X)}{Q(X)},\ Q(X) \neq 0
	\end{align*}
	avec $\deg(F) = \deg(P) - \deg(Q)$. \\
	Si $\deg(P) > \deg(Q)$, alors il existe $E(X)$ appelée la \emph{partie entière} et des polynômes $R(X), S(X)$ tels que $\deg(R) < \deg(S)$ et : 
	\begin{align*}
		F(X) = E(X) + \frac{R(X)}{S(X)}.
	\end{align*}
\end{definition}

\begin{theorem}[Décomposition en éléments simples]
	Soient $F(X) = \frac{P(X)}{Q(X)}$ et $m, n, \alpha, \beta \in \N^*$.
	\\
	Sur $\R$ :
	\begin{align*}
		Q(X) &= (X - a_1)^{\alpha_1} \ldots (X - a_n)^{\alpha_n} (X^2 + b_1 X + c_1)^{\beta_1} \ldots (X^2 + b_m X + c_m)^{\beta_m} \\ 
            &= \prod_{i = 0}^{n} (X - a_i)^{\alpha_i} \cdot \prod_{k = 0}^{m} (X^2 + b_k X + c_k)^{\beta_k}
	\end{align*}
	Alors $\forall A, B, C \in \R$, $F$ s'écrit de manière unique sous la forme :
	\begin{align*}
		F(X) &= \left[ \left( \frac{A_{11}}{(X - a_1)^1} + \cdots + \frac{A_{1n}}{(X - a_1)^{\alpha_1}} \right) + \cdots + \left( \frac{A_{n1}}{(X - a_n)^1} + \cdots + \frac{A_{nn}}{(X - a_n)^{\alpha_n}}  \right) \right] \\
		+ &\left[ \left( \frac{B_{11}X + C_{11}}{(X^2 + b_1 X + c_1)^1} + \cdots + \frac{B_{1m}X + C_{1m}}{(X^2 + b_1 X + c_1)^{\beta_1}} \right) + \cdots + \right. \\
		&\left. \left( \frac{B_{m1}X + C_{m1}}{(X^2 + b_m X + c_m)^1} + \cdots + \frac{B_{mm}X + C_{mm}}{(X^2 + b_m X + c_m)^{\beta_m}}  \right) \right] \\
        F(X) &=  \sum_{i = 1}^{n} \sum_{j = 1}^{\alpha_i} \frac{A_{ij}}{\left(X - a_i\right)^{j} } + \sum_{k = 1}^{m} \sum_{l = 1}^{\beta_k} \frac{B_{kl}X + C_{kl}}{(X^2 + b_kX + c_k)^{l}}
	\end{align*}
	Sur $\C$ :
	\begin{align*}
		Q(X) &= (X - a_1)^{\alpha_1} \ldots (X - a_n)^{\alpha_n} \\
            &= \prod_{i = 0}^{n} (X - a_i)^{n}
	\end{align*}
	Alors $\forall A, B, C \in \C, \ \forall \alpha \in \N^*$  $F$ s'écrit de manière unique sous la forme :
	\begin{align*}
		F(X) &= \left[ \left( \frac{A_{11}}{(X - a_1)^1} + \cdots + \frac{A_{1n}}{(X - a_1)^{\alpha_1}} \right) + \cdots + \left( \frac{A_{n1}}{(X - a_n)^1} + \cdots + \frac{A_{nn}}{(X - a_n)^{\alpha_n}}  \right) \right] \\
        F(X) &=  \sum_{i = 1}^{n} \sum_{j = 1}^{\alpha_i} \frac{A_{ij}}{\left(X - a_i\right)^{j} }
	\end{align*}    
\end{theorem}

\begin{proof}\cite{math-sup.fr_decomp}
    Sur $\C$ : il est recommandé de consulter la section sur les espaces vectoriels avant (voir \autoref{chap:espaces_vectoriels}).
    Rappelons tout d'abord :
    \[ \C[X] \coloneqq \left\{ \sum_{k = 0}^{n} a_k X^k : n \in \N,\ a_k \in \C \right\} \]
    Notons l'espace des fractions rationnelles :
    \[ \C(X) \coloneqq \left\{ \frac{P(X)}{Q(X)} : P, Q \in \C[X],\ Q \neq 0_{\C[X]} \right\} \]
    La stratégie de la démonstration consiste à montrer une égalité de deux ensembles :
    \[ E = \left\{ \frac{P(X)}{Q(X)} : P \in \C[X],\ \deg(P) < \deg(Q) \right\} \]
    \[ F = \left\{ \sum_{i = 1}^{n} \sum_{j = 1}^{m_i} \frac{a_{ij}}{(X - \alpha_i)^{j}} : a_{ij} \in \C \right\} \]
    \emph{\'Etude de $E$ :} Posons $d \coloneqq \deg(Q)$
    \begin{align*}
        E &= \left\{ \frac{P(X)}{Q(X)} : P \in \C[X],\ \deg(P) < d \right\} \\
          &= \left\{ \frac{a_0 + a_1X + \cdots + a_{d - 1} X^{d - 1}}{Q(X)} : a_0, \ldots, a_{d - 1} \in \C \right\} \\
          &= \left\{ a_0 \frac{1}{Q(X)} + a_1 \frac{X}{Q(X)} + \cdots + a_{d - 1} \frac{X^{d-1}}{Q(X)} : a_0, \ldots, a_{d - 1} \in \C \right\} \\
          &= \operatorname{Vect}\left\{ \frac{1}{Q(X)}, \frac{X}{Q(X)}, \ldots, \frac{X^{d-1}}{Q(X)} \right\}
    \end{align*}
    Notons : $\mathcal{F} \coloneqq \operatorname{Vect}\left\{ \frac{1}{Q(X)}, \frac{X}{Q(X)}, \ldots, \frac{X^{d-1}}{Q(X)} \right\}$
    Nous avons réussi à exprimer $E$ sous la forme d'une famille de vecteurs, nous en déduisons que $E$ est un espace vectoriel admettant $\mathcal{F}$ comme famille génératrice, montrons ensuite que $\mathcal{F}$ est libre :
    Soient $ \lambda_0,\ldots,\lambda_{d-1} \in \C$
    \begin{align*}
        \lambda_0 \frac{1}{Q(X)} + \cdots + \lambda_{d - 1} \frac{X^{d-1}}{Q(X)} &= \frac{\lambda_0 + \lambda_{1} X + \cdots + \lambda_{d - 1} X^{d - 1}}{Q(X)}
    \end{align*}
    \[ \frac{\lambda_0 + \lambda_{1} X + \cdots + \lambda_{d - 1} X^{d - 1}}{Q(X)} = 0_{\C(X)} \iff \lambda_0 + \lambda_{1} X + \cdots + \lambda_{d - 1} X^{d - 1} = 0_{\C[X]} \]
    Un polynôme est nul si tous ses coefficients sont nuls, ainsi :
    \[ \lambda_0 = \cdots = \lambda_{d-1} = 0 \]
    Ainsi $\mathcal{F}$ est libre. $\mathcal{F}$ est libre et génératrice, elle forme donc une base de $E$ et ainsi : 
    \[ \dim(E) = d \]
    \emph{\'Etude de F :}
    \[ F = \left\{ \sum_{i = 1}^{n} \sum_{j = 1}^{m_i} \frac{a_{ij}}{(X - \alpha_i)^{j}} : a_{ij} \in \C \right\} \]
    On peut ainsi écrire $F$ ainsi : 
    \begin{align*}
        F &= \left\{ \frac{a_{11}}{(X - \alpha_1)} + \cdots + \frac{a_{1m_1}}{(X - a_1)^{m_1}} + \cdots + \frac{a_{n1}}{(X - \alpha_n)} + \cdots + \frac{a_{nm_n}}{(X - \alpha_n)^{m_n}} : a_{ij} \in \C \right\} \\
        &= \operatorname{Vect} \left\{ \frac{1}{(X - \alpha_1)}, \ldots, \frac{1}{(X - \alpha_1)^{m_1}}, \ldots, \frac{1}{(X - \alpha_n)}, \ldots, \frac{1}{(X - \alpha_n)^{m_n}}  \right\}
    \end{align*}
    Notons $\mathcal{F}_2 \coloneqq \operatorname{Vect} \left\{ \frac{1}{(X - \alpha_1)}, \ldots, \frac{1}{(X - \alpha_1)^{m_1}}, \ldots, \frac{1}{(X - \alpha_n)}, \ldots, \frac{1}{(X - \alpha_n)^{m_n}}  \right\}$. Nous en déduisons que $F$ est un espace vectoriel admettant $\mathcal{F}_2$ comme famille génératrice, montrons que $\mathcal{F}_2$ est libre : Soient $a_{ij} \in \C$ tels que :
    \begin{align*}
        \begin{split}
            a_{11} \frac{1}{(X - \alpha_1)} + a_{12} \frac{1}{(X - \alpha_1)^2} + \cdots + a_{1(m_1 -1)} \frac{1}{(X - \alpha_1)^{m_1 - 1}} + \cdots \\
            + a_{n1} \frac{1}{(X - \alpha_n)} + a_{n2} \frac{1}{(X - \alpha_n)^2} + \cdots + a_{nm_n} \frac{1}{(X - \alpha_n)^{m_n}} = 0_{\C(X)}
        \end{split}
    \end{align*}
    En multipliant l'équation par $(X - \alpha_1)^{m_1}$ :
    \begin{align*}
        \begin{split}
            a_{11} (X - \alpha_1)^{m_1 - 1} + a_{12} (X - \alpha_1)^{m_1 - 2} + \cdots + a_{1m_1} + \\
            (X - \alpha_1)^{m_1}
            \left( 
            a_{21} \frac{1}{(X - \alpha_2)} + \cdots + a_{n1} \frac{1}{(X - \alpha_n)} + \cdots + a_{nm_n} \frac{1}{(X - \alpha_n)^{m_n}}
            \right)
            = 0_{\C(X)}
         \end{split}
    \end{align*}
    En posant $X = \alpha_1$, on trouve que $a_{1m_1} = 0$.
    En remplaçant $a_{1m_1}$ par sa valeur dans l'équation initiale, celle-ci devient : 
    \begin{align*}
        \begin{split}
            a_{11} \frac{1}{(X - \alpha_1)} + a_{12} \frac{1}{(X - \alpha_1)^2} + \cdots + a_{1(m_1 - 1)} \frac{1}{(X - \alpha_1)^{m_1 - 1}} + \cdots \\
            + a_{n1} \frac{1}{(X - \alpha_n)} + a_{n2} \frac{1}{(X - \alpha_n)^2} + \cdots + a_{nm_n} \frac{1}{(X - \alpha_n)^{m_n}} = 0_{\C(X)}
        \end{split}
    \end{align*}
    En multipliant l'équation par $(X - \alpha_1)^{m_1 - 1}$ et en posant $X = \alpha_1$, on trouve que $a_{1(m_1 - 1)} = 0$. On procède de la même manière jusqu'à prouver $a_{11} = \cdots = a_{1m_1} = 0$.
    On continue ainsi de suite pour montrer que tous les coefficients sont nuls et donc que $\mathcal{F}_2$ est libre. $\mathcal{F}_2$ est libre et génératrice, elle forme donc une base de $F$.
    \[ \dim(F) = m_1 + \cdots + m_n \]
    Mais aussi : 
    \[ \deg(Q) = \deg \left( \prod_{i=1}^{n} (X - \alpha_i)^{m_i} \right) = m_1 + \cdots + m_n \]
    Ainsi : \[ \dim(F) = \deg(Q) = d \]
    \emph{Inclusion de $F$ dans $E$ :}
    \\
    Un élément de $F$ est de la forme suivante : 
    \[ \frac{a_{11}}{(X - \alpha_1)} + \cdots + \frac{a_{1m_1}}{(X - \alpha_1)^{m_1}} + \cdots + \frac{a_{n1}}{(X - \alpha_n)} + \cdots + \frac{a_{nm_n}}{(X - \alpha_n)^{m_n}} \]
    En mettant toutes les fractions sur le même dénominateur, on obtient :
    \[ \frac{a_{11}(X - \alpha_1)^{m_1 - 1} \prod_{i = 2}^{n}(X - \alpha_i)^{m_i} + \cdots + a_{1m_1} \prod_{i = 2}^{n}(X - \alpha_i)^{m_i} + \cdots + a_{nm_n} \prod_{i = 1}^{n - 1}(X - \alpha_i)^{m_i}}{Q(X)} \]
    Le degré du numérateur est strictement inférieur à celui de $Q$. Donc c'est un élément de $E$. Ainsi :
    \[ F \subset E \]
    Nous avons montré que $\dim(E) = \dim(F)$ et que $F \subset E$, ainsi $E = F$.
    Ce qui revient à dire que toute fraction rationnelle 
    \begin{align*}
        \frac{P(X)}{Q(X)},\ \deg(P) < \deg(Q)
    \end{align*}
    s'exprime sous la forme 
    \begin{align*}
        \sum_{i=1}^n \sum_{j = 1}^{m_i} \frac{a_{ij}}{(X - \alpha_i)^j}
    \end{align*}
    L'unicité de la décomposition découle du fait que tout élément d'un espace vectoriel s'exprime par une unique combinaison linéaire des vecteurs d'une de ses bases.
\end{proof}
\chapter{Systèmes linéaires et matrices}
\def\arraystretch{1}

\section{Définitions et opérations élémentaires}
\begin{definition}[Matrice]
	Soient $m, n \in \N^*$ et $a_{1,1}, \ldots, a_{m,n} \in \K$. 
	\\
    Une \textbf{matrice} est un tableau de données appartenant à $\mathcal{M}_{m,n} (\K)$ :
    \begin{align*}
        \begin{pmatrix}
            a_{1,1} & a_{1,2} & \cdots & a_{1,n} \\
            a_{2,1} & a_{2,2} & \cdots & a_{2,n} \\
            \vdots & \vdots & \ddots & \vdots \\
            a_{m,1} & a_{m,2} & \cdots & a_{m,n}
        \end{pmatrix}
        = 
        (a_{i,j})_{\substack{1 \leq i \leq m \\ 1 \leq j \leq n}}
    \end{align*}
\end{definition}

\begin{definition}[Système linéaire]
    Soient $m, n \in \N^*$, $x_1, \ldots, x_n \in \K$ $n$ inconnues, $a_{1, 1}, \ldots, a_{m,n} \in \K$ et $b_1, \ldots, b_n \in \K$.
    \\
    Un \textbf{système linéaire} est décrit par :
    \begin{align*}
        \begin{cases}
            a_{1,1} \cdot x_1 + a_{1,2} \cdot x_2 + \cdots + a_{1,n} \cdot x_n = b_1 \\
            a_{2,1} \cdot x_1 + a_{2,2} \cdot x_2 + \cdots + a_{2,n} \cdot x_n = b_2 \\
            \vdots \\
            a_{m,1} \cdot x_1 + a_{m,2} \cdot x_2 + \cdots + a_{m,n} \cdot x_n = b_m
        \end{cases}
    \end{align*}        
    Sa matrice associée est : 
    \begin{align*}
        \left(
        \begin{matrix}    
            a_{1,1} & a_{1,2} & \cdots & a_{1,n} \\
            a_{2,1} & a_{2,2} & \cdots & a_{2,n} \\
            \vdots & \vdots & \ddots & \vdots \\
            a_{m,1} & a_{m,2} & \cdots & a_{m,n}
        \end{matrix}
        \
        \middle|
        \
        \begin{matrix}
            b_1 \\
            b_2 \\
            \vdots \\
            b_m
        \end{matrix}
        \right)
    \end{align*}
\end{definition}

\begin{definition}[Opérations élémentaires]
	Soient $i, j$ tels que $i \neq j$ des numéros de ligne et $\lambda \in \K^*$.
	\begin{enumerate}
		\item $L_i \leftarrow L_i + \lambda L_j$ 
		\item $L_i \leftrightarrow L_j$
		\item $L_i \leftarrow \lambda \cdot L_i$
	\end{enumerate}
	Opérations analogues sur les colonnes.
\end{definition}

\begin{definition}
	\begin{enumerate}
		\item Une matrice est \textbf{échelonnée} en lignes si et seulement si le nombre de zéros commençant une ligne croît strictement ligne par ligne jusqu'à ce qu'il ne reste plus que des zéros.
		\item Une matrice est dite dite \textbf{échelonnée réduite} en lignes si les pivots valent 1 et si les autres coefficients dans les colonnes des pivots sont nuls.
	\end{enumerate}
\end{definition}

\begin{example}
	$
	\begin{pmatrix}
		2 & 1 & 3 \\
		0 & 1 & 1 \\
		0 & 0 & 0
	\end{pmatrix}
	$
	est une matrice échelonnée.
	\\
	$
	\begin{pmatrix}
		1 & 1 & 0 \\
		0 & 0 & 1 \\
		0 & 0 & 0
	\end{pmatrix}
	$
	est une matrice échelonnée réduite.
\end{example}

\par Décrivons maintenant l'algorithme du pivot de Gauss utilisé pour résoudre des systèmes. On expliquera la méthode sur les lignes mais les principes sont analogues pour les colonnes.
\par \noindent La méthode consiste à appliquer les opérations élémentaires sur le système ou la matrice afin d'en faire un système ou une matrice échelonnée réduite.

Considérons le système suivant :
\[
\begin{cases}
	a_{1,1} x_1 + a_{1,2} x_2 + \cdots + a_{1,n} x_n = b_1 \\
	a_{2,1} x_1 + a_{2,2} x_2 + \cdots + a_{2,n} x_n = b_2 \\
	\vdots \\
	a_{m,1} x_1 + a_{m,2} x_2 + \cdots + a_{m,n} x_n = b_n
\end{cases}
\]

\par Le but est d'annuler tous les coefficients $a_{2,1}, a_{3,1}, \ldots, a_{m,1}$ à partir de $a_{1,1}$ puis d'annuler les coefficients $a_{3,2}, a_{4,2}, \ldots, a_{m,2}$ à partir de $a_{2,2}$ et ainsi de suite jusqu'à que $a_{m,n} = 1$.
\\
\par Pour annuler le coefficient $a_{i,j}$ à partir du pivot $a_{k,j}$ avec $k < i$, on applique l'opération élémentaire suivante :
\[ L_i \leftarrow L_i - \frac{a_{i,j}}{a_{k,j}} L_k \]

\`A la fin de toutes ces opérations, on obtient un système linéaire de la forme :

\[
\left\{
\begin{array}{rrrrrrrrr}
	a_{1,1} x_1 & + & a_{1,2} x_2  & + & \cdots & + & a_{1,n} x_n  & = & b_1  \\
	            &   & a_{2,2}' x_2 & + & \cdots & + & a_{2,n}' x_n & = & b_2' \\
	            &   &              & \vdots  &  &   &              &   &      \\
	            &   &              &   &        &   & x_n          & = & b_n'
\end{array}
\right.
\]

On va ensuite chercher à remonter dans le système en appliquant les opérations élémentaires avec le même principe que lors de la descente. Pour ensuite obtenir un système de la forme 

\[
\left\{
\begin{array}{ccccccccc}
	x_1 &  &     &        &  &  &     & = & b_1'  \\
	    &  & x_2 &        &  &  &     & = & b_2'' \\
	    &  &     & \ddots &  &  &     &   &       \\
	    &  &     &        &  &  & x_n & = & b_n'
\end{array}
\right.
\]

\begin{example}
	Utilisons l'algorithme du pivot de Gauss pour résoudre le système suivant.
	\[
	\systeme{
	-3x -2y + 2z = 0,
	-2x -y + z = 2,
	7x + 5y - 4z = -3
	} 
	\]
	\begin{align*}
		\systeme{
			-3x -2y + 2z = 0,
			-2x -y + z = 2,
			7x + 5y - 4z = -3
		} 
		&\overset{\substack{L_2 \leftarrow L_2 - \frac{-2}{-3} L_1 \\ L_3 \leftarrow L_3 - \frac{7}{-3} L_1 }}{\iff}
		\systeme{
			-3x -2y + 2z = 0, 
			\frac{1}{3}y - \frac{1}{3}z = 2,
			\frac{1}{3}y + \frac{2}{3}z = -3
		} 
		\\
		&\overset{\substack{L_2 \leftarrow 3 L_2 \\ L_3 \leftarrow 3 L_3}}{\iff} 
		\systeme{
			-3x -2y + 2z = 0,
			y - z = 6,
			y + 2z = -9
		}
		\\
		&\overset{L_3 \leftarrow L_3 - L_2}{\iff}
		\systeme{
			-3x -2y + 2z = 0,
			y - z = 6,
			3z = -15
		}
		\\
		&\overset{L_3 \leftarrow \frac{1}{3} L_3}{\iff}
		\systeme{
			-3x -2y + 2z = 0,
			y - z = 6,
			z = -5
		}
		\\
		&\overset{L_2 \leftarrow L_2 + L_3}{\iff}
		\systeme{
			-3x - 2y + 2z = 0,
			y = 1,
			z = -5
		}
		\\
		&\overset{L_1 \leftarrow L_1 + 2 L_2 - 2 L_3}{\iff}
		\systeme{
			-3x = 12,
			y = 1,
			z = -5
		}
		\\
		&\overset{-\frac{1}{3}L_1}{\iff}
		\systeme{
		x = -4,
		y = 1,
		z = -5
		}
	\end{align*}
\end{example}

\begin{definition}[Rang d'une matrice]
	Le rang d'une matrice $A$ est son nombre de lignes non nulles après échelonnage. Il est noté $\rg(A)$.
\end{definition}

\begin{theorem}
    Soient $\mathcal{S}$ un système linéaire de $m$ lignes et $n$ inconnues, $A \in \mathcal{M}_{m,n}(\K)$ et $B \in \mathcal{M}_{m,1}(\K)$ telles que $A|B$ forme la matrice associée à $\mathcal{S}$. $\mathcal{S}$ est solvable si et seulement si :
    \[ \rg(A) = \rg(A|B) \]
\end{theorem}

\section{Opérations sur les matrices}
\begin{definition}[Opérations sur les matrices]
	Soient $m, n, p, q \in \N^*$, $A \in \mathcal{M}_{m, n}(\K)$, $B \in \mathcal{M}_{p, q}(\K)$ et $\lambda  \in \K$ tels que :
	\begin{align*}
		A &=
		\begin{pmatrix}
			a_{1,1} & \cdots & a_{1,n} \\
			\vdots & \ddots & \vdots \\
			a_{m,1} & \cdots & a_{m,n}
		\end{pmatrix}
		&
		B &= 
		\begin{pmatrix}
			b_{1,1} & \cdots & b_{1,q} \\
			\vdots & \ddots & \vdots \\
			b_{p,1} & \cdots & b_{p,q}
		\end{pmatrix}
	\end{align*}
	\begin{enumerate}
		\item Si $m = p$ et $n = q$ alors on peut définir l'addition entre $A$ et $ B$ et la multiplication par $\lambda$.
		\begin{align*}
			A + \lambda B &= 
			\begin{pmatrix}
				(a_{1,1} + \lambda b_{1,1}) & \cdots & (a_{1,n} + \lambda b_{1,q}) \\
				\vdots & \ddots & \vdots \\
				(a_{m,1} + \lambda b_{p,1}) & \cdots & (a_{m,n} + \lambda b_{p,q})
			\end{pmatrix}
		\end{align*}
		\item Si $n = p$ alors on peut définir la multiplication entre $A$ et $B$.\\
		Soit $C = AB$. Chaque coefficient $c_{i,j}$ de $C$ est défini par :
		\begin{align*}
			c_{i,j} = \sum_{k=1}^{n} a_{i,k} b_{k,j}
		\end{align*}
        Autrement dit, le coefficient $c_{i,j}$ pour $C = AB$ est donné par :
        \begin{align*}
            &\begin{pmatrix}
                b_{1,1} & \cdots & \color{red}b_{1,j} & \cdots & b_{1,q} \\
                b_{2,1} & \cdots & \color{green}b_{2,j} & \cdots & b_{2,q} \\
                \vdots & \cdots & \vdots & \cdots & \vdots \\
                b_{p,1} & \cdots & \color{blue}b_{pj} & \cdots & b_{p,q}
            \end{pmatrix}
            \\
            \begin{pmatrix}
                a_{1,1} & a_{1,2} & \cdots & a_{1,n} \\ 
                \vdots & \vdots & \vdots & \vdots \\
                \color{red}a_{i,1} & \color{green}a_{i,2} & \cdots & \color{blue}a_{i,n} \\ 
                \vdots & \vdots & \vdots & \vdots \\
                a_{m,1} & a_{m,2} & \cdots & a_{m,n}
            \end{pmatrix}
            &\begin{pmatrix}
                c_{1,1} & c_{1,2} & \cdots & \cdots & c_{1,q} \\
                \vdots & \vdots & \vdots & \vdots & \vdots \\
                \cdots & \cdots & \boxed{c_{i,j}} & \cdots & c_{i,q} \\
                \vdots & \vdots & \vdots & \vdots & \vdots \\
                c_{m,1} & \cdots & \cdots & \cdots & c_{m,q}
            \end{pmatrix}
        \end{align*}
        \[ \boxed{c_{i,j}} = \color{red}a_{i,1} \cdot b_{1,j} \color{black}+ \color{green} a_{i,2} \cdot b_{2,j} \color{black}+ \color{blue}\cdots + a_{i,n} \cdot b_{p,j} \]
		Attention, la multiplication n'est pas commutative $(AB \neq BA)$.
	\end{enumerate}
\end{definition}

\begin{example}
	Soient
	$A =
	\begin{pmatrix}
		1 & 2 \\
		3 & 4
	\end{pmatrix}
	$
	et 
	$B =
	\begin{pmatrix}
		5 & 6 \\
		7 & 8
	\end{pmatrix}
	$
	\begin{align*}
		A + B &= 
		\begin{pmatrix}
			1 + 5 & 2 + 6 \\
			3 + 7 & 4 + 8
		\end{pmatrix}
		\\
		&= 
		\begin{pmatrix}
			6 & 8 \\
			10 & 12
		\end{pmatrix}
	\end{align*}
	\begin{align*}
		&\begin{pmatrix}
			5 & 6 \\
			7 & 8
		\end{pmatrix}
		\\
		\begin{pmatrix}
			1 & 2 \\
			3 & 4
		\end{pmatrix}
		&
		\begin{pmatrix}
			1 \cdot 5 + 2 \cdot 7 & 1 \cdot 6 + 2 \cdot 8 \\
			3 \cdot 5 + 4 \cdot 7 & 3 \cdot 6 + 4 \cdot 8
		\end{pmatrix} = AB
		\\
		AB = 
		&\begin{pmatrix}
			19 & 22 \\
			43 & 50
		\end{pmatrix}
	\end{align*}
\end{example}

\begin{proposition}
	Soient $m, n, k, l \in \N^*$.
	\begin{enumerate}
		\item Soient $A, B, C \in \mathcal{M}_{m,n} (\K)$.
		\[A + B = B + A\]
		\[A + (B + C) = (A + B) + C\]
		\item Soient $A \in \mathcal{M}_{m, k}(\K),\ B \in \mathcal{M}_{k, l}(\K),\ C \in \mathcal{M}_{l, n}(\K)$.
		\[A \cdot (B \cdot C) = (A \cdot B) \cdot C\]
		\item Soient $A \in \mathcal{M}_{m, k}(\K),\ B, C \in \mathcal{M}_{k, n}(\K),\ \lambda \in \K$.
		\[A \cdot (B + \lambda C) = A \cdot B + \lambda \cdot A \cdot C\]
		\[(A + \lambda B) \cdot C = A \cdot C + \lambda \cdot B \cdot C\]
	\end{enumerate}
\end{proposition}

\begin{proposition}[Binôme de Newton]
	Soient $n \in \N^*$, $A, B \in \mathcal{M}_n(\K)$ telles que $AB = BA$ et $p \in \N$.
	\[ (A + B)^p = \sum_{k = 0}^{p} \binom{p}{k} A^k B^{p-k} \]
\end{proposition}

\begin{definition}[Matrice nulle]
	La matrice nulle est la matrice dont tous les coefficients sont $0$. On la note $0_{m,n}$, $m$ étant le nombre de lignes, $n$ le nombre de colonnes. 
\end{definition}

\begin{definition}[Symbole de Kronecker]
	On définit le symbole de Kronecker ainsi pour $i, j \in \N^*$ :
	\[ \delta_{i,j} = 
	\begin{cases}
		1 \text{ si } i = j \\
		0 \text{ si } i \neq j
	\end{cases}
	\]
\end{definition}

\begin{definition}[Matrice identité]
	La matrice identité est la matrice dont tous les coefficients sont $0$ à l'exception de ceux de la diagonale principale à $1$. On la note $I_n$, $n$ étant le nombre de lignes.
	\begin{align*}
		I_n = 
		\begin{pmatrix}
			1 & 0 & \cdots & 0 \\
			0 & 1 & \cdots & 0 \\
			\vdots & \ddots & \ddots & \vdots \\
			0 & \cdots & 0 & 1
		\end{pmatrix}
	\end{align*}
	On peut également dire que :
	\[ I_n = (\delta_{i,j})_{\substack{1 \leq i \leq n \\ 1 \leq j \leq n}} \]
\end{definition}

\begin{remark}
    La matrice identité est parfois notée : $\mathds{1}_n$.
\end{remark}

\begin{lemma}
	Soient $m,n \in \N^*$.
	\begin{multicols}{2}
	    \begin{enumerate}
    		\item $\forall A \in \mathcal{M}_{m, n}(\K),\ A + 0_{m,n} = A$
    		\item $\forall A \in \mathcal{M}_{n}(\K),\ A \cdot I_n = I_n \cdot A = A$
    	\end{enumerate}
	\end{multicols}
\end{lemma}

\begin{definition}[Transposée]
	Soient $m, n \in \N^*$, $A \in \mathcal{M}_{m, n}(\K)$ de coefficients $a_{1,1}, \ldots, a_{m,n} \in \K$.
	\\
	La transposée de $A$ est notée $A^T$. $A^T \in \mathcal{M}_{n,m}(\K)$ et pour $a^T_{1,1}, \ldots, a^T_{n,m} \in \K$ ses coefficients :
	\[ (a^T_{j,i})_{\substack{1 \leq j \leq n \\ 1 \leq i \leq m}} = (a_{i,j})_{\substack{1 \leq i \leq m \\ 1 \leq j \leq n}} \]
\end{definition}

\begin{example}
	Soit $A$ une matrice telle que :
	\[ 
	A = 
	\begin{pmatrix}
	1 & 2 & 3 \\
	4 & 5 & 6
	\end{pmatrix}
	\]
	\[
	A^T =
	\begin{pmatrix}
	1 & 4 \\
	2 & 5 \\
	3 & 6
	\end{pmatrix}
	\]
\end{example}

\begin{definition}
	Soient $n \in \N^*$ et $M \in \mathcal{M}_n(\K)$.
	    \begin{enumerate}
    		\item $M$ est symétrique $\iff M^T = M$
    		\item $M$ est anti-symétrique $\iff M^T = -M$
    	\end{enumerate}
\end{definition}

\begin{definition}
	Soit $A \in \mathcal{M}_{m,n}(\C)$.
	\\
	On définit $A^* = (\overline{A})^T = (\overline{A^T})$ la \textbf{matrice adjointe / transconjuguée}.
	Autrement dit : on prend la transposée de $A$ et on prend les conjugués des coefficients.
\end{definition}

\begin{example}
	Soit $A$ telle que :
	\[ A =
	\begin{pmatrix}
		1 + i & 3 - i \\
		2 & -2 \\
		i & 5i
	\end{pmatrix}
	\]
	On a $A^T$ :
	\[ 
	A^T = 
	\begin{pmatrix}
		1 + i  & 2 & i \\
		3 - i & -2 & 5i
	\end{pmatrix}
	\]
	et finalement $A^*$ :
	\[
	A^* =
	\begin{pmatrix}
		1 - i & 2 & -i \\
		3 + i & -2 & -5i
	\end{pmatrix}
	\]
\end{example}

Pour calculer le déterminant d'une matrice carrée noté $\det(A)$ ou s'il n'y a pas d'ambiguïté $\abs{A}$, il existe plusieurs méthodes.
\\
Soit $A$ une matrice carrée de taille $n \times n$ de coefficients $a_{1,1},\ldots,a_{n,n} \in \K$. 
\begin{enumerate}
    \item $n = 2$. 
    \[ 
    A = 
    \begin{pmatrix}
        a_{1,1} & a_{1,2} \\
        a_{2,1} & a_{2,2}
    \end{pmatrix}
    \]
    \[ \det(A) = a_{1,1} a_{2,2} - a_{2,1} a_{1,2} \]
    \item $n = 3$
    \[
    A = 
    \begin{pmatrix}
        a_{1,1} & a_{1,2} & a_{1,3} \\
        a_{2,1} & a_{2,2} & a_{2,3} \\
        a_{3,1} & a_{3,2} & a_{3,3}
    \end{pmatrix}
    \]
    \[ \det(A) = (a_{1,1} a_{2,2} a_{3,3} + a_{1,2} a_{2,3} a_{3,1} + a_{1,3} a_{2,1} a_{3,2} ) - (a_{1,3} a_{2,2} a_{3,1} + a_{1,2} a_{2,1} a_{3,3} + a_{1,1} a_{2,3} a_{3,2}) \]
    \item $n \geq 3$ \cite{wikipedia_determinant}.
    On peut calculer le déterminant d'une matrice en calculant à l'aide des déterminants des matrices de taille $n-1$. Pour se faire on développe en lignes ou en colonnes.
    Notons $A_{i,j}$ ma matrice obtenue en enlevant à $A$ sa $i$-ème ligne et sa $j$-ème colonne.
    \begin{align*}
       	A_{i,j} = 
       	\begin{pmatrix}
       		a_{1,1} & \cdots & a_{1, j-1} & a_{1, j+1} & \cdots & a_{1,n} \\
       		\vdots & & \vdots & \vdots & & \vdots \\
       		a_{i-1, 1} & \ldots & a_{i-1, j-1} & a_{i-1, j+1} & \cdots & a_{i-1,n} \\
       		a_{i+1, 1} & \cdots & a_{i+1, j-1} & a_{i+1, j+1} & \cdots & a_{i+1, n} \\
       		\vdots & & \vdots & \vdots & & \vdots \\
       		a_{n,1} & \cdots & a_{n,j-1} & a_{n,j+1} & \cdots & a_{n,n}
       	\end{pmatrix}
    \end{align*}
    On peut alors développer le calcul du déterminant  suivant une ligne ou une colonne : 
    \begin{enumerate}
       	\item Suivant la ligne $i$ : 
       	\[ \det(A) = \sum_{j=1}^n a_{i,j} (-1)^{i+j} \det(A_{i,j}) \]
       	\item Suivant la colonne $j$ : 
       	\[ \det(A) = \sum_{i=1}^n a_{i,j} (-1)^{i+j} \det(A_{i,j}) \]
    \end{enumerate}
\end{enumerate}

\begin{example}
	Calculons le déterminant de la matrice 
	\begin{align*}
		A = 
		\begin{pmatrix}
		2 & 1 & -2 \\
		1 & 1 & -2 \\
		1 & 2 & -1
		\end{pmatrix}
	\end{align*}
	en développant suivant la première ligne.
	\begin{align*}
		\det(A) &= 2 (-1)^{1+1} 
		\begin{vmatrix}
			1 & -2 \\
			2 & -1	
		\end{vmatrix}			
		+ 1 (-1)^{1 + 2}
		\begin{vmatrix}
			1 & -2 \\
			1 & -1
		\end{vmatrix}
		+ (-2) (-1)^{1 + 3}
		\begin{vmatrix}
			1 & 1 \\
			1 & 2
		\end{vmatrix}
		\\
		&=
		2[1 \cdot (-1) - (-2) \cdot 2] -1[1 \cdot (-1) - (-2) \cdot 1] -2 [1 \cdot 2 - 1 \cdot 1]
		\\
		&= 6 - 1 - 2 \\
		&= 3
	\end{align*}
\end{example}

\begin{proposition}
	Soient $n \in \N^*$, $A, B \in \mathcal{M}_n(\K)$ et $\lambda \in \K$.
	\begin{multicols}{2}
	    \begin{enumerate}
		\item $\det(A \cdot B) = \det(A) \cdot \det(B)$
		\item $\det(\lambda A) = \lambda^n \det(A)$
		\item $\det(A^T) = \det(A)$
		\item $\det(A^{-1}) = \frac{1}{\det(A)}$ si $\det(A) \neq 0$
	\end{enumerate}
	\end{multicols}
\end{proposition}

\begin{definition}[Matrice inversible]
	Soient $n \in \N^*$ et $A \in \mathcal{M}_n(\K)$.
	\\
	$A$ est inversible si et seulement si existe une unique matrice $B = A^{-1}$ telle que 
	\[AB = BA = I_n.\]
\end{definition}

\begin{proposition}
	Une matrice est inversible si et seulement si son déterminant est non nul.
\end{proposition}

\begin{definition}[Comatrice]
	La comatrice $\com(A)$ d'une matrice $A \in \mathcal{M}_n(\K)$ est une matrice est définie ainsi :
	\[ (\com(A))_{i,j} = (-1)^{i+j} \det(A_{i,j}) \]
\end{definition}

\par Lorsque nous devons calculer l'inverse d'une matrice, plusieurs cas sont possibles.
Si on a une matrice
$A =
\begin{pmatrix}
	a & b \\
	c & d
\end{pmatrix}
$.
\[
A^{-1} = \frac{1}{\det(A)} 
\begin{pmatrix}
	d & -b \\
	-c & a
\end{pmatrix}
\]
Une méthode générale pour calculer l'inverse d'une matrice $A$ est d'utiliser l'algorithme de Gauss-Jordan.
On part de la matrice 
\[ A|I_n \]
et on applique les opérations élémentaires pour avoir une matrice de la forme
\[I_n|B\]
$B = A^{-1}$ est l'inverse de $A$.
\\
Une autre méthode générale est d'utiliser la comatrice de $A$.
\[ A^{-1} = \frac{1}{\det(A)} \left( \com(A) \right)^T \]

\begin{example}
	Soit 
	$
	A =
	\begin{pmatrix}
		1 & 2 & 1 \\
		-1 & -1 & 2 \\
		-2 & -2 & 6
	\end{pmatrix}
	$.
	\\
	Posons $A | I_n$ :
	\begin{align*}
		A | I_n = 
		&\left( 
		\begin{matrix}
			1 & 2 & 1 \\
			-1 & -1 & 2 \\
			-2 & -2 & 6
		\end{matrix}
		\ \middle| \
		\begin{matrix}
			1 & 0 & 0 \\
			0 & 1 & 0 \\
			0 & 0 & 1
		\end{matrix}
		\right) \\
		\overset{\substack{L_2 \leftarrow L_2 + L_1 \\ L_3 \leftarrow L_3 + 2 L_1}}{\iff}
		&\left(
		\begin{matrix}
			1 & 2 & 1 \\
			0 & 1 & 3 \\
			0 & 2 & 8
		\end{matrix}
		\ \middle| \ 
		\begin{matrix}
			1 & 0 & 0 \\
			1 & 1 & 0 \\
			2 & 0 & 1
		\end{matrix}
		\right) \\
		\overset{\substack{L_3 \leftarrow L_3 - 2 L_1}}{\iff} 
		&\left(
		\begin{matrix}
			1 & 2 & 1 \\
			0 & 1 & 3 \\
			0 & 0 & 2
		\end{matrix}
		\ \middle| \ 
		\begin{matrix}
			1 & 0 & 0 \\
			1 & 1 & 0 \\
			0 & -2 & 1
		\end{matrix}
		\right) \\
		\overset{\substack{L_3 \leftarrow \frac{1}{2} L_3}}{\iff} 
		&\left(
		\begin{matrix}
			1 & 2 & 1 \\
			0 & 1 & 3 \\
			0 & 0 & 1
		\end{matrix}
		\ \middle| \
		\begin{matrix}
			1 & 0 & 0 \\
			1 & 1 & 0 \\
			0 & -1 & \frac{1}{2}
		\end{matrix}
		\right) \\
		\overset{\substack{L_2 \leftarrow L_2 - 3 L_3}}{\iff}
		&\left(
		\begin{matrix}
			1 & 2 & 1 \\
			0 & 1 & 0 \\
			0 & 0 & 1
		\end{matrix}
		\ \middle| \
		\begin{matrix}
			1 & 0 & 0 \\
			1 & 4 & -\frac{3}{2} \\
			0 & -1 & \frac{1}{2}
		\end{matrix}
		\right) \\
		\overset{\substack{L_1 \leftarrow L_1 - 2 L_2 - L_3}}{\iff}
		&\left(
		\begin{matrix}
			1 & 0 & 0 \\
			0 & 1 & 0 \\
			0 & 0 & 1
		\end{matrix}
		\ \middle| \
		\begin{matrix}
			-1 & -7 & \frac{5}{2} \\
			1 & 4 & -\frac{3}{2} \\
			0 & -1 & \frac{1}{2}
		\end{matrix}
		\right)
	\end{align*}
	Ainsi on a :
	\[
	A^{-1} = 
	\begin{pmatrix}
		-1 & -7 & \frac{5}{2} \\
		1 & 4 & -\frac{3}{2} \\
		0 & -1 & \frac{1}{2}
	\end{pmatrix}
	\]
\end{example}

\begin{definition}[Trace d'une matrice]
	Soient $A \in \mathcal{M}_n(\K)$ et $a_{1,1}, \ldots, a_{n,n}$ les coefficients de A. 
	\\
	La trace de $A$, notée $\tr(A)$ est définie par l'application suivante.
	\begin{center}
		$
		\appli{\tr}{\mathcal{M}_n(\K)}{\K}{A}{\sum_{i=1}^n a_{i,i}}
		$
	\end{center}
\end{definition}

\begin{example}
	Soit 
	$
	A = 
	\begin{pmatrix}
		1 & 2 & 3 \\
		4 & 5 & 6 \\
		7 & 8 & 9
	\end{pmatrix}
	$,
	on a $\tr(A) = 1 + 5 + 9 = 15$.
\end{example}

\begin{lemma}
	\begin{multicols}{2}
	    \begin{enumerate}
    		\item $\tr(A^T) = \tr(A)$
    		\item $\tr(A B) = \tr(B A)$
    	\end{enumerate}
	\end{multicols}
\end{lemma}

\begin{proof}
	\leavevmode
    \begin{enumerate}
        \item Par application directe de la définition de la trace et de la transposée.
        \item Les coefficients $a_{1,1}, \ldots, a_{n,n}$ et $b_{1,1}, \ldots, b_{n,n}$ sont respectivement les coefficients des matrices $A$ et $B$. Les coefficients $(ab)_{1,1}, \ldots, (ab)_{n,n}$ et $(ba)_{1,1}, \ldots, (ba)_{n,n}$ sont respectivement ceux des matrices $AB$ et $BA$.
        \begin{align*}
        	\tr(AB) &= \sum_{i=1}^n (ab)_{i,i} \\
        			&= \sum_{i=1}^n \left( \sum_{k=1}^n a_{i,k} b_{k,i} \right) \\
        			&= \sum_{i=1}^n (a_{i,1} b_{1,i} + \cdots + a_{i,n} b_{n,i}) \\
        			&= (a_{1,1} b_{1,1} + \cdots + a_{1,n} b_{n,1}) + \cdots + (a_{n,1} b_{1,n} + \cdots + a_{n,n} b_{n,n}) \\
        			&= (b_{1,1}a_{1,1} + \cdots + b_{1,n}a_{n,1}) + \cdots + (b_{n,1} a_{1,n} + \cdots + b_{n,n} a_{n,n}) \\
        			&= \sum_{i=1}^n \left( b_{i,1} a_{1,i} + \cdots + b_{i,n} a_{n,i} \right) \\
        			&= \sum_{i=1}^n \left( \sum_{k=1}^n b_{i,k} a_{k,i} \right) \\
        			&= \sum_{i=1}^n (ba)_{i,i} \\
        			&= \tr(BA)
        \end{align*}
    \end{enumerate}
\end{proof}

\begin{definition}[Produit scalaire dans $\R^n$] 
	Soient $\vec{x} = \begin{pmatrix} x_1 \\ \vdots \\ x_n \end{pmatrix}$ et $\vec{y} = \begin{pmatrix} y_1 \\ \vdots \\ y_n \end{pmatrix}$ deux vecteurs de $\R^n$.
	\\
	On définit le produit scalaire ainsi :
	\begin{center}
		$
		\appli{\langle \cdot | \cdot \rangle}{\R^n \times \R^n}{\R}{(\vec{x}, \vec{y})}{\sum_{i = 1}^n x_i y_i}
		$
	\end{center}
\end{definition}

\begin{definition}[Produit vectoriel dans $\R^3$]
	Soient $\vec{x} = \begin{pmatrix} x_1 \\ x_2 \\ x_3 \end{pmatrix}$ et $\vec{y} = \begin{pmatrix} y_1 \\ y_2 \\ y_3 \end{pmatrix}$ deux vecteurs de $\R^3$.
	\\
	On définit le produit vectoriel ainsi :
	\begin{center}
		$
		\appli{\wedge}{\R^3 \times \R^3}{\R^3}{(\vec{x}, \vec{y})}{
		\begin{pmatrix}
			x_2 y_3 - x_3 y_2 \\
			x_1 y_3 - x_3 y_1 \\
			x_1 y_2 - x_2 y_1
		\end{pmatrix}
		}
		$
	\end{center}
\end{definition}

Pour passer d'une forme cartésienne à une forme paramétrique, on applique le pivot de Gauss sur les lignes.
Pour passer d'une forme paramétrique à une forme cartésienne, on utilise le déterminant ou le produit vectoriel.



\chapter{Espaces vectoriels}\label{chap:espaces_vectoriels}
\def\arraystretch{1}

\section{Définitions}
\begin{definition}[Loi de composition]
	Soient $E$, $F$ deux ensembles et $f$ une application.
	\begin{enumerate}
		\item On dit que $f$ est une loi de composition \textbf{interne} si et seulement si : $f : E \times E \to E$.
		\item On dit que $f$ est une loi de composition \textbf{externe} si et seulement si : $f : E \times F \to E$.
	\end{enumerate}
\end{definition}

\begin{definition}[Magma]
	On appelle \textbf{magma} un ensemble $M$ muni d'une loi de composition \textbf{interne} \og $*$ \fg. 
	\\
	On le note $(M, *)$.
\end{definition}

\begin{definition}[Groupe]
	Soit $(G, *)$ un magma.
	\\ 
	On dit que $(G, *)$ est un groupe si et seulement si pour $g_1, g_2, g_3 \in G$ : 
    \begin{enumerate}
    		\item $(g_1 * g_2) * g_3 = g_1 * (g_2 * g_3)$
    		\item $\exists 0_G \in G,\ g_1 * 0_G = 0_G * g_1 = g_1$
    		\item $\exists g^{-1} \in G,\ g_1 * g^{-1} = g^{-1} * g_1 = 0_G$
    	\end{enumerate}
    \noindent De plus, si et seulement :
    \[ g_1 * g_2 = g_2* g_1 \]
    on dit que $(G, *)$ est un groupe \textbf{commutatif} ou \textbf{abélien}.
\end{definition}

\begin{definition}[Anneau \cite{bibmath_resume_groupes}]
	Soit $A$ un ensemble muni de deux lois de compositions \textbf{internes} \og $+$ \fg et \og $\cdot$ \fg sur $A$ telles que pour $a, b, c \in A$ : 
	\begin{enumerate}
			\item $(A, +)$ est un groupe \textbf{commutatif}.
			\item $a \cdot (b \cdot c) = (a \cdot b) \cdot c$.
			\item $a \cdot (b + c) = a \cdot b + a \cdot c$.
			\item $(b + c) \cdot a = b \cdot a + c \cdot a$.
			\item \og $\cdot$ \fg possède un \textbf{élément neutre}.
		\end{enumerate}
	\noindent On dit que $A$ est \textbf{intègre} si :
	\begin{enumerate}
			\item $A$ est commutatif : $a \cdot b = b \cdot a$.
			\item $a \cdot b = 0 \implies a = 0 \lor b = 0$.
		\end{enumerate}
\end{definition}

\begin{definition}[Corps \cite{bibmath_resume_groupes}]
    Un corps est un anneau \textbf{commutatif} dans lequel tout élément non nul est inversible.
\end{definition}

\begin{definition}[$\K$-espace vectoriel]
	Soit $\K$ un corps. 
	\\
	Un $\K$-espace vectoriel est un ensemble $E$ composé d'une loi de composition \textbf{interne} \og $+$ \fg et d'une loi de composition \textbf{externe} \og $\cdot$ \fg telles que :
	\begin{center}
		$
		\appli{+}{E \times E}{E}{(x,y)}{x+y}
		\qquad 
		\appli{\cdot}{\K \times E}{E}{(\lambda, u)}{\lambda \cdot u}
		$	
	\end{center}		
	
	\begin{enumerate}
		\item $(E, +)$ est un groupe commutatif.
		\item $\forall \lambda_1, \lambda_2 \in \K,\ u, v \in E :$
		\begin{enumerate}
			\item $\lambda_1 \cdot (u + v) = \lambda_1 \cdot u + \lambda_1 \cdot v$
			\item $(\lambda_1 + \lambda_2) \cdot u = \lambda_1 \cdot u + \lambda_2 \cdot u$
			\item $(\lambda_1 \cdot \lambda_2) \cdot u = \lambda_1 \cdot (\lambda_2 \cdot u)$
			\item $1 \cdot u = u$
		\end{enumerate}
	\end{enumerate}
\end{definition}

\begin{definition}[Sous-espace vectoriel]
	Soit $E$ un espace-vectoriel, $F$ est un sous-espace vectoriel de $E$ si :
	\begin{enumerate}
    		\item $F \subset E$
    		\item $F \neq \varnothing$
    		\item $\forall u, v \in F,\ \lambda \in \K,\ u + \lambda v \in F$
    	\end{enumerate}
\end{definition}

\begin{definition}[Somme directe]
	Soient $F_1, F_2 \subseteq E$. On dit que $F_1$ et $F_2$ sont en \textbf{somme directe} ou qu'ils sont \textbf{supplémentaires} dans $E$ si et seulement si :
	\begin{multicols}{2}
	    \begin{enumerate}
		\item $F_1 \cap F_2 = \{ 0_E \}$
		\item $F_1 + F_2 = E$
	\end{enumerate}
	\end{multicols}
	On note alors :
	\[ F_1 \oplus F_2 = E \]
\end{definition}

\begin{definition}
	Soient $n \in \N^*$, $E$ un espace vectoriel et $\mathcal{F} = (u_1, \ldots, u_n)$ une famille de vecteurs de $E$.
	\begin{enumerate}
		\item On dit que $\mathcal{F}$ est \textbf{libre} si et seulement si :
		\begin{align*}
			\forall (\lambda_1, \ldots, \lambda_n) \in \K^n,\ \sum_{i = 1}^n \lambda_i u_i = 0_E \implies \lambda_1 = \cdots = \lambda_n = 0
		\end{align*}
		\item On dit que $\mathcal{F}$ est \textbf{génératrice} si et seulement si :
		\begin{align*}
			\forall x \in E,\ \exists (\lambda_1, \ldots, \lambda_n) \in \K^n,\ x = \sum_{i = 1}^n \lambda_i u_i 
		\end{align*}
	\end{enumerate}
\end{definition}

\begin{remark}
	Si une famille n'est pas libre, on dit qu'elle est liée.
\end{remark}

\section{Base et dimension}
\begin{definition}[Base]
	Une famille de vecteurs est une base si elle est \textbf{libre} et \textbf{génératrice}.
\end{definition}

\begin{proposition}
	Soient un espace vectoriel $E$ et $\mathcal{F} = (u_1, \ldots, u_n)$ une famille de vecteurs de $E$.
	\[ 
	\mathcal{F} \text{ est une base} \iff 
	\forall x \in E,\ \exists ! (\lambda_1, \ldots, \lambda_n) \in \K^n,\ x = \sum_{i = 1}^n \lambda_i u_i
	\]
\end{proposition}

\begin{proof}
	L'existence est évidente car $\mathcal{F}$ est génératrice. 
	\\
	Montrons l'unicité : Soient $\lambda_1, \ldots, \lambda_n, \mu_1, \ldots, \mu_n \in \K$.
	On a d'une part :
	\[ \forall x \in E,\ x = \sum_{i=1}^{n} \lambda_i u_i \]
	puis d'autre part :
	\[ \forall x \in E,\ x = \sum_{i=1}^{n} \mu_i u_i \]
	Donc on a :
	\begin{align*}
		&\sum_{i=1}^{n} \lambda_i u_i = \sum_{i=1}^{n} \mu_i u_i \\
		\iff &\sum_{i=1}^{n} (\lambda_i - \mu_i) u_i = 0 
	\end{align*}
	Or $\mathcal{F}$ est libre, donc pour tout $i \in \llbracket 1, n \rrbracket,\ \lambda_i - \mu_i = 0 \implies \lambda_i = \mu_i$.
\end{proof}

\begin{proposition}
	Soit $\mathcal{F} = (u_1, \ldots, u_n)$ une famille de vecteurs de $\R^n$. $\mathcal{F}$ est une base de $\R^n$ si et seulement si :
	\[ \det(\mathcal{F}) \neq 0 \]
\end{proposition}

\begin{definition}[Dimension d'un espace vectoriel]
	Soit $E$ un espace vectoriel, on appelle dimension de $E$, notée $\dim(E)$, le nombre d'éléments d'une base de $E$. 
\end{definition}


\begin{proposition}
	Soient $E$ un espace vectoriel de dimension $n$ et $\mathcal{F} = (u_1, \ldots, u_n)$ une famille de vecteurs de $E$. Alors on a :
	\begin{enumerate}
		\item $\mathcal{F}$ est une base.
		\item $\mathcal{F}$ est libre.
		\item $\mathcal{F}$ est génératrice.
	\end{enumerate}
	\noindent Ainsi il suffit de montrer que $\mathcal{F}$ est libre pour montrer les deux autres propriétés.
\end{proposition}

\begin{theorem}[Théorème de la base incomplète]
	Toute famille libre peut être complétée en une base.
\end{theorem}

\begin{theorem}[Théorème de la base extraite]
    De toute famille génératrice, on peut extraire une base.
\end{theorem}

\begin{theorem}
    Chaque espace vectoriel admet une base.
\end{theorem}

\begin{proposition}
	Soient $E$ un espace vectoriel et $F, G\subset E$.
	\begin{multicols}{2}
	    \begin{enumerate}
    		\item $\dim(F \oplus G) = \dim(F) + \dim(G)$ 
    		\item $\dim(F + G) \leq \dim(E)$
    	\end{enumerate}
	\end{multicols}
\end{proposition}

\begin{theorem}[Théorème de Grassman]
	Soient $F$ et $G$ deux espaces vectoriels.
	\[ \dim(F+G) = \dim(F) + \dim(G) - \dim(F\cap G) \]
\end{theorem}


\chapter{Applications linéaires}
\def\arraystretch{1}

\section{Définitions}
\begin{definition}[Application linéaire]
    Soient $E$ et $F$ deux $\K$-espaces vectoriels et $f : E \to F$.
    On dit que $f$ est une \textbf{application linéaire} si et seulement si $\forall x_1, x_2 \in E,\ \lambda \in \K$ :
    \[ f(x_1 + \lambda x_2) = f(x_1) + \lambda f(x_2) \]
\end{definition}

\begin{definition}
    Soient $E$ et $F$ deux $\K$-espaces vectoriels et $f : E \to F$ une application linéaire.
    \\
    On définit le \textbf{noyau} de $f$, noté $\ker(f)$ tel que :
    \[ \ker(f) = \{ x \in E : f(x) = 0_E \} \] 
    On définit l'\textbf{image} de $f$, notée $\Im(f)$ telle que :
    \[ \Im(f) = \{ y \in F : \exists x \in E : y = f(x) \} \]
\end{definition}

\begin{theorem}
	Soient $E, F$ des espaces vectoriels et $f : E \to F$ une application linéaire.
	\begin{enumerate}
		\item $\ker(f) \subseteq E$.
		\item $\Im(f) \subseteq F$.
	\end{enumerate}
\end{theorem}

\begin{proof}
	\leavevmode 
	\begin{enumerate}
		\item $\forall x_1, x_2 \in \ker(f),\ \lambda \in \K$.
		\begin{align*}
			f(x_1 + \lambda x_2) &= f(x_1) + \lambda f(x_2) \\ 
								 &= 0 + \lambda 0 \\
								 &= 0
		\end{align*}
		\item $\forall y_1, y_2 \in \Im(f),\ \lambda \in \K$.
		\[ y_1 \in \Im(f) \iff \exists x_1 : f(x_1) = y_1 \]
		\[ y_2 \in \Im(f) \iff \exists x_2 : f(x_2) = y_2 \]
		\begin{align*}
			y_1 + \lambda y_2 &= f(x_1) + \lambda f(x_2) \\
			                  &= f(x_1 + \lambda x_2)
		\end{align*}
		Ainsi $y_1 + \lambda y_2 \in \Im(f)$.
	\end{enumerate}
\end{proof}

\begin{definition}
    \par \noindent Soient $E$ et $F$ deux $\K$-espaces vectoriels et $f$ une application linéaire de $E$ vers $F$.
    \begin{enumerate}
        \item On dit que $f$ est un \textbf{morphisme} de $E$ vers $F$, on note $f \in \mathcal{L}(E, F)$.
        \item Si $E = F$, on dit que $f$ est un \textbf{endomorphisme} de $E$, on note $f \in \mathcal{L}(E)$ ou $\End(E)$.
        \item Si $f \in \mathcal{L}(E, F)$ est une \textbf{bijection}, alors $f$ est un \textbf{isomorphisme}. On le note $f : E \overset{\sim}{\to} F$.
        \item Si $f \in \mathcal{L}(E)$ est un \textbf{isomorphisme}, on dit que $f$ est un \textbf{automorphisme} de $E$ et on le note $\Aut(E)$.
        \\
        $E$ et $F$ sont appelés \textbf{isomorphes} s'il existe un \textbf{isomorphisme} de l'un vers l'autre, on écrit parfois $E \cong F$.
    \end{enumerate}
\end{definition}

Soit $\varepsilon = \{ f \in \mathcal{C}^N(I) : f^{(n)} + a_{n - 1} f^{(n - 1)} + \cdots + a_1 f' + a_0 f = 0 \},\ n, N \in \N^*,\ a_i \in \mathcal{C}^0(I),\ I$ un intervalle ouvert.

\begin{proposition}
    Soit $n \in \N$.
    \begin{multicols}{2}
        \begin{enumerate}
            \item $\varepsilon \subset \mathcal{C}^n(I)$
            \item $\dim(E) = n$
        \end{enumerate}
    \end{multicols}
\end{proposition}

\begin{definition}[Rang d'un morphisme]
    Soit $f \in \mathcal{L}(E, F)$. On appelle $\rg(f) = \dim(\Im(f))$ le rang de $f$.
\end{definition}

\begin{theorem}[Théorème du rang]
    Soient $E, F$ des espaces vectoriels de dimensions finies ou infinies et $f \in \mathcal{L}(E, F)$, alors :
    \[ \dim(E) = \dim(\ker(f)) + \rg(f) \]
\end{theorem}

\begin{corollary}
    Soit $f \in \mathcal{L}(E)$.
    \begin{align*}
        \begin{cases}
            \ker(f) + \Im(f) = \ker(f) \oplus \Im(f) \\
            \ker(f) \oplus \Im(f) = E
        \end{cases}
        \iff 
        \ker(f) \cap \Im(f) = \{0_E\}
    \end{align*}
\end{corollary}

\begin{proposition}
    Soit $f \in \mathcal{L}(E, F)$ et $(b_1, \ldots, b_n)$ une base de $E$. Si l'on connait $f(b_i) \in F$, $\forall i \in \llbracket 1, n \rrbracket$, on connait toute l'application $f$.
\end{proposition}

\begin{corollary}
    Soit $(b_1, \ldots, b_n)$ une base de $E$. Alors $\varphi \in \mathcal{L}(E, \R^n)$ est un isomorphisme donné par :
    \[ \varphi(b_i) = e_i,\ \forall i \in \llbracket 1, n \rrbracket \]
    \[ \varphi \left( \sum_{i = 1}^{n} x_i b_i \right) = \sum_{i = 1}^{n} x_i e_i = 
    \begin{pmatrix}
        x_1 \\
        \vdots \\
        x_n
    \end{pmatrix}
    \in \R^n
     \]
\end{corollary}

\begin{definition}
    Soient $f \in \mathcal{L}(E, F)$ et $\mathcal{B} = (b_1, \ldots, b_n)$ une base de $E$ et $\mathcal{C} = (c_1, \ldots, c_n)$ une base de $F$.
    Alors on peut définir une matrice $A \in \mathcal{M}_{m,n}(\mathbb{K}),\ \forall i \in \llbracket 1, n \rrbracket$. 
    \[ f(b_i) = \sum_{j = 1}^{n} a_{j,i} c_j \]
\end{definition}

\begin{lemma} Soient $\K$ un corps et $E$ un espace-vectoriel de dimension $n \in \N$.
    \[ \dim(E) = n \implies E \cong \K^n \]
\end{lemma}

\begin{lemma}
    Soient $\mathcal{B} = (b_1, \ldots, b_n)$ une base de $E$, $\lambda_1, \ldots, \lambda_n \in \K$ et $\varphi \in \mathcal{L}(E, \K^n)$  alors :
    \begin{center}
    	$
    	\appli{\varphi}{E}{\K^n}{\sum_{i = 1}^{n} \lambda_i b_i}{
    	\begin{pmatrix}
    		\lambda_1 \\
    		\vdots \\
    		\lambda_n
    	\end{pmatrix}
    	}
    	$
    \end{center}
    est une bijection, alors c'est un isomorphisme.
\end{lemma}

\begin{proposition}
    Soient $E, F$ des espaces vectoriels $f \in \mathcal{L}(E, F)$ un morphisme et $A$ sa matrice associée. 
    \[ y = f(x) \iff y = A \cdot x \]
\end{proposition}

\section{Projecteurs et symétries}
\begin{definition}
    Soit $f \in \mathcal{L}(E)$. Notons $f^2 = f \circ f$.
    \begin{enumerate}
        \item On dit que $f$ \textbf{est idempotente/une projection} si et seulement si : $f^2 = f$.
        \item On dit que $f$ \textbf{est involutive/une symétrie linéaire} si et seulement si : $f^2 = id_E$.
    \end{enumerate}
\end{definition}

\begin{proposition} 
    Soit $p \in \mathcal{L}(E)$.
    \begin{enumerate}
        \item $id_E - p \text{ est une projection} \iff p \text{ est une projection}$
        \item $2p - id_E \text{ est une symétrie} \iff p \text{ est une projection}$
    \end{enumerate}
\end{proposition}

\begin{proof}
	\leavevmode
    \begin{enumerate}
        \item Posons $f(x) = id_E(x) - p(x)$. Montrons que $f(f(x)) = id_E(x) - p(x)$ si et seulement si $p(p(x))$.
        \begin{align*}
            f(f(x)) &= id_E(id_E(x) - p(x)) - p(id_E(x) - p(x)) \\ 
            &= id_E(id_E(x)) - id_E(p(x)) - p(id_E(x)) + p(p(x)) \\ 
            &= id_E(x) - p(x) - p(x) + p(p(x)) \\ 
            &= id_E(x) - 2p(x) + p(p(x)) \\ 
            &= id_E(x) - [2p(x) - p(p(x))]
        \end{align*}
        \begin{align*}
            id_E(x) - [2p(x) - p(p(x))] = id_E(x) - p(x) &\iff p(p(x)) = p(x) \\ 
            &\iff 2p(x) - p(p(x)) = p(x) \\ 
            &\iff 2p(x) = p(x) + p(p(x)) \\
            &\iff p(x) = p(p(x))
        \end{align*}
        \item Posons $s(x) = 2p(x) - id_E(x)$. 
        \\
        Montrons que $s(s(x)) = id_E(x)$ si et seulement si $p(p(x)) = p(x)$.
        \begin{align*}
            s(s(x)) &= s(2p(x) - id_E(x)) \\
            &= 2p(2p(x) - id_E(x)) - id_E(2p(x) - id_E(x)) \\
            &= 4p(p(x)) - 2p(id_E(x)) - 2id_E(p(x)) + id_E(id_E(x))
        \end{align*}
        \begin{align*}
            4p(p(x)) - 4p(x) + id_E(x) = id_E(x) &\iff 4p(p(x)) - 4p(x) = 0 \\
            &\iff 4p(p(x)) = 4p(x) \\
            &\iff p(p(x)) = p(x)
        \end{align*}
    \end{enumerate}
\end{proof}

\begin{definition}
    Soient $E = F \oplus G$, $u \in F$, $v \in G$, alors l'application 
    \begin{center}
    	$
    	\appli{p_F}{E}{E}{u+v}{u}
    	$
    \end{center}
    est appelée un \textbf{projecteur} sur $F$ \textbf{parallélement} à $G$.
\end{definition}

\begin{proposition}
    Soit $p_F$ définie comme dans la définition précédente.
    \begin{enumerate}
        \item $p_F$ est une projection.
        \item Soit $p$ une projection, $p$ est un projecteur sur $\Im(p)$ parallélement à son noyau $\ker(p)$.
    \end{enumerate}
\end{proposition}

\section{Rotations.}
\begin{definition}
	Soit $n \in \N^*$.
    \[ \operatorname{GL}(\R, n) = \{ A \in \mathcal{M}_n(\R) : \det(A) \neq 0 \} \]
\end{definition}

\begin{definition}
	Soit $n \in \N^*$.
    \[ \operatorname{SL}(\R, n) = \{ A \in \mathcal{M}_n(\R) : \det(A) = 1 \} \]
\end{definition}

\begin{definition}
	Soit $n \in \N^*$.
    \[ \operatorname{O}(n) = \{ R \in \mathcal{M}_n(\R) : \forall (x, y) \in (\R^n)^2,\ \langle Rx|Ry \rangle = \langle x|y \rangle \} \]
\end{definition}

\begin{definition}
    Soit $n \in \N^*$.
    \[ \operatorname{SO}(n) = \{ R \in \operatorname{O}(n) : \det(R) = 1 \} \]
\end{definition}

\begin{proposition}
	Soit $R \in \mathcal{M}_n(\K)$.
    \begin{enumerate}
        \item $R \in \operatorname{O}(n) \iff R^T \cdot R = I_n$.
        \item $R \in \operatorname{O}(n) \implies \det(R) \in \{ \pm 1 \}$.
    \end{enumerate}
\end{proposition}

\begin{corollary}
	Soit $n \in \N^*$.
    \[ \operatorname{O}(n) = \{ R \in \mathcal{M}_n(\K) : R^T \cdot R = I_n \} \]
\end{corollary}

\section{Changements de bases et matrices associées aux applications linéaires.}

\begin{definition}
    Soient $n, p \in \N^*$, $E$ et $F$ deux $\K$-espaces vectoriels, $\mathcal{B} = (b_1, \ldots, b_p)$ une base de $E$, $\mathcal{B}' = (b_1', \ldots, b_n')$ une base de $F$ et $f \in \mathcal{L}(E, F)$.
    \\
    On peut définir une matrice $\operatorname{Mat}_{\mathcal{B},\mathcal{B}'} (f) \in \mathcal{M}_{n,p}(\K)$ telle que :
    \begin{align*}
        \hspace{0.4cm}
        \begin{matrix}
            f(b_1) & \cdots & f(b_p) 
        \end{matrix}
    \end{align*}
    \vspace{-0.8cm}
    \begin{align*}
        \begin{matrix}
            b_1' \\
            \vdots \\ 
            b_n'
        \end{matrix}
        \begin{pmatrix}
            a_{1,1} & \cdots & a_{1,p} \\
            \vdots & \ddots & \vdots \\ 
            a_{n,1} & \cdots & a_{n,p}
        \end{pmatrix}
    \end{align*}
    avec les coefficients $a_{1,1}, \ldots, a_{n,p} \in \K$ tels que : 
    \[ f(b_j) = \sum_{i = 0}^n a_{i,j} b_i' \]
\end{definition}

\begin{definition}[Matrice de passage]
    Soient $n \in \N^*$, $E$ et $F$ deux $\K$-espaces vectoriels de dimension $n$, $\mathcal{B} = (b_1, \ldots, b_n)$ et $\mathcal{B}' = (b_1', \ldots, b_n')$ deux bases de $E$. On appelle \textbf{matrice de passage} de $\mathcal{B}$ à $\mathcal{B}'$ la matrice carrée de taille $n$ dont la $j-$ième colonne est formée des coordonnées de $b_j'$ dans la base $\mathcal{B}$.
\end{definition}

\par \noindent Concrètement, écrivons : 
\begin{align*}
    \hspace{0.45cm}
    \begin{matrix}
        b_1' & \cdots & b_n' 
    \end{matrix}
\end{align*}
\vspace{-1.1cm}
\begin{align*}
    \begin{matrix}
        b_1 \\
        \vdots \\ 
        b_n
    \end{matrix}
    \begin{pmatrix}
        a_{1,1} & \cdots & a_{1,n} \\
        \vdots & \ddots & \vdots \\ 
        a_{n,1} & \cdots & a_{n,n}
    \end{pmatrix}
\end{align*}
Avec les coefficients $a_{1,1}, \ldots, a_{n,n} \in \K$ tels que :
\[ b_j' = \sum_{i=1}^{n} a_{i,j} b_i  \]

\begin{definition}
    Soient $A, B \in \mathcal{M}_{m,n}(\K)$.
    \begin{enumerate}
        \item On dit que $A$ et $B$ sont \textbf{équivalentes} si et seulement si : 
        \[ \exists P \in \operatorname{GL}(n),\ Q \in \operatorname{GL}(m),\ B = Q^{-1} A P \]
        \item On dit que $A$ et $B$ sont \textbf{similaires} si et seulement si :
        \[ \exists P \in \operatorname{GL}(n),\ B = P^{-1} A P \]
    \end{enumerate}
\end{definition}



\part{Annexes}
\def\arraystretch{1.5}
\begin{table}[!h]
    \centering
    \begin{tabular}{cc}
        \toprule
        Soient $f, g$ des fonctions \\
        \midrule
        $C \in \R,\ C f(x)$ & $C f'(x)$  \\
        $(f + g)'(x)$ & $f'(x) + g'(x)$ \\
        $(f \cdot g)'(x)$ & $f'(x) \cdot g(x) + f(x) \cdot g'(x)$ \\
        $\left( \frac{f}{g} \right)'(x)$ & $\frac{f'(x) \cdot g(x) - f(x) \cdot g'(x)}{(g(x))^2}$ \\
        $(f \circ g)'(x)$ & $f'(g(x)) \cdot g'(x) $ \\
        \bottomrule
    \end{tabular}
    \caption{Formules de dérivation.}
    \label{tab:formules_derivation}
\end{table}

\begin{table}[!h]
    \centering
    \begin{tabular}{ccc}
        \toprule
        $f(x)$ & $f'(x)$ & $\mathcal{D}_f$ \\ 
        \midrule
        $C \in \R$ & $0$ & $\R$ \\
        $x^a,\ a \in \R$ & $a x^{a - 1}$ & $\R^* \text{ si } a \in \Z^-, \R_+^* \text{ sinon }$ \\
        $\sqrt{x}$ & $\frac{1}{2\sqrt{x}}$ & $\R_+^*$ \\
        $e^x$ & $e^x$ & $\R$ \\
        $\ln{\abs{x}}$ & $\frac{1}{x}$ & $\R_+^*$ \\
        $\cos{x}$ & $-\sin{x}$ & $\R$ \\
        $\sin{x}$ & $\cos{x}$ & $\R$ \\
        $\tan{x}$ & $\frac{1}{\cos^2{x}} = 1 + \tan^2{x}$ & $\forall k \in \Z, \R \backslash \left\{ \frac{\pi}{2} + k \pi \right\}$ \\
        $\cosh{x}$ & $\sinh{x}$ & $\R$ \\
        $\sinh{x}$ & $\cosh{x}$ & $\R$ \\
        $\tanh{x}$ & $1 - \tanh^2{x} = \frac{1}{\cosh^2{x}}$ & $\R$ \\
        $\arctan{x}$ & $\frac{1}{1 + x^2}$ & $\R$ \\
        $\arcsin{x}$ & $\frac{1}{\sqrt{1 - x^2}}$ & $]-1, 1[$ \\
        $\arccos{x}$ & $-\frac{1}{\sqrt{1 - x^2}}$ & $]-1, 1[$ \\
        \bottomrule
    \end{tabular}
    \caption{Dérivées usuelles.}
    \label{tab:derivees_usuelles}
\end{table}

\begin{table}[!h]
    \centering
    \begin{tabular}{ccc}
         \toprule
         $f(x)$ & $F(x)$ & $I$ \\
         \midrule
         $x^a,\ a \in \R \backslash \{1\}$ & $\frac{1}{a + 1} x^{a+1} + C$ & $\R_+^* \text{ ou } \R_-^* \text{ si } a \in \Z^-,\ \R_+^* \text{ sinon}$ \\
         $\frac{1}{x}$ & $\ln(\abs{x}) + C$ & $\R_+^* \text{ ou } \R_-^*$ \\
         $e^{ax},\ a \in \R$ & $\frac{1}{a} e^{ax} + C$ & $\R$ \\
         $\cos{x}$ & $\sin{x} + C$ & $\R$ \\
         $\sin{x}$ & $-\cos{x} + C$ & $\R$ \\
         $1 + \cos^2{x} = 1 + \tan^2{x}$ & $\tan{x} + C$ & $\forall k \in \Z,\ \left] -\frac{\pi}{2} + k \pi, \frac{\pi}{2} + k \pi \right[$ \\
         $\cosh{x}$ & $\sinh{x} + C$ & $\R$ \\
         $\sinh{x}$ & $\cosh{x} + C$ & $\R$ \\
         $\frac{1}{\cosh^2{x}} \text{ ou } 1 + \tanh^2{x}$ & $\tanh{x} + C$ & $\R$ \\
         $\frac{1}{1 + x^2}$ & $\arctan{x} + C$ & $\R$ \\
         $\frac{1}{\sqrt{1 - x^2}}$ & $-\arccos{x} + C \text{ ou } \arcsin{x} + C$ & $]-1, 1[$ \\
         \bottomrule
    \end{tabular}
    \caption{Primitives usuelles, $C \in \R$.}
    \label{tab:primitives_usuelles}
\end{table}

\begin{table}[!h]
    \centering
    \begin{tabular}{cc}
         \toprule
         $f(x)$ & DL \\ 
         \midrule
         $e^x$ & $\sum_{k=0}^n \frac{x^k}{k!} + o(x^n)$ \\
         $\cosh(x)$ & $ \sum_{k=0}^n \frac{x^{2k}}{(2k)!} + o(x^n)$ \\
         $\sinh(x)$ & $ \sum_{k=0}^n \frac{x^{2k + 1}}{(2k + 1)!} + o(x^{2n + 1})$ \\
         $\cos(x)$ & $ \sum_{k=0}^n (-1)^k \frac{x^{2k}}{(2k)!} + o(x^{2n})$ \\
         $\sin(x)$ & $ \sum_{k=0}^n (-1)^k \frac{x^{2k+1}}{(2k+1)!} + o(x^{2n + 1})$ \\
         $(1 + x)^{a}$ & $ \sum_{k=0}^n \binom{a}{k} x^k + o(x^n)$ \\ 
         $\frac{1}{1 - x}$ & $ \sum_{k=0}^n x^k + o(x^n)$ \\ 
         $\frac{1}{1 + x}$ & $ \sum_{k=0}^n (-1)^k x^k + o(x^n)$ \\ 
         $\ln(1 + x)$ & $ \sum_{k=1}^n (-1)^{k - 1} \frac{x^k}{k} + o(x^n)$ \\ 
         $\ln(1 - x)$ & $ \sum_{k=1}^n \frac{x^k}{k} + o(x^n)$ \\ 
         $\arctan(x)$ & $ \sum_{k=0}^n (-1)^k \frac{x^{2k+1}}{2k+1} + o(x^{2n+1})$ \\
         \bottomrule
    \end{tabular}
    \caption{Développements limités usuels en 0.}
    \label{tab:dl_usuels}
\end{table}

\begin{table}[!h]
    \centering
    \begin{tabular}{cc}
         \toprule 
         $e^x - 1 \underset{x \to 0}{\sim} x$ \\
         $\ln(1+x) \underset{x \to 0}{\sim} x$ \\
         $\ln(x) \underset{x \to 1}{\sim} x - 1$ \\ 
         $\sin(x) \underset{x \to 0}{\sim} \sinh(x) \underset{x \to 0}{\sim} \tan(x) \underset{x \to 0}{\sim} \tanh(x) \underset{x \to 0}{\sim} \arcsin(x) \underset{x \to 0}{\sim} \arctan(x) \underset{x \to 0}{\sim} x$ \\ 
         $1 - \cos(x) \underset{x \to 0}{\sim} \frac{x^2}{2}$ \\ 
         $\cosh(x) - 1 \underset{x \to 0}{\sim} \frac{x^2}{2}$ \\ 
         $\arccos(x) \underset{\substack{x \to 1 \\ x < 1}}{\sim} \sqrt{2}\sqrt{1-x}$ \\ 
         $(1 + x)^a - 1 \underset{x \to 0}{\sim} ax,\ a \in \R^*$ \\ 
         $\frac{1}{1-x} - 1 \underset{x \to 0}{\sim} x$ \\ 
         $\cosh(x) \underset{x \to +\infty}{\sim} \sinh(x) \underset{x \to +\infty}{\sim} \frac{e^x}{2}$ \\ 
         \bottomrule
    \end{tabular}
    \caption{\'Equivalents usuels}
    \label{tab:equiv_usuels}
\end{table}

\printbibliography

\end{document}
